% 电力系统分析课程笔记 —— 主文件（正文排版在 0.样式.tex，章节文件只放内容）
% 编译：xelatex 1.main.tex 两遍；成品见 build.ps1
\documentclass[11pt,a4paper,oneside,openany]{book}
\usepackage[UTF8,fontset=windows,scheme=plain]{ctex}
\usepackage{fontspec}
\setmainfont{Times New Roman}
\setsansfont{Microsoft YaHei}
\newcommand{\YHBbooktrue}{}
\newcommand{\YHBcourse}{电力系统分析}
\newcommand{\YHBversion}{v0.1}
\newcommand{\YHBquote}{电压、功率与稳定，三者不可兼得。}
\newcommand{\YHBcoverSize}{40pt}
\newcommand{\YHBcoverLead}{48pt}
\usepackage{xcolor}
\definecolor{yhbAccent}{HTML}{1E4E3A}
\definecolor{yhbAccent2}{HTML}{3F7A5C}
% ============================================================
% 0.样式.tex —— 课程笔记统一正文排版
% 用法：主文件先定义课程信息与主色，再 % ============================================================
% 0.样式.tex —— 课程笔记统一正文排版
% 用法：主文件先定义课程信息与主色，再 % ============================================================
% 0.样式.tex —— 课程笔记统一正文排版
% 用法：主文件先定义课程信息与主色，再 \input{0.样式.tex}
%   \newcommand{\YHBcourse}{微积分A}      % 课程名（页眉与书签用）
%   \newcommand{\YHBversion}{v1.2}        % 版本号（封面用）
%   \newcommand{\YHBquote}{……}            % 封面题句
%   \definecolor{yhbAccent}{HTML}{24465C} % 主色（深）
%   \definecolor{yhbAccent2}{HTML}{347A80}% 辅色（浅）
%   \YHBbooktrue / \YHBbookfalse          % 是否 book 文档类（有 \chapter）
% 章节文件只写内容，不要重复 \usepackage、不要改字号。
% ============================================================

% ---------- 页面与字体 ----------
\usepackage{geometry}
\geometry{a4paper,left=24mm,right=24mm,top=23mm,bottom=24mm,headheight=15pt,headsep=7mm,footskip=13mm}
\usepackage{microtype}
\usepackage{amsmath,amssymb,amsfonts,mathtools}
\usepackage{esint}
\usepackage{graphicx}
\usepackage{adjustbox}
\usepackage{array,booktabs,multirow,makecell,tabularx,longtable}
\usepackage{xcolor,colortbl}
\usepackage{enumitem}
\usepackage{needspace,etoolbox}
\usepackage{titlesec,titletoc}
\usepackage{fancyhdr}
\usepackage{caption}
\usepackage[most]{tcolorbox}

% ---------- 颜色别名（章节里用的旧名继续有效） ----------
\colorlet{mainblue}{yhbAccent}
\colorlet{secondcyan}{yhbAccent2}
\colorlet{thirdteal}{yhbAccent2}
\colorlet{highlight}{yhbAccent!6!white}
\colorlet{tablehead}{yhbAccent!12!white}
\definecolor{textgray}{HTML}{4A5057}

% ---------- 课程标签 ----------
\providecommand{\YHBdoclabel}{课程笔记}
% ---------- 正文节奏 ----------
\linespread{1.3}
\setlength{\parindent}{0pt}
\setlength{\parskip}{6pt plus 1.5pt minus 1.5pt}
\setlength{\emergencystretch}{2.5em}
\raggedbottom
\widowpenalty=10000
\clubpenalty=10000
\setlength{\tabcolsep}{7pt}
\renewcommand{\arraystretch}{1.3}
\allowdisplaybreaks[2]
\AtBeginDocument{%
  \setlength{\abovedisplayskip}{10pt plus 2pt minus 2pt}%
  \setlength{\belowdisplayskip}{10pt plus 2pt minus 2pt}%
  \setlength{\abovedisplayshortskip}{6pt}%
  \setlength{\belowdisplayshortskip}{6pt}%
}
\captionsetup{font=small,labelfont=bf,skip=6pt}

% ---------- 列表 ----------
\setlist{leftmargin=2em,labelsep=0.6em,itemsep=4pt,parsep=2pt,topsep=6pt,partopsep=0pt}
\setlist[itemize,2]{label=\textbullet,leftmargin=1.6em}

% ---------- 标题层级 ----------
\ifdefined\YHBbooktrue
  \titleformat{\chapter}[display]
    {\sffamily\bfseries\color{yhbAccent}}
    {\normalsize\sffamily 第\,\thechapter\,章}
    {6pt}
    {\LARGE}
    [{\vspace{5pt}{\color{yhbAccent!55}\titlerule[1pt]}}]
  \titlespacing*{\chapter}{0pt}{14pt}{24pt}
  \titleformat{\section}{\Large\sffamily\bfseries\color{yhbAccent}}{\thesection}{0.7em}{}
  \titlespacing*{\section}{0pt}{18pt plus 3pt minus 2pt}{8pt}
\else
  \titleformat{\section}{\Large\sffamily\bfseries\color{yhbAccent}}{\thesection}{0.7em}{}[{\vspace{3pt}{\color{yhbAccent!45}\titlerule[0.7pt]}}]
  \titlespacing*{\section}{0pt}{18pt plus 3pt minus 2pt}{9pt}
\fi
\titleformat{\subsection}{\large\sffamily\bfseries\color{yhbAccent!85!black}}{\thesubsection}{0.7em}{}
\titleformat{\subsubsection}{\normalsize\bfseries\color{black!85}}{\thesubsubsection}{0.7em}{}
\titlespacing*{\subsection}{0pt}{12pt plus 2pt minus 1pt}{5pt}
\titlespacing*{\subsubsection}{0pt}{9pt plus 2pt minus 1pt}{4pt}
\pretocmd{\section}{\Needspace{6\baselineskip}}{}{}
\pretocmd{\subsection}{\Needspace{5\baselineskip}}{}{}
\pretocmd{\subsubsection}{\Needspace{4\baselineskip}}{}{}

\ifdefined\YHBbooktrue
  \setcounter{secnumdepth}{2}
  \setcounter{tocdepth}{1}
\else
  \setcounter{secnumdepth}{3}
  \setcounter{tocdepth}{2}
\fi
\renewcommand{\contentsname}{目录}
% 目录编号颜色由 titletoc 的 titlecontents 控制，不再重定义 \numberline

% ---------- 目录样式 ----------
\ifdefined\YHBbooktrue
  \titlecontents{chapter}[0pt]
    {\addvspace{6pt}\sffamily\bfseries\color{yhbAccent}}
    {\contentslabel{2.2em}}{}{\titlerule*[0.6em]{.}\contentspage}
  \titlecontents{section}[2.2em]
    {\addvspace{2pt}}
    {\contentslabel{2.2em}}{}{\titlerule*[0.6em]{.}\contentspage}
\else
  \titlecontents{section}[0pt]
    {\addvspace{5pt}\sffamily\bfseries\color{yhbAccent}}
    {\contentslabel{2.2em}}{}{\titlerule*[0.6em]{.}\contentspage}
  \titlecontents{subsection}[2.2em]
    {\addvspace{2pt}}
    {\contentslabel{2.4em}}{}{\titlerule*[0.6em]{.}\contentspage}
\fi

% ---------- 页眉页脚 ----------
\pagestyle{fancy}
\fancyhf{}
\fancyhead[L]{\footnotesize\sffamily\color{textgray}\YHBcourse\ \YHBdoclabel}
\ifdefined\YHBbooktrue
  \fancyhead[R]{\footnotesize\sffamily\color{textgray}\nouppercase{\leftmark}}
  \renewcommand{\chaptermark}[1]{\markboth{#1}{}}
\else
  \fancyhead[R]{\footnotesize\sffamily\color{textgray}\nouppercase{\rightmark}}
  \renewcommand{\sectionmark}[1]{\markright{#1}}
\fi
\fancyfoot[C]{\footnotesize\color{textgray}\thepage}
\renewcommand{\headrulewidth}{0.35pt}
\renewcommand{\headrule}{{\color{yhbAccent!45}\hrule width\headwidth height\headrulewidth}}
\fancypagestyle{plain}{%
  \fancyhf{}%
  \fancyfoot[C]{\footnotesize\color{textgray}\thepage}%
  \renewcommand{\headrulewidth}{0pt}%
}

% ---------- 公式与考点框 ----------
\newcommand{\keyformula}[1]{%
  \begin{tcolorbox}[enhanced,breakable,colback=highlight,colframe=yhbAccent!30!white,boxrule=0.5pt,arc=2.5pt,left=11pt,right=11pt,top=9pt,bottom=9pt,boxsep=0pt,before skip=11pt,after skip=11pt,halign=center,width=\linewidth]
  \adjustbox{max width=\linewidth-2em}{$\displaystyle #1$}
  \end{tcolorbox}}
\newcommand{\corepoint}[1]{%
  \begin{tcolorbox}[enhanced,breakable,colback=highlight,colframe=yhbAccent2,boxrule=0pt,borderline west={2.4pt}{0pt}{yhbAccent2},arc=0pt,left=11pt,right=11pt,top=8pt,bottom=8pt,before skip=11pt,after skip=11pt]
  \textbf{\color{yhbAccent}核心考点}\quad #1
  \end{tcolorbox}}
\newcommand{\symb}[2]{\textbf{#1}：#2\par}

% ---------- 超链接（务必最后加载） ----------
\usepackage[unicode=true,colorlinks=true,linkcolor=yhbAccent,citecolor=yhbAccent,urlcolor=yhbAccent!85!black,pdfborder={0 0 0},bookmarksnumbered=true,pdfstartview=FitH]{hyperref}
\usepackage{hypcap}
\hypersetup{pdftitle={\YHBcourse 课程笔记},pdfauthor={浮萍},pdfsubject={课程学习与复习笔记},pdfcreator={XeLaTeX}}

%   \newcommand{\YHBcourse}{微积分A}      % 课程名（页眉与书签用）
%   \newcommand{\YHBversion}{v1.2}        % 版本号（封面用）
%   \newcommand{\YHBquote}{……}            % 封面题句
%   \definecolor{yhbAccent}{HTML}{24465C} % 主色（深）
%   \definecolor{yhbAccent2}{HTML}{347A80}% 辅色（浅）
%   \YHBbooktrue / \YHBbookfalse          % 是否 book 文档类（有 \chapter）
% 章节文件只写内容，不要重复 \usepackage、不要改字号。
% ============================================================

% ---------- 页面与字体 ----------
\usepackage{geometry}
\geometry{a4paper,left=24mm,right=24mm,top=23mm,bottom=24mm,headheight=15pt,headsep=7mm,footskip=13mm}
\usepackage{microtype}
\usepackage{amsmath,amssymb,amsfonts,mathtools}
\usepackage{esint}
\usepackage{graphicx}
\usepackage{adjustbox}
\usepackage{array,booktabs,multirow,makecell,tabularx,longtable}
\usepackage{xcolor,colortbl}
\usepackage{enumitem}
\usepackage{needspace,etoolbox}
\usepackage{titlesec,titletoc}
\usepackage{fancyhdr}
\usepackage{caption}
\usepackage[most]{tcolorbox}

% ---------- 颜色别名（章节里用的旧名继续有效） ----------
\colorlet{mainblue}{yhbAccent}
\colorlet{secondcyan}{yhbAccent2}
\colorlet{thirdteal}{yhbAccent2}
\colorlet{highlight}{yhbAccent!6!white}
\colorlet{tablehead}{yhbAccent!12!white}
\definecolor{textgray}{HTML}{4A5057}

% ---------- 课程标签 ----------
\providecommand{\YHBdoclabel}{课程笔记}
% ---------- 正文节奏 ----------
\linespread{1.3}
\setlength{\parindent}{0pt}
\setlength{\parskip}{6pt plus 1.5pt minus 1.5pt}
\setlength{\emergencystretch}{2.5em}
\raggedbottom
\widowpenalty=10000
\clubpenalty=10000
\setlength{\tabcolsep}{7pt}
\renewcommand{\arraystretch}{1.3}
\allowdisplaybreaks[2]
\AtBeginDocument{%
  \setlength{\abovedisplayskip}{10pt plus 2pt minus 2pt}%
  \setlength{\belowdisplayskip}{10pt plus 2pt minus 2pt}%
  \setlength{\abovedisplayshortskip}{6pt}%
  \setlength{\belowdisplayshortskip}{6pt}%
}
\captionsetup{font=small,labelfont=bf,skip=6pt}

% ---------- 列表 ----------
\setlist{leftmargin=2em,labelsep=0.6em,itemsep=4pt,parsep=2pt,topsep=6pt,partopsep=0pt}
\setlist[itemize,2]{label=\textbullet,leftmargin=1.6em}

% ---------- 标题层级 ----------
\ifdefined\YHBbooktrue
  \titleformat{\chapter}[display]
    {\sffamily\bfseries\color{yhbAccent}}
    {\normalsize\sffamily 第\,\thechapter\,章}
    {6pt}
    {\LARGE}
    [{\vspace{5pt}{\color{yhbAccent!55}\titlerule[1pt]}}]
  \titlespacing*{\chapter}{0pt}{14pt}{24pt}
  \titleformat{\section}{\Large\sffamily\bfseries\color{yhbAccent}}{\thesection}{0.7em}{}
  \titlespacing*{\section}{0pt}{18pt plus 3pt minus 2pt}{8pt}
\else
  \titleformat{\section}{\Large\sffamily\bfseries\color{yhbAccent}}{\thesection}{0.7em}{}[{\vspace{3pt}{\color{yhbAccent!45}\titlerule[0.7pt]}}]
  \titlespacing*{\section}{0pt}{18pt plus 3pt minus 2pt}{9pt}
\fi
\titleformat{\subsection}{\large\sffamily\bfseries\color{yhbAccent!85!black}}{\thesubsection}{0.7em}{}
\titleformat{\subsubsection}{\normalsize\bfseries\color{black!85}}{\thesubsubsection}{0.7em}{}
\titlespacing*{\subsection}{0pt}{12pt plus 2pt minus 1pt}{5pt}
\titlespacing*{\subsubsection}{0pt}{9pt plus 2pt minus 1pt}{4pt}
\pretocmd{\section}{\Needspace{6\baselineskip}}{}{}
\pretocmd{\subsection}{\Needspace{5\baselineskip}}{}{}
\pretocmd{\subsubsection}{\Needspace{4\baselineskip}}{}{}

\ifdefined\YHBbooktrue
  \setcounter{secnumdepth}{2}
  \setcounter{tocdepth}{1}
\else
  \setcounter{secnumdepth}{3}
  \setcounter{tocdepth}{2}
\fi
\renewcommand{\contentsname}{目录}
% 目录编号颜色由 titletoc 的 titlecontents 控制，不再重定义 \numberline

% ---------- 目录样式 ----------
\ifdefined\YHBbooktrue
  \titlecontents{chapter}[0pt]
    {\addvspace{6pt}\sffamily\bfseries\color{yhbAccent}}
    {\contentslabel{2.2em}}{}{\titlerule*[0.6em]{.}\contentspage}
  \titlecontents{section}[2.2em]
    {\addvspace{2pt}}
    {\contentslabel{2.2em}}{}{\titlerule*[0.6em]{.}\contentspage}
\else
  \titlecontents{section}[0pt]
    {\addvspace{5pt}\sffamily\bfseries\color{yhbAccent}}
    {\contentslabel{2.2em}}{}{\titlerule*[0.6em]{.}\contentspage}
  \titlecontents{subsection}[2.2em]
    {\addvspace{2pt}}
    {\contentslabel{2.4em}}{}{\titlerule*[0.6em]{.}\contentspage}
\fi

% ---------- 页眉页脚 ----------
\pagestyle{fancy}
\fancyhf{}
\fancyhead[L]{\footnotesize\sffamily\color{textgray}\YHBcourse\ \YHBdoclabel}
\ifdefined\YHBbooktrue
  \fancyhead[R]{\footnotesize\sffamily\color{textgray}\nouppercase{\leftmark}}
  \renewcommand{\chaptermark}[1]{\markboth{#1}{}}
\else
  \fancyhead[R]{\footnotesize\sffamily\color{textgray}\nouppercase{\rightmark}}
  \renewcommand{\sectionmark}[1]{\markright{#1}}
\fi
\fancyfoot[C]{\footnotesize\color{textgray}\thepage}
\renewcommand{\headrulewidth}{0.35pt}
\renewcommand{\headrule}{{\color{yhbAccent!45}\hrule width\headwidth height\headrulewidth}}
\fancypagestyle{plain}{%
  \fancyhf{}%
  \fancyfoot[C]{\footnotesize\color{textgray}\thepage}%
  \renewcommand{\headrulewidth}{0pt}%
}

% ---------- 公式与考点框 ----------
\newcommand{\keyformula}[1]{%
  \begin{tcolorbox}[enhanced,breakable,colback=highlight,colframe=yhbAccent!30!white,boxrule=0.5pt,arc=2.5pt,left=11pt,right=11pt,top=9pt,bottom=9pt,boxsep=0pt,before skip=11pt,after skip=11pt,halign=center,width=\linewidth]
  \adjustbox{max width=\linewidth-2em}{$\displaystyle #1$}
  \end{tcolorbox}}
\newcommand{\corepoint}[1]{%
  \begin{tcolorbox}[enhanced,breakable,colback=highlight,colframe=yhbAccent2,boxrule=0pt,borderline west={2.4pt}{0pt}{yhbAccent2},arc=0pt,left=11pt,right=11pt,top=8pt,bottom=8pt,before skip=11pt,after skip=11pt]
  \textbf{\color{yhbAccent}核心考点}\quad #1
  \end{tcolorbox}}
\newcommand{\symb}[2]{\textbf{#1}：#2\par}

% ---------- 超链接（务必最后加载） ----------
\usepackage[unicode=true,colorlinks=true,linkcolor=yhbAccent,citecolor=yhbAccent,urlcolor=yhbAccent!85!black,pdfborder={0 0 0},bookmarksnumbered=true,pdfstartview=FitH]{hyperref}
\usepackage{hypcap}
\hypersetup{pdftitle={\YHBcourse 课程笔记},pdfauthor={浮萍},pdfsubject={课程学习与复习笔记},pdfcreator={XeLaTeX}}

%   \newcommand{\YHBcourse}{微积分A}      % 课程名（页眉与书签用）
%   \newcommand{\YHBversion}{v1.2}        % 版本号（封面用）
%   \newcommand{\YHBquote}{……}            % 封面题句
%   \definecolor{yhbAccent}{HTML}{24465C} % 主色（深）
%   \definecolor{yhbAccent2}{HTML}{347A80}% 辅色（浅）
%   \YHBbooktrue / \YHBbookfalse          % 是否 book 文档类（有 \chapter）
% 章节文件只写内容，不要重复 \usepackage、不要改字号。
% ============================================================

% ---------- 页面与字体 ----------
\usepackage{geometry}
\geometry{a4paper,left=24mm,right=24mm,top=23mm,bottom=24mm,headheight=15pt,headsep=7mm,footskip=13mm}
\usepackage{microtype}
\usepackage{amsmath,amssymb,amsfonts,mathtools}
\usepackage{esint}
\usepackage{graphicx}
\usepackage{adjustbox}
\usepackage{array,booktabs,multirow,makecell,tabularx,longtable}
\usepackage{xcolor,colortbl}
\usepackage{enumitem}
\usepackage{needspace,etoolbox}
\usepackage{titlesec,titletoc}
\usepackage{fancyhdr}
\usepackage{caption}
\usepackage[most]{tcolorbox}

% ---------- 颜色别名（章节里用的旧名继续有效） ----------
\colorlet{mainblue}{yhbAccent}
\colorlet{secondcyan}{yhbAccent2}
\colorlet{thirdteal}{yhbAccent2}
\colorlet{highlight}{yhbAccent!6!white}
\colorlet{tablehead}{yhbAccent!12!white}
\definecolor{textgray}{HTML}{4A5057}

% ---------- 课程标签 ----------
\providecommand{\YHBdoclabel}{课程笔记}
% ---------- 正文节奏 ----------
\linespread{1.3}
\setlength{\parindent}{0pt}
\setlength{\parskip}{6pt plus 1.5pt minus 1.5pt}
\setlength{\emergencystretch}{2.5em}
\raggedbottom
\widowpenalty=10000
\clubpenalty=10000
\setlength{\tabcolsep}{7pt}
\renewcommand{\arraystretch}{1.3}
\allowdisplaybreaks[2]
\AtBeginDocument{%
  \setlength{\abovedisplayskip}{10pt plus 2pt minus 2pt}%
  \setlength{\belowdisplayskip}{10pt plus 2pt minus 2pt}%
  \setlength{\abovedisplayshortskip}{6pt}%
  \setlength{\belowdisplayshortskip}{6pt}%
}
\captionsetup{font=small,labelfont=bf,skip=6pt}

% ---------- 列表 ----------
\setlist{leftmargin=2em,labelsep=0.6em,itemsep=4pt,parsep=2pt,topsep=6pt,partopsep=0pt}
\setlist[itemize,2]{label=\textbullet,leftmargin=1.6em}

% ---------- 标题层级 ----------
\ifdefined\YHBbooktrue
  \titleformat{\chapter}[display]
    {\sffamily\bfseries\color{yhbAccent}}
    {\normalsize\sffamily 第\,\thechapter\,章}
    {6pt}
    {\LARGE}
    [{\vspace{5pt}{\color{yhbAccent!55}\titlerule[1pt]}}]
  \titlespacing*{\chapter}{0pt}{14pt}{24pt}
  \titleformat{\section}{\Large\sffamily\bfseries\color{yhbAccent}}{\thesection}{0.7em}{}
  \titlespacing*{\section}{0pt}{18pt plus 3pt minus 2pt}{8pt}
\else
  \titleformat{\section}{\Large\sffamily\bfseries\color{yhbAccent}}{\thesection}{0.7em}{}[{\vspace{3pt}{\color{yhbAccent!45}\titlerule[0.7pt]}}]
  \titlespacing*{\section}{0pt}{18pt plus 3pt minus 2pt}{9pt}
\fi
\titleformat{\subsection}{\large\sffamily\bfseries\color{yhbAccent!85!black}}{\thesubsection}{0.7em}{}
\titleformat{\subsubsection}{\normalsize\bfseries\color{black!85}}{\thesubsubsection}{0.7em}{}
\titlespacing*{\subsection}{0pt}{12pt plus 2pt minus 1pt}{5pt}
\titlespacing*{\subsubsection}{0pt}{9pt plus 2pt minus 1pt}{4pt}
\pretocmd{\section}{\Needspace{6\baselineskip}}{}{}
\pretocmd{\subsection}{\Needspace{5\baselineskip}}{}{}
\pretocmd{\subsubsection}{\Needspace{4\baselineskip}}{}{}

\ifdefined\YHBbooktrue
  \setcounter{secnumdepth}{2}
  \setcounter{tocdepth}{1}
\else
  \setcounter{secnumdepth}{3}
  \setcounter{tocdepth}{2}
\fi
\renewcommand{\contentsname}{目录}
% 目录编号颜色由 titletoc 的 titlecontents 控制，不再重定义 \numberline

% ---------- 目录样式 ----------
\ifdefined\YHBbooktrue
  \titlecontents{chapter}[0pt]
    {\addvspace{6pt}\sffamily\bfseries\color{yhbAccent}}
    {\contentslabel{2.2em}}{}{\titlerule*[0.6em]{.}\contentspage}
  \titlecontents{section}[2.2em]
    {\addvspace{2pt}}
    {\contentslabel{2.2em}}{}{\titlerule*[0.6em]{.}\contentspage}
\else
  \titlecontents{section}[0pt]
    {\addvspace{5pt}\sffamily\bfseries\color{yhbAccent}}
    {\contentslabel{2.2em}}{}{\titlerule*[0.6em]{.}\contentspage}
  \titlecontents{subsection}[2.2em]
    {\addvspace{2pt}}
    {\contentslabel{2.4em}}{}{\titlerule*[0.6em]{.}\contentspage}
\fi

% ---------- 页眉页脚 ----------
\pagestyle{fancy}
\fancyhf{}
\fancyhead[L]{\footnotesize\sffamily\color{textgray}\YHBcourse\ \YHBdoclabel}
\ifdefined\YHBbooktrue
  \fancyhead[R]{\footnotesize\sffamily\color{textgray}\nouppercase{\leftmark}}
  \renewcommand{\chaptermark}[1]{\markboth{#1}{}}
\else
  \fancyhead[R]{\footnotesize\sffamily\color{textgray}\nouppercase{\rightmark}}
  \renewcommand{\sectionmark}[1]{\markright{#1}}
\fi
\fancyfoot[C]{\footnotesize\color{textgray}\thepage}
\renewcommand{\headrulewidth}{0.35pt}
\renewcommand{\headrule}{{\color{yhbAccent!45}\hrule width\headwidth height\headrulewidth}}
\fancypagestyle{plain}{%
  \fancyhf{}%
  \fancyfoot[C]{\footnotesize\color{textgray}\thepage}%
  \renewcommand{\headrulewidth}{0pt}%
}

% ---------- 公式与考点框 ----------
\newcommand{\keyformula}[1]{%
  \begin{tcolorbox}[enhanced,breakable,colback=highlight,colframe=yhbAccent!30!white,boxrule=0.5pt,arc=2.5pt,left=11pt,right=11pt,top=9pt,bottom=9pt,boxsep=0pt,before skip=11pt,after skip=11pt,halign=center,width=\linewidth]
  \adjustbox{max width=\linewidth-2em}{$\displaystyle #1$}
  \end{tcolorbox}}
\newcommand{\corepoint}[1]{%
  \begin{tcolorbox}[enhanced,breakable,colback=highlight,colframe=yhbAccent2,boxrule=0pt,borderline west={2.4pt}{0pt}{yhbAccent2},arc=0pt,left=11pt,right=11pt,top=8pt,bottom=8pt,before skip=11pt,after skip=11pt]
  \textbf{\color{yhbAccent}核心考点}\quad #1
  \end{tcolorbox}}
\newcommand{\symb}[2]{\textbf{#1}：#2\par}

% ---------- 超链接（务必最后加载） ----------
\usepackage[unicode=true,colorlinks=true,linkcolor=yhbAccent,citecolor=yhbAccent,urlcolor=yhbAccent!85!black,pdfborder={0 0 0},bookmarksnumbered=true,pdfstartview=FitH]{hyperref}
\usepackage{hypcap}
\hypersetup{pdftitle={\YHBcourse 课程笔记},pdfauthor={浮萍},pdfsubject={课程学习与复习笔记},pdfcreator={XeLaTeX}}

\usepackage{tikz}
\usepackage{circuitikz}
\begin{document}
% ============================================================
% 2.cover.tex —— 统一封面（信息全部来自主文件定义）
% 需要的宏：\YHBcourse、\YHBversion、\YHBquote、yhbAccent、yhbAccent2
% 可选：\YHBcoverSize / \YHBcoverLead（课程名字号与行距）、\YHBcoversub（副标题）
% ============================================================
\providecommand{\YHBcoverSize}{38pt}
\providecommand{\YHBcoverLead}{46pt}
\providecommand{\YHBcoversub}{课程笔记}
\begin{titlepage}
\thispagestyle{empty}
\pdfbookmark[0]{封面}{cover}

\begin{tikzpicture}[remember picture,overlay]
  % 顶部深色块
  \fill[yhbAccent] (current page.north west) rectangle ([yshift=-132mm]current page.north east);
  % 分隔细条（辅色）
  \fill[yhbAccent2] ([yshift=-132mm]current page.north west) rectangle ([yshift=-135mm]current page.north east);
  % 极细横线，弱装饰，不压字
  \foreach \y in {20,34,48,62,76,90,104,118}{%
    \draw[white,opacity=0.08,line width=0.5pt]
      ([xshift=16mm,yshift=-\y mm]current page.north west) --
      ([xshift=-16mm,yshift=-\y mm]current page.north east);}
\end{tikzpicture}

\vspace*{32mm}
\noindent
{\xeCJKsetup{CJKecglue={}}\color{white}\sffamily\bfseries\fontsize{\YHBcoverSize}{\YHBcoverLead}\selectfont \YHBcourse\par}
\vspace{6mm}
{\color{white}\sffamily\fontsize{25}{32}\selectfont \YHBcoversub\par}
\vspace{9mm}
{\color{white}\fontsize{12.5}{18}\selectfont 作者\quad 浮萍\qquad 版本\quad \YHBversion\par}

\vspace{20mm}
{\sffamily\bfseries\fontsize{13}{20}\selectfont\color{yhbAccent}\YHBquote\par}

\vfill
{\color{yhbAccent}\rule{\textwidth}{1.4pt}}\par
\vspace{7mm}
{\fontsize{11.5}{17}\selectfont\sffamily
\textbf{QQ 邮箱}\quad\href{mailto:3071974740@qq.com}{3071974740@qq.com}\par
\vspace{2mm}
\textbf{GitHub}\quad\href{https://github.com/Duckweed-yhb}{github.com/Duckweed-yhb}\par}
\vspace{7mm}
{\fontsize{10.5}{16}\selectfont\color{black!78}
本笔记为个人学习整理，内容仍在持续完善。如发现内容、公式、文字或排版问题，欢迎通过上述 QQ 邮箱反馈；请尽量注明版本、页码与问题描述，便于核对和修订。\par}
\vspace*{12mm}
\end{titlepage}

\newpage
\tableofcontents
\newpage
% ===== 章节目录：在此按教材章节顺序 \input =====
% ============================================================
% 第 1 章  电力系统的基本概念（教材：何仰赞、温增银《电力系统分析》上册 第 1 章）
% 结构对齐教材 1.1—1.5；例题、习题不录
% ============================================================

\chapter{电力系统的基本概念}

电力系统分析研究的是“电能在生产、输送、分配、消费全过程中各运行变量之间的定量关系”。
本章先把对象、单位制（额定值）和网络结构定下来，后面各章的等值电路、潮流、短路、稳定，
都是在这套对象和约定之上展开的。

\section{电力系统的组成}

\subsection{三个层次：动力系统、电力系统、电力网}

\begin{description}[leftmargin=4.5em,style=nextline]
  \item[动力系统] 电力系统加上发电厂的动力部分（锅炉、汽轮机、水轮机、核反应堆、燃气轮机，
        以及燃料输送、除灰、供水等辅助系统）。
  \item[电力系统] 由发电厂、变电所、输电线路、配电线路和用电设备连接而成，完成电能生产、
        输送、分配与消费的统一整体。
  \item[电力网] 电力系统中除去发电机和用电设备以后的部分，即变压器与各级电压线路组成的网络，
        其任务是输送和分配电能。
\end{description}

三者是包含关系：
\[
  \text{动力系统}\;\supset\;\text{电力系统}\;\supset\;\text{电力网}.
\]

\subsection{电能的四个环节}

电能从一次能源到用户，依次经过“发—输—配—用”四个环节：

\begin{enumerate}
  \item \textbf{发电}：火电、水电、核电以及风电、光伏等新能源电源，把一次能源转换为电能。
        发电机端电压一般为 6～20~kV，需经升压变压器升至电网电压。
  \item \textbf{输电}：由最高一级或两级电压的线路构成，把大容量电源的电能送到负荷中心，
        我国常见为 220~kV、330~kV、500~kV、750~kV（交流）及直流输电。
  \item \textbf{配电}：把电能分配到各用户，电压一般为 0.4～35~kV，
        负荷密度大的城市也用 110~kV、220~kV 配电。
  \item \textbf{用电}：异步电动机、同步电动机、电热装置、整流装置、照明设备等。
\end{enumerate}

\begin{table}[htbp]
  \centering
  \caption{电力系统各环节与主要设备}
  \small
  \begin{tabular}{@{}llll@{}}
    \toprule
    环节 & 主要设备 & 典型电压 & 分析中的处理方式 \\
    \midrule
    发电 & 同步发电机、励磁系统 & 6～20~kV & 电势源 $+$ 电抗（第 3 章） \\
    升压/降压 & 双绕组、三绕组、自耦变压器 & 各级 & 阻抗 $+$ 理想变压器（第 2 章） \\
    输电 & 架空线路、电缆 & 110～1000~kV & $\pi$ 型等值电路（第 2 章） \\
    配电 & 架空线路、电缆、配电变压器 & 0.4～35~kV & 简化等值电路 \\
    用电 & 异步电动机为主的综合负荷 & 0.38~kV 等 & 恒定功率或阻抗型负荷（第 9 章） \\
    \bottomrule
  \end{tabular}
\end{table}

\subsection{电力系统运行的三个特点}

\begin{enumerate}
  \item \textbf{电能不能大量存储}。电能的生产、输送、分配和消费几乎同时进行，
        任一时刻发电机发出的有功功率必须等于该时刻负荷所需功率与网络损耗之和。
        这一“功率实时平衡”的约束，是潮流计算、频率调整和稳定问题的共同起点。
  \item \textbf{暂态过程非常短促}。电磁暂态以毫秒计（短路后 $5\sim 10$~ms 内电流即达冲击值），
        机电暂态以秒计。故分析中必须区分电磁暂态与机电暂态，采用不同的数学模型。
  \item \textbf{与国民经济各部门关系极为密切}。供电突然中断后果严重，
        这决定了可靠性要求高于单纯的经济性考虑。
\end{enumerate}

\corepoint{“电能不可存储”是全书的第一性约束：负荷变化必然引起频率与电压变化，
而频率靠有功平衡、电压靠无功平衡来维持。凡是讨论“平衡”的章节（第 11、12、13 章），
根子都在这一条特点上。}

\section{电力系统的额定电压和额定频率}

\subsection{额定电压的定义与约定}

额定电压是电气设备设计和制造的基本技术参数，也是衡量电能质量的基本指标之一。
我国规定，\textbf{电力网的额定电压}就是线路首末两端用电设备所要求的电压，
以 $U_N$ 表示。各类设备的额定电压必须与电网额定电压相适应，但又不完全相同，
原因是电压沿线路有降落、变压器内部有阻抗压降。

\begin{table}[htbp]
  \centering
  \caption{各类电气设备额定电压的确定}
  \small
  \begin{tabular}{@{}p{0.17\linewidth}p{0.34\linewidth}p{0.36\linewidth}@{}}
    \toprule
    设备 & 额定电压的取法 & 理由 \\
    \midrule
    用电设备 & $U_N$（等于所在电网额定电压） & 用电设备直接接于电网，按电网电压设计 \\
    发电机 & $1.05\,U_N$ & 发电机接在线路首端，需补偿线路上的电压降落 \\
    变压器一次绕组 & 等于所接电网的额定电压；\newline 直接与发电机相连时取 $1.05\,U_N$ & 一次绕组相当于用电设备 \\
    变压器二次绕组 & $1.1\,U_N$（短路电压较小或线路很短时可取 $1.05\,U_N$） & 二次额定电压指\textbf{空载}电压；带额定负荷时内部压降约 $5\%$，故空载须高 $10\%$ \\
    \bottomrule
  \end{tabular}
\end{table}

\keyformula{\Delta U\% = \frac{U-U_N}{U_N}\times 100\%,\qquad
            \delta f\% = \frac{f-f_N}{f_N}\times 100\%}

变压器高压绕组常设有分接头，主抽头为 $U_N$，其余为 $\pm 5\%$ 或 $\pm 2\times 2.5\%$
（即 $\pm 5\%$、$\pm 2.5\%$ 若干档），用以在运行中改变变比调压（第 12 章）。
\textbf{实际变比}应取运行时两侧实际抽头的空载线电压之比，而不是铭牌额定电压之比。

\subsection{电压等级与平均额定电压}

我国采用的交流电压等级（线电压有效值）为
\[
  3,\;6,\;10,\;35,\;60,\;110,\;220,\;330,\;500,\;750\ \text{kV},
\]
其中 330~kV 主要用于西北电网，750~kV 用于超高压外送。电压等级按 $\sqrt{3}$ 倍左右的
几何级数配置，使相邻两级之间的输送能力匹配。

在标幺制短路计算中（第 6 章）需要统一基准电压。工程上取各级电网的\textbf{平均额定电压}
$U_{av}\approx 1.05U_N$ 作为基准电压，我国规定为
\begin{equation}
  U_{av}=3.15,\;6.3,\;10.5,\;15.75,\;37,\;115,\;230,\;345,\;525,\;1050\ \text{kV}.
\end{equation}
取平均额定电压作基准后，变压器的额定电压比与基准电压比恰好相等，
变压器的标幺变比等于 $1$，等值电路中理想变压器可以消去——这是简化计算的依据。

\subsection{额定频率与电能质量指标}

我国电力系统的额定频率为 $50$~Hz（工频）。正常运行时允许的频率偏移为
$\pm 0.2\sim \pm 0.5$~Hz（大系统取严、小系统取宽）。供电电压的允许偏移为：
\begin{itemize}
  \item 35~kV 及以上电压级：$\pm 5\%$；
  \item 10~kV 及以下电压级：$\pm 7\%$；
  \item 电压正弦波形畸变率（各次谐波有效值平方和的方根值与基波有效值之比）：
        6～10~kV 供电电压不超过 $4\%$，0.38~kV 不超过 $5\%$。
\end{itemize}

\begin{table}[htbp]
  \centering
  \caption{电能质量的主要指标}
  \small
  \begin{tabular}{@{}lll@{}}
    \toprule
    指标 & 允许值 & 主要影响因素 \\
    \midrule
    频率偏移 & $\pm 0.2\sim\pm 0.5$~Hz & 系统有功功率平衡 \\
    电压偏移 & 35~kV 及以上 $\pm 5\%$；10~kV 及以下 $\pm 7\%$ & 系统无功功率平衡与网络压降 \\
    波形畸变率 & 6～10~kV $\le 4\%$；0.38~kV $\le 5\%$ & 非线性负荷（整流、变频、电弧） \\
    三相不平衡度 & 按国标规定 & 单相负荷、不对称故障 \\
    \bottomrule
  \end{tabular}
\end{table}

\corepoint{频率取决于\textbf{有功}功率平衡，电压取决于\textbf{无功}功率平衡。
系统有功不足则频率偏低，无功不足则电压偏低。因此“保证电能质量”这件事，
在分析上就落实为两个功率平衡问题加两套调整手段（调频：第 13 章；调压：第 12 章）。}

\section{对电力系统运行的基本要求}

对电力系统运行的基本要求可概括为四条：
\begin{enumerate}
  \item 保证安全可靠地供电；
  \item 有合乎要求的电能质量；
  \item 有良好的经济性；
  \item 尽可能减小对生态环境的有害影响。
\end{enumerate}

\subsection{可靠性要求与负荷分级}

按用户对供电可靠性的不同要求，我国将负荷分为三级：

\begin{table}[htbp]
  \centering
  \caption{负荷分级及对供电的要求}
  \small
  \begin{tabular}{@{}p{0.10\linewidth}p{0.45\linewidth}p{0.35\linewidth}@{}}
    \toprule
    级别 & 中断供电的后果 & 供电要求 \\
    \midrule
    第一级 & 危及人身安全；关键设备难以修复的损坏；生产秩序长期不能恢复；
             市政生活重要部门混乱 & 由两个独立电源供电，任一电源故障时另一电源应能保证供电 \\
    第二级 & 大量减产；城市大量居民正常活动受影响 & 一般由双回路供电，条件许可时设备用电源 \\
    第三级 & 不属于第一、二级，短时停电不造成重大损失 & 可由单回路无备用网络供电 \\
    \bottomrule
  \end{tabular}
\end{table}

提高可靠性的措施是“两手”：一是配备足够的有功、无功电源备用，完善网络结构以提高
抵抗扰动的能力（即提高稳定性）；二是严密监视与维修设备、提高运行水平，
并借助计算机进行安全监视与控制。

\subsection{经济性要求}

经济性集中在三个损耗率指标上：
\begin{itemize}
  \item 发电厂的煤耗率（水耗率）——反映发电环节的效率；
  \item 厂用电率——反映发电厂自身用电所占比重；
  \item 电力网的损耗率（网损率）——反映输配电环节的效率。
\end{itemize}
降低网损的措施（提高电压等级、升高运行电压、无功就地补偿、合理选择导线截面、
经济调度等）将在第 14 章系统讨论。

\subsection{电能质量与功率平衡的关系}

电压和频率是电气设备设计与制造的基本参数，也是衡量电能质量的两个基本指标。
它们超差不仅造成废品和减产，还会影响设备安全，严重时危及整个系统。
而要保证电能质量，关键在于“功率平衡 + 配置得当 + 有调整手段”：
\begin{itemize}
  \item 有功功率在额定频率下平衡 $\Rightarrow$ 频率合格；
  \item 无功功率在额定电压下平衡 $\Rightarrow$ 电压合格；
  \item 对谐波污染源进行限制和治理 $\Rightarrow$ 波形合格。
\end{itemize}

\section{电力系统的接线方式}

电力系统的接线包括发电厂主接线、变电所主接线和电力网接线。本节只讨论\textbf{电力网}的
接线方式，它按供电可靠性分为无备用与有备用两大类。

\subsection{无备用接线（开式网络）}

每个负荷只能沿\textbf{唯一路径}取得电能的网络称为\textbf{开式网络}，包括：
\begin{itemize}
  \item \textbf{放射式}：电源点经一条线路直接引至一个负荷点，故障影响范围小、控制方便；
  \item \textbf{干线式}：一条干线沿途分支供电，节省线路走廊；
  \item \textbf{树状}：多级分支，结构最简单。
\end{itemize}
共同优点：简单、设备费用少、运行方便。共同缺点：可靠性低，任一段线路故障或检修
都要中断部分用户供电；干线式和树状网络线路较长时，末端电压往往偏低。

\subsection{有备用接线（闭式网络）}

\begin{itemize}
  \item \textbf{双回路放射式、双回路干线式}：在无备用网络的每一段线路上改用双回路。
        简单、运行方便，可靠性和电压质量明显提高，缺点是设备费用增加很多。
  \item \textbf{环形网络}：由一个或几个电源点与一个或几个负荷点经线路连成环。
        可靠性令人满意且比较经济，缺点是运行调度复杂；单电源环网中一段线路故障开环时，
        正常线段可能过负荷，负荷点电压明显降低。
  \item \textbf{两端供电网络}：可靠性相当于有两个电源的环形网络。
\end{itemize}

环形网络和两端供电网络中，每个负荷点至少经两条线路从不同方向取得电能，
这类网络统称\textbf{闭式网络}。

\begin{table}[htbp]
  \centering
  \caption{两类接线方式的比较}
  \small
  \begin{tabular}{@{}p{0.14\linewidth}p{0.22\linewidth}p{0.22\linewidth}p{0.28\linewidth}@{}}
    \toprule
    接线方式 & 可靠性 & 经济性 & 适用场合 \\
    \midrule
    无备用（放射、干线、树状） & 低，故障即停电 & 最好，设备最省 & 第三级负荷 \\
    双回路 & 较高 & 较差，线路投资大 & 第一、二级负荷比重大的用户 \\
    环网 & 高 & 较好 & 负荷中心的大电厂与枢纽变电所互联 \\
    两端供电 & 高（相当于双电源环网） & 较好 & 有两个独立电源的负荷区 \\
    \bottomrule
  \end{tabular}
\end{table}

\subsection{输电网络与配电网络}

\begin{itemize}
  \item \textbf{输电网络}：任务是把大容量发电厂的电能可靠而经济地送到负荷集中地区，
        通常由系统中最高的一级或两级电压线路组成，连接区域发电厂（经升压站）与枢纽变电所。
        要求：足够的可靠性、满足稳定要求、便于经济调度、对运行方式变更和系统发展有适应性。
        远离负荷中心的大电厂输电干线、向缺乏电源地区供电的干线常采用双回路或多回路；
        位于负荷中心的大电厂和枢纽变电所一般通过环网互联。
  \item \textbf{配电网络}：任务是分配电能，电压一般为 0.4～35~kV，
        大城市负荷密度大时也用 110~kV、220~kV。电源点是相应电压等级的母线，
        负荷点是低一级变电所或用电设备。
\end{itemize}

选择接线方式时须综合考虑：满足用户对可靠性和电压质量的要求、运行灵活方便、
良好的经济指标；一般要对多个可行方案做技术经济比较后确定。

表~\ref{tab:capacity} 给出各级电压单回架空线路输送容量与输送距离的适宜范围，
它反映了“电压等级要与系统规模（容量和供电范围）相适应”这一基本原则。

\begin{table}[htbp]
  \centering
  \caption{各级电压架空线路的输送能力（教材表 1-2）}
  \label{tab:capacity}
  \small
  \begin{tabular}{@{}cccc@{}}
    \toprule
    额定电压/kV & 输送容量/(MV$\cdot$A) & 额定电压/kV & 输送容量/(MV$\cdot$A) \\
    \midrule
    \multicolumn{2}{c}{（输送距离/km）} & \multicolumn{2}{c}{（输送距离/km）} \\
    3   & $0.1\sim1.0$\quad（$1\sim3$）      & 110 & $10\sim50$\quad（$50\sim150$） \\
    6   & $0.1\sim1.2$\quad（$4\sim15$）     & 220 & $100\sim500$\quad（$100\sim300$） \\
    10  & $0.2\sim2$\quad（$6\sim20$）       & 330 & $200\sim800$\quad（$200\sim600$） \\
    35  & $2\sim10$\quad（$20\sim50$）       & 500 & $1000\sim1500$\quad（$150\sim850$） \\
    60  & $3.5\sim30$\quad（$30\sim100$）    & 1000 & $3500\sim5000$\quad（$1000$ 以上） \\
    \bottomrule
  \end{tabular}
\end{table}

\corepoint{无备用接线只允许向第三级负荷供电；第一、二级负荷占较大比重的用户必须由
有备用网络供电。答题时若问“某接线方式是否可行”，先判负荷级别，再判是否满足
$N-1$（任一元件退出仍能供电）。}

\section{电力系统分析课程的主要内容}

\subsection{运行状态与运行变量}

当系统中各种发电、变电、输配电及用电设备之间的联接情况已经确定时，
系统的运行状态由一组\textbf{运行变量}（运行参数）描述：功率、频率、电压、电流、
磁链、电动势以及发电机转子间的相对位移角等。

\begin{table}[htbp]
  \centering
  \caption{稳态与暂态的划分}
  \small
  \begin{tabular}{@{}lll@{}}
    \toprule
    状态 & 含义 & 典型分析内容 \\
    \midrule
    正常稳态 & 三相对称运行，参数在恒定平均值附近微小变化 & 潮流、无功/有功平衡、经济运行 \\
    故障稳态 & 短路进入稳态后的状态（对称或不对称） & 短路电流实用计算 \\
    故障后稳态 & 故障切除后的新稳态 & 稳定分析中的初始点 \\
    暂态 & 从一种运行状态到另一种状态的过渡 & 短路暂态过程、暂态稳定 \\
    振荡 & 运行参数周期性变化的运行状态 & 低频振荡、静态稳定 \\
    \bottomrule
  \end{tabular}
\end{table}

严格说，由于负荷等随机扰动，绝对稳态并不存在；工程上把运行参数只在某恒定平均值附近
作微小变化的状态视为稳态。

\subsection{数学模型：运行变量与系统参数}

分析运行状态的基本方法是：把待研究的状态用数学方程式描述出来，再用适当的数学方法
和计算工具求解。方程式中的未知量是运行变量，系数是\textbf{系统参数}。

系统参数指各元件或其组合在运行中反映其物理特性的参数：电阻、电感（电抗）、
电容（电纳）、时间常数、变压器变比、输入阻抗、转移阻抗等。参数主要取决于元件结构，
也与额定参数密切相关。

以 $n$ 个节点的网络为例，其节点方程为
\keyformula{\dot I_i=\sum_{j=1}^{n}Y_{ij}\dot U_j,\qquad i=1,2,\dots,n}
其中 $Y_{ij}$ 为节点导纳矩阵元素（系统参数），$\dot U_j$ 为节点电压（运行变量）。
这一个式子同时是第 4 章（网络数学模型）、第 6 章（短路电流计算）和第 11 章（潮流计算）
的出发点，区别只在于已知量与未知量的选择不同。

\subsection{上、下册的两条主线}

\begin{table}[htbp]
  \centering
  \caption{课程内容的两条主线}
  \small
  \begin{tabular}{@{}p{0.10\linewidth}p{0.22\linewidth}p{0.28\linewidth}p{0.30\linewidth}@{}}
    \toprule
    分册 & 基本物理量 & 元件处理方式 & 核心问题 \\
    \midrule
    上册（第 1～8 章） & 电流 & 元件当作电路元件，用等值电路代替 & 短路电流计算：各序电抗的物理意义与计算、标幺参数归算、
                                           输入阻抗与转移阻抗、节点导纳/阻抗矩阵 \\
    下册（第 9～19 章） & 功率 & 发电与用电设备从功率转换角度看特性；
                          电网元件着重电压降落与功率损耗 & 稳态：潮流（核心）、无功与电压、有功与频率、经济运行；
                                                                     稳定：功角、静态与暂态稳定、提高稳定的措施 \\
    \bottomrule
  \end{tabular}
\end{table}

上册的关键在\textbf{阻抗（工程计算中主要是电抗）的计算}；下册的关键在\textbf{功率的分布}。
两者共用同一套网络方程，只是把“电流激励—电压响应”换成了“功率注入—电压分布”。

\corepoint{记课程地图：第 2 章给参数、第 3 章给发电机方程、第 4 章给网络方程，
这三章合起来是“建模”；第 5～8 章是“短路”（上册主线）；第 9～14 章是“稳态运行”
（潮流为核）；第 15～19 章是“稳定”。考试大题的骨架通常就是“建模 $\to$ 计算 $\to$ 判据”。}
          % 第 1 章  电力系统的基本概念
% ============================================================
% 第 2 章  电力网络各元件的等值电路和参数计算
% （教材：何仰赞、温增银《电力系统分析》上册 第 2 章）
% 结构对齐 2.1—2.4；例题、习题不录
% ============================================================

\chapter{电力网络各元件的等值电路和参数计算}

电力系统的运行状态基本上是三相对称的（正常运行），或者可以化为三相对称的（用对称分量法），
因此只需研究一相。元件的三相有星形和三角形接法，等值电路相应地也有星形和三角形两种，
为便于用一相等值电路计算，常把三角形化为星形。\textbf{等值电路中的参数是计及了其余两相
影响（相间互感等）的一相等值参数}，这一点是本章所有算式的共同前提。

\section{架空输电线路的参数}

线路的四个参数各有明确物理意义：电阻反映有功损耗，电感反映磁场效应，电导反映泄漏电流
与电晕引起的有功损耗，电容反映电场效应。参数沿全长均匀分布，单位长度参数记为
$r_0$、$L_0$（或 $x_0$）、$g_0$、$C_0$（或 $b_0$）。电缆参数由厂家给出或实测，本节只讨论架空线。

\subsection{电阻}

有色金属导线单位长度的直流电阻
\begin{equation}
  r=\rho/S ,
\end{equation}
其中 $\rho$ 为电阻率（$\Omega\cdot\text{mm}^2/\text{km}$），$S$ 为载流部分标称截面积（$\text{mm}^2$）。
工程上不用材料的标准电阻率，而用略为增大的\textbf{计算值}：铜 $18.8$、铝 $31.5\
\Omega\cdot\text{mm}^2/\text{km}$。原因有三：
\begin{enumerate}
  \item 工频交流下集肤效应与邻近效应使交流电阻略大于直流电阻；
  \item 多股绞线扭绞后实际长度比线路长度长 $2\%\sim3\%$；
  \item 制造中实际截面常比标称截面略小。
\end{enumerate}

手册查得或按上式计算的值均为 $20\,^\circ$C 时的值，其他温度下
\keyformula{r_t=r_{20}\left[1+\alpha(t-20)\right]}
其中温度系数 $\alpha$：铜 $0.00382/^\circ$C，铝 $0.00361/^\circ$C。
$220$~kV 以下线路一般可忽略集肤效应，直接取手册值。

\subsection{电感（电抗）}

\subsubsection{基本算式}

设磁路磁导率为常数，则磁链与电流成正比，自感 $L=\varPsi/i$，互感 $M_{AB}=\varPsi_{AB}/i_B$。
非铁磁材料圆柱形长导线（长度 $l\gg$ 半径 $r$）每单位长度的自感
\begin{equation}
  L=\frac{\mu_0}{2\pi}\left(\ln\frac{2l}{r}-1\right)=\frac{\mu_0}{2\pi}\ln\frac{2l}{D_s},
  \qquad D_s=re^{-1/4}=0.779\,r ,
\end{equation}
$D_s$ 称为导线的\textbf{自几何均距}。两根平行长导线（轴线间距 $D$）每单位长度的互感
\begin{equation}
  M=\frac{\mu_0}{2\pi}\left(\ln\frac{2l}{D}-1\right)=\frac{\mu_0}{2\pi}\ln\frac{2l}{D}.
\end{equation}
这两式是计算多相线路电感的基础。

\subsubsection{三相输电线路的一相等值电感}

等边三角形排列、相间距离 $D$ 的三相线路，利用 $i_a+i_b+i_c=0$ 消去互感项后得
\begin{equation}
  L_1=\frac{\mu_0}{2\pi}\ln\frac{D}{D_s}\quad(\text{H/m}).
\end{equation}
三相排列不对称时，各相交链磁链不等，须用\textbf{导线换位}（一个整循环换位分三段，
每段三相各占一个位置）使各相平均电感接近相等。经整循环换位后
\begin{equation}
  L_1=\frac{\mu_0}{2\pi}\ln\frac{D_{eq}}{D_s},
  \qquad D_{eq}=\sqrt[3]{D_{ab}D_{bc}D_{ca}},
\end{equation}
$D_{eq}$ 称为三相导线的\textbf{互几何均距}。三相水平排列时 $D_{eq}=\sqrt[3]{D\cdot 2D\cdot D}=1.26D$。

导线自几何均距的取值：非铁磁材料单股线 $D_s=0.779r$；多股线 $D_s=(0.724\sim0.771)r$；
钢芯铝线 $D_s=(0.77\sim0.9)r$（查手册）。

\subsubsection{分裂导线的等值电感}

每相分裂成 $n$ 根（一般 $n\le4$），布置在正多边形顶点上，边长（分裂间距）为 $d$。
因相间距离远大于 $d$，可认为各相分裂导线重心间的距离即相间距离，于是只需用分裂导线的
自几何均距 $D_{sb}$ 代替 $D_s$：
\begin{equation}
  L_1=\frac{\mu_0}{2\pi}\ln\frac{D_{eq}}{D_{sb}},\qquad
  D_{sb}=\sqrt[n]{D_s\,d^{\,n-1}} .
\end{equation}
具体写开为
\begin{align}
  n=2:&\quad D_{sb}=\sqrt{D_s d},\\
  n=3:&\quad D_{sb}=\sqrt[3]{D_s d^{2}},\\
  n=4:&\quad D_{sb}=\sqrt[4]{D_s d^{3}}=1.09\sqrt[4]{D_s d^{3}} .
\end{align}
$d\gg D_s$，故 $D_{sb}$ 大于单导线的 $D_s$，分裂导线等值电感较小。

\subsubsection{等值电抗}

额定频率 $f_N=50$~Hz、$\mu_0=4\pi\times10^{-7}$~H/m 时，单导线线路
\keyformula{x_1=0.0628\ln\frac{D_{eq}}{D_s}=0.1445\lg\frac{D_{eq}}{D_s}\quad(\Omega/\text{km})}
分裂导线线路把 $D_s$ 换成 $D_{sb}$：
\begin{equation}
  x_1=0.1445\lg\frac{D_{eq}}{D_{sb}}\quad(\Omega/\text{km}).
\end{equation}
相间距离、导线截面等结构参数都在对数符号内，故各种线路的电抗值变化不大：
单导线约 $0.4\ \Omega/\text{km}$；$2$、$3$、$4$ 分裂时分别约 $0.33$、$0.30$、$0.28\ \Omega/\text{km}$。
钢导线的集肤效应与内部磁导率随电流变化，其阻抗是非线性的，只能用实验测定，不能用上式计算。

\subsection{电导}

电导反映泄漏电流和空气游离引起的有功损耗。绝缘良好时泄漏可忽略，主要考虑\textbf{电晕}：
导线表面电场强度超过空气击穿强度时，附近空气游离而产生局部放电（嘶嘶声、臭氧、夜间紫晕）。
开始出现电晕的电压称为临界电压，三相导线等边三角形排列时
\begin{equation}
  U_{cr}=49.3\,m_1 m_2 \delta\, r\lg\frac{D}{r}\quad(\text{kV}),
\end{equation}
其中 $m_1$ 为导线表面状况系数（多股绞线 $0.83\sim0.87$），$m_2$ 为气象系数（干燥晴朗取 $1$，
雨雪雾 $0.8\sim1$），$r$ 为导线计算半径（cm），$\delta$ 为空气相对密度
\begin{equation}
  \delta=\frac{3.92p}{273+t},
\end{equation}
$p$ 为大气压力（Pa），$t$ 为温度（$^\circ$C）；$t=25\,^\circ$C、$p=76$~Pa 时 $\delta=1$。
水平排列时，两边线临界电压比上式计算值高 $6\%$，中间相低 $4\%$。

若三相线路每公里电晕损耗为 $\Delta P_g$（MW/km），则每相等值电导
\begin{equation}
  g_1=\frac{\Delta P_g}{U_L^{2}}\times10^{-3}\quad(\text{S/km}),
\end{equation}
$U_L$ 为线电压（kV）。

设计时总尽量避免正常气象下发生电晕。由 $U_{cr}$ 算式可见：$D$ 在对数内，增大它对
$U_{cr}$ 影响不大却大幅增加杆塔造价；$U_{cr}$ 差不多与 $r$ 成正比，故\textbf{增大导线半径}
是减小电晕损耗的有效方法。220~kV 以下按避免电晕选导线半径；220~kV 及以上用分裂导线
（或扩径导线）增大每相等值半径。因此一般系统计算中可认为 $g\approx0$。

\subsection{电容（电纳）}

由带电平行长导线的电位叠加可得导线 $A$ 的对地电位
$U_A=\dfrac{q}{2\pi\varepsilon}\ln\dfrac{D}{r}$，据此两线间单位长度电容
$C=\pi\varepsilon/\ln(D/r)$。三相线路须计及大地影响（用镜像代替），经整循环换位并略去
与大地影响有关的次要项后，每相等值电容
\begin{equation}
  C_1=\frac{0.0241}{\lg\dfrac{D_{eq}}{r}}\times10^{-6}\quad(\text{F/km}),
\end{equation}
分裂导线用一相导线组的\textbf{等值半径} $r_{eq}$ 代替 $r$：
\begin{equation}
  r_{eq}=\sqrt[n]{r\,d^{\,n-1}},
\end{equation}
即 $n=2$ 时 $r_{eq}=\sqrt{rd}$，$n=3$ 时 $\sqrt[3]{rd^{2}}$，$n=4$ 时 $1.09\sqrt[4]{rd^{3}}$，
单导线 $r_{eq}=r$。

额定频率下的等值电纳
\keyformula{b_1=2\pi f_N C_1=\frac{7.58}{\lg\dfrac{D_{eq}}{r_{eq}}}\times10^{-6}\quad(\text{S/km})}
单导线约 $2.8\times10^{-6}$~S/km；2、3、4 分裂分别约 $3.4$、$3.8$、$4.1\times10^{-6}$~S/km。
与电抗同理，结构参数在对数内，故各电压等级线路电纳差别不大。

\begin{table}[htbp]
  \centering
  \caption{架空线路四个参数的算式与数量级汇总}
  \small
  \begin{tabular}{@{}p{0.08\linewidth}p{0.42\linewidth}p{0.18\linewidth}p{0.22\linewidth}@{}}
    \toprule
    参数 & 算式 & 数量级 & 物理意义 \\
    \midrule
    $r$ & $\rho/S$，$\rho$ 取计算值（铜 18.8、铝 31.5） & $0.05\sim0.2\ \Omega$/km & 有功损耗 \\
    $x$ & $0.1445\lg(D_{eq}/D_s)$（分裂用 $D_{sb}$） & $0.28\sim0.42\ \Omega$/km & 磁场效应 \\
    $g$ & $\Delta P_g/U_L^{2}$，正常可忽略 & $\approx0$ & 电晕与泄漏损耗 \\
    $b$ & $7.58\times10^{-6}/\lg(D_{eq}/r_{eq})$ & $2.8\sim4.1\ \mu$S/km & 电场效应 \\
    \bottomrule
  \end{tabular}
\end{table}

\corepoint{三个易错点：（1） $D_{eq}$ 用互几何均距 $1.26D$（水平排列），不是 $D$；
（2） 分裂导线只改自几何均距 $D_{sb}$（电抗）与等值半径 $r_{eq}$（电纳），$D_{eq}$ 不变；
（3） $r_{eq}$ 用于电纳，$D_{sb}$ 用于电抗，两者虽然形式相似但一个是“半径”一个是“均距”，不能互串。}

\section{架空输电线的等值电路}

\subsection{输电线路的方程式}

设单位长度阻抗 $z_0=r_0+j\omega L_0$、导纳 $y_0=g_0+j\omega C_0$，沿线均匀分布。
在距末端 $x$ 处取微段 $\mathrm{d}x$，由电压降和分流关系得
\begin{equation}
  \frac{\mathrm{d}\dot U}{\mathrm{d}x}=z_0\dot I,\qquad
  \frac{\mathrm{d}\dot I}{\mathrm{d}x}=y_0\dot U ,
\end{equation}
对第一式求导并代入第二式，得二阶常系数齐次微分方程
$\dfrac{\mathrm{d}^{2}\dot U}{\mathrm{d}x^{2}}=z_0y_0\dot U$，通解为
$\dot U=A_1e^{\gamma x}+A_2e^{-\gamma x}$。式中
\keyformula{\gamma=\sqrt{z_0y_0}=\beta+j\alpha,\qquad Z_c=\sqrt{z_0/y_0}=|Z_c|e^{j\theta_c}}
$\gamma$ 称\textbf{传播常数}（实部 $\beta$ 为衰减常数，虚部 $\alpha$ 为相位常数，
因 $z_0$、$y_0$ 的辐角均在 $0^\circ\sim90^\circ$，故 $\beta$、$\alpha$ 均为正）；
$Z_c$ 称\textbf{波阻抗}（特性阻抗），二者只与线路参数和频率有关。
高压架空线 $r_0\ll\omega L_0$、$g_0\approx0$，故 $Z_c$ 接近纯电阻而略呈电容性；
略去 $r_0$、$g_0$ 时 $\gamma=j\alpha=j\omega\sqrt{L_0C_0}$、$Z_c=R_c=\sqrt{L_0/C_0}$。
单导线架空线波阻抗 $370\sim410\ \Omega$，分裂导线 $270\sim310\ \Omega$，电缆 $30\sim50\ \Omega$。

由边界条件（$x=0$ 处 $\dot U=\dot U_2$、$\dot I=\dot I_2$）定出积分常数，写成双曲函数形式：
\begin{align}
  \dot U&=\dot U_2\operatorname{ch}\gamma x+\dot I_2Z_c\operatorname{sh}\gamma x,\\
  \dot I&=\frac{\dot U_2}{Z_c}\operatorname{sh}\gamma x+\dot I_2\operatorname{ch}\gamma x .
\end{align}
取 $x=l$ 得首末端关系，与二端口网络通用方程
$\dot U_1=A\dot U_2+B\dot I_2$、$\dot I_1=C\dot U_2+D\dot I_2$ 比较，有
\begin{equation}
  A=D=\operatorname{ch}\gamma l,\qquad B=Z_c\operatorname{sh}\gamma l,\qquad
  C=\frac{\operatorname{sh}\gamma l}{Z_c},
\end{equation}
且 $AD-BC=1$，故输电线是\textbf{对称的无源二端口网络}，可用对称等值电路表示。

\subsection{集中参数等值电路}

\subsubsection{$\pi$ 型与 T 型}

$\pi$ 型电路（两端各并 $Y'/2$、中间串 $Z'$）与 T 型电路均可作等值电路，其精确参数为
\begin{align}
  \pi\ \text{型}:&\quad Z'=Z_c\operatorname{sh}\gamma l,\qquad
     \frac{Y'}{2}=\frac{\operatorname{ch}\gamma l-1}{Z_c\operatorname{sh}\gamma l},\\
  \text{T 型}:&\quad Z_T'=Z_c\operatorname{th}\gamma l,\qquad Y_T'=\frac{\operatorname{sh}\gamma l}{Z_c}.
\end{align}
工程上绝大多数采用 $\pi$ 型。

\subsubsection{修正系数}

复数双曲函数计算不便，令 $Z=z_0l$、$Y=y_0l$ 为全线总阻抗与总导纳，则
\keyformula{Z'=K_zZ,\qquad Y'=K_yY,\qquad
  K_z=\frac{\operatorname{sh}\gamma l}{\gamma l},\quad
  K_y=\frac{2(\operatorname{ch}\gamma l-1)}{\gamma l\operatorname{sh}\gamma l}}
略去电导并作级数展开，得常用的\textbf{简化修正公式}
\begin{equation}
  Z'=\left(k_r r_0+jk_x x_0\right)l,\qquad Y'\approx jk_b b_0 l,
\end{equation}
\begin{align}
  k_r&=1-\tfrac{1}{3}x_0b_0l^{2},\\
  k_x&=1-\tfrac{1}{6}\left(x_0b_0-\frac{r_0^{2}b_0}{x_0}\right)l^{2},\\
  k_b&=1+\tfrac{1}{12}x_0b_0l^{2}.
\end{align}
可见 $k_r<1$、$k_x<1$ 而 $k_b>1$：修正使等值电阻、电抗略减、电纳略增，
这正是分布参数“沿线逐渐泄漏”的集总化补偿。

\subsubsection{三种参数与适用范围}

\begin{table}[htbp]
  \centering
  \caption{$\pi$ 型电路三种参数取法及适用长度}
  \small
  \begin{tabular}{@{}p{0.13\linewidth}p{0.30\linewidth}p{0.44\linewidth}@{}}
    \toprule
    取法 & 算式 & 适用与说明 \\
    \midrule
    近似参数 & $Z'=Z$,\ $Y'=Y$ & 长度不超过 $300$~km 可用一个 $\pi$ 型电路；更长时用串级的多个 $\pi$ 型，
                                    每个代替 $200\sim300$~km \\
    修正参数 & $Z'=(k_rr_0+jk_xx_0)l$,\ $Y'=jk_bb_0l$ & 一个 $\pi$ 型可代替 $500\sim600$~km 的线路，
                                    精度已接近精确参数 \\
    精确参数 & $Z'=K_zZ$,\ $Y'=K_yY$ & 双曲函数计算，用于长线校核；上述处理仅适用于工频稳态计算 \\
    \bottomrule
  \end{tabular}
\end{table}

误差规律：随线路长度增加而增大；相对而言电阻误差最大，电抗次之，电纳最小。
修正后各参数与精确值已非常接近，修正效果显著。即使 $g_0=0$，精确参数中仍有一个数值很小的
电导，实际计算可忽略。

\corepoint{短线路（$\le100$~km、$U_N\le110$~kV）可只保留串联阻抗、忽略并联电纳；
中等长度（$100\sim300$~km）用 $\pi$ 型近似参数；长线（$>300$~km）必须用修正或精确参数，
或串级 $\pi$ 型。判断“能否用一个 $\pi$ 型”看的是线路长度与波长之比，不是电压等级。}

\section{变压器的一相等值电路和参数}

\subsection{等值电路}

双绕组变压器：把二次绕组的电阻和漏抗折算到一次侧，与一次绕组的合并为 $R_T+jX_T$，
励磁支路 $G_T-jB_T$ \textbf{前移}到电源侧（因励磁电流很小，前移引起的误差可忽略）。
三绕组变压器：采用励磁支路前移的星形等值电路，三个绕组阻抗 $R_i+jX_i$ 均折算到一次侧。
自耦变压器的等值电路与普通变压器相同。

\subsection{双绕组变压器的参数计算}

参数由铭牌四个数据算出：短路损耗 $\Delta P_s$、短路电压百分数 $U_s\%$（由短路试验确定
$R_T$、$X_T$），空载损耗 $\Delta P_0$、空载电流百分数 $I_0\%$（由空载试验确定 $G_T$、$B_T$）。

\subsubsection{电阻 $R_T$}

短路试验时外加电压小、铁耗可忽略，故短路损耗即额定电流下原副方绕组的铜耗
$\Delta P_s=3I_N^2R_T$。改用三相额定容量 $S_N$ 与额定线电压 $U_N$ 表示，得
\keyformula{R_T=\frac{\Delta P_s U_N^{2}}{1000\,S_N^{2}}\quad(\Omega)}
其中 $\Delta P_s$ 单位 kW，$S_N$ 单位 MV$\cdot$A，$U_N$ 单位 kV（取折算侧额定线电压）。

\subsubsection{电抗 $X_T$}

短路电压百分数 $U_s\%$ 是额定电流下阻抗压降占额定电压的百分数。大容量变压器
$R_T\ll X_T$，可近似认为 $U_x\%\approx U_s\%$，于是
\keyformula{X_T=\frac{U_s\%}{100}\cdot\frac{U_N^{2}}{S_N}\quad(\Omega)}
（$S_N$ 单位 MV$\cdot$A，$U_N$ 单位 kV）。

\subsubsection{电导 $G_T$ 与电纳 $B_T$}

铁耗近似等于空载损耗，故
\begin{equation}
  G_T=\frac{\Delta P_0}{U_N^{2}}\times10^{-3}\quad(\text{S}),
\end{equation}
$\Delta P_0$ 单位 kW。空载电流的有功分量很小，无功分量近似等于空载电流，故
\begin{equation}
  B_T=\frac{I_0\%}{100}\cdot\frac{S_N}{U_N^{2}}\times10^{-3}\quad(\text{S}).
\end{equation}

\subsubsection{变比 $k_T$}

三相系统中变比指两侧绕组\textbf{空载线电压}之比，与同一铁芯柱上的匝数比有区别：
Y,y 与 D,d 接法 $k_T=U_{1N}/U_{2N}=W_1/W_2$；Y,d 接法 $k_T=U_{1N}/U_{2N}=\sqrt{3}W_1/W_2$。
运行中的\textbf{实际变比}应取两侧实际抽头的空载线电压之比（抽头有 $\pm5\%$、$\pm2\times2.5\%$ 等）。

\begin{table}[htbp]
  \centering
  \caption{双绕组变压器四个参数（铭牌数据 $\to$ 等值参数）}
  \small
  \begin{tabular}{@{}llll@{}}
    \toprule
    参数 & 算式 & 由哪个试验数据 & 数量级（110~kV、20~MV$\cdot$A 级） \\
    \midrule
    $R_T$ & $\Delta P_sU_N^{2}/(1000S_N^{2})$ & 短路损耗 $\Delta P_s$ & 几欧 \\
    $X_T$ & $U_s\%\,U_N^{2}/(100S_N)$ & 短路电压 $U_s\%$ & 几十欧 \\
    $G_T$ & $\Delta P_0\times10^{-3}/U_N^{2}$ & 空载损耗 $\Delta P_0$ & $10^{-6}$~S \\
    $B_T$ & $I_0\%\,S_N\times10^{-3}/(100U_N^{2})$ & 空载电流 $I_0\%$ & $10^{-5}$~S \\
    \bottomrule
  \end{tabular}
\end{table}

\corepoint{$R_T$、$X_T$ 依赖“折算到哪一侧”，$G_T$、$B_T$ 也随所选 $U_N$ 变化，四者必须
折算到\textbf{同一侧}才能配套使用。铭牌容量 $S_N$ 恒为三相容量，$U_N$ 恒为线电压——
这两条单位约定是本章计算题最常见的失分点。}

\subsection{三绕组变压器的参数计算}

\subsubsection{电阻}

三绕组变压器的短路试验是依次让一个绕组开路、按双绕组方式做的，测得
$\Delta P_{s(1\text{-}2)}$、$\Delta P_{s(2\text{-}3)}$、$\Delta P_{s(3\text{-}1)}$，则
\begin{align}
  \Delta P_{s(1\text{-}2)}&=\Delta P_{s1}+\Delta P_{s2},\\
  \Delta P_{s(2\text{-}3)}&=\Delta P_{s2}+\Delta P_{s3},\\
  \Delta P_{s(3\text{-}1)}&=\Delta P_{s3}+\Delta P_{s1},
\end{align}
解出各绕组短路损耗
\begin{equation}
  \Delta P_{s1}=\tfrac{1}{2}\left(\Delta P_{s(1\text{-}2)}+\Delta P_{s(3\text{-}1)}
  -\Delta P_{s(2\text{-}3)}\right),
\end{equation}
其余两式按“$1,2,3$”轮换。再按双绕组公式
\begin{equation}
  R_i=\frac{\Delta P_{si}U_N^{2}}{1000\,S_N^{2}}\quad(i=1,2,3).
\end{equation}

\textbf{容量比的处理}：三个绕组额定容量可以不等，我国常见容量比（高/中/低）为
$100/100/100$、$100/100/50$、$100/50/100$。铭牌额定容量指容量最大的绕组（高压绕组）。
短路试验受较小容量绕组额定电流限制，故工厂给的 $\Delta P'_{s}$ 须先折算：
\begin{equation}
  \Delta P_{s(1\text{-}2)}=\Delta P'_{s(1\text{-}2)}
  \left(\frac{S_N}{\min\{S_{1N},S_{2N}\}}\right)^{2},
\end{equation}
其余两组按所涉两绕组中较小的容量折算。若厂家只提供最大短路损耗 $\Delta P_{s.\max}$
（两个 $100\%$ 容量绕组通额定电流、第三绕组空载时的损耗），则额定容量绕组的电阻为
\begin{equation}
  R_{(S_N)}=\frac{\Delta P_{s.\max}U_N^{2}}{2000\,S_N^{2}},
\end{equation}
容量为 $S'_N$ 的绕组 $R'_{(S'_N)}=R_{(S_N)}\left(S_N/S'_N\right)^{2}$。

\subsubsection{电抗}

近似认为电抗上的电压降等于短路电压。给出 $U_{s(1\text{-}2)}\%$、$U_{s(2\text{-}3)}\%$、
$U_{s(3\text{-}1)}\%$ 后，各绕组短路电压为
\begin{equation}
  \begin{aligned}
    U_{s1}\%&=\tfrac{1}{2}\left(U_{s(1\text{-}2)}\%+U_{s(3\text{-}1)}\%-U_{s(2\text{-}3)}\%\right),\\
    U_{s2}\%&=\tfrac{1}{2}\left(U_{s(1\text{-}2)}\%+U_{s(2\text{-}3)}\%-U_{s(3\text{-}1)}\%\right),\\
    U_{s3}\%&=\tfrac{1}{2}\left(U_{s(2\text{-}3)}\%+U_{s(3\text{-}1)}\%-U_{s(1\text{-}2)}\%\right),
  \end{aligned}
\end{equation}
\begin{equation}
  X_i=\frac{U_{si}\%}{100}\cdot\frac{U_N^{2}}{S_N}\quad(i=1,2,3).
\end{equation}
注意：手册和厂家提供的短路电压一般已折算到变压器额定容量，可直接用上述两式计算，
\textbf{不必再折算}（这与电阻的折算形成对比，是常考的区别）。

\subsubsection{各绕组电抗的相对大小}

电抗大小取决于三个绕组在铁芯上的排列。高压绕组因绝缘要求排在最外层，中压或低压绕组
可能位于中层。\textbf{排在中层的绕组等值电抗较小，甚至可为不大负值}：
低压在中层（升压变压器常用）时有利于功率由低压向高、中压侧传送；
中压在中层（降压变压器常用）时有利于功率由高压向中压侧传送，且 $X_1$、$X_3$ 较大，
有利于限制低压侧短路电流。

导纳 $G_T-jB_T$ 与变比 $k_{12}$、$k_{13}$、$k_{23}$ 的计算与双绕组变压器相同。

\subsection{自耦变压器的参数计算}

等值电路与参数计算原理同普通变压器。三绕组自耦变压器的第三绕组（低压绕组）总接成
三角形，用以消除铁芯饱和引起的三次谐波，其容量小于变压器额定容量（高、中压绕组的
通过容量）。故计算等值电阻时须折算短路试验数据；若手册给的短路电压是未折算值，
计算等值电抗时也要先折算：
\begin{equation}
  U_{s(2\text{-}3)}\%=U'_{s(2\text{-}3)}\%\left(\frac{S_N}{S_{3N}}\right),
  \qquad
  U_{s(3\text{-}1)}\%=U'_{s(3\text{-}1)}\%\left(\frac{S_N}{S_{3N}}\right).
\end{equation}

\subsection{变压器的 $\Gamma$ 型等值电路}

用 $R_T+jX_T$（折算到原方）加理想变压器 $k_T:1$ 表示变压器时，副方求得的电流、电压都是
折算值，与副方相接的其他元件参数也须用折算值。实际计算常需副方真实值，为此在等值电路中
\textbf{增添只反映变比的理想变压器}（无损、无漏磁、无需励磁电流）。

设阻抗 $Z_T$ 在电源侧、理想变压器变比 $k$，则
\begin{equation}
  \dot U_1-Z_T\dot I_1=k\dot U_2,\qquad \dot I_1=-\dot I_2/k ,
\end{equation}
消去理想变压器得两端口的导纳方程
\keyformula{
  \begin{bmatrix}\dot I_1\\ \dot I_2\end{bmatrix}
  =\frac{1}{Z_T}
  \begin{bmatrix}1 & -k\\ -k & k^{2}\end{bmatrix}
  \begin{bmatrix}\dot U_1\\ \dot U_2\end{bmatrix}}
相应的 $\pi$ 型等值电路为：串联支路阻抗 $Z_T/k$，节点 1 对地阻抗 $Z_T/(1-k)$，
节点 2 对地阻抗 $Z_T/\left[k(k-1)\right]$。$k>1$ 时节点 1 对地阻抗为负值——
这是“非标准变比”变压器的典型特征，写节点导纳矩阵时直接用它即可（第 4 章）。

\section{标幺制}

\subsection{标幺制的概念}

标幺值 $=$ 有名值 $/$ 基准值（同单位），无量纲；百分值 $=100\times$ 标幺值。
例如取 $U_B=10.5$~kV，则 $10.5$~kV 的标幺值为 $1$。同一有名值，基准值不同则标幺值不同，
\textbf{说标幺值必须同时说明基准值}。

\subsection{基准值的选择}

四个基准值 $U_B$、$I_B$、$S_B$、$Z_B$ 应使其关系与有名值关系同形，即满足
\begin{equation}
  U_B=\sqrt{3}I_BZ_B,\qquad S_B=\sqrt{3}U_BI_B .
\end{equation}
两个方程约束四个量，故一般选定 $S_B$ 与 $U_B$（三相容量、线电压），其余两个导出：
\keyformula{I_B=\frac{S_B}{\sqrt{3}U_B},\qquad
  Z_B=\frac{U_B^{2}}{S_B},\qquad Y_B=\frac{S_B}{U_B^{2}}}
注意 $U_B$ 是线电压基准，相电压基准为 $U_B/\sqrt{3}$；$Z_B$、$Y_B$ 是\textbf{一相}的基准值。

\subsection{不同基准值的标幺值间的换算}

手册查得的电机、电器阻抗是以其自身额定容量与额定电压为基准的额定标幺阻抗。
统一到全网基准（$S_B$、$U_B$）时，先把额定标幺值还原为有名值再归算：
\keyformula{X_{*(B)}=X_{*(N)}\cdot\frac{S_B}{S_N}\cdot\left(\frac{U_N}{U_B}\right)^{2}}
此式适用于发电机与变压器。\textbf{电抗器}的额定标幺电抗以额定电压、额定电流为基准，
其换算为
\begin{equation}
  X_{R*(B)}=\frac{X_R\%}{100}\cdot\frac{U_{RN}}{\sqrt{3}I_{RN}}\cdot\frac{S_B}{U_B^{2}} .
\end{equation}
变压器的额定标幺电抗常用 $X_{T*(N)}=U_s\%/100$。

\subsection{有几级电压的网络中各元件参数标幺值的计算}

各电压等级的线路段经变压器（仅有磁耦合、无直接电气联系）相联，故可\textbf{对各段分别选
基准电压，而全系统共用同一个基准功率} $S_B$。于是
\begin{equation}
  X_{*}=\frac{X_{\text{有名值}}}{Z_B}=\frac{X_{\text{有名值}}\,S_B}{U_B^{2}} ,
\end{equation}
理想变压器的变比也要用标幺值表示：
\keyformula{k_{T*}=\frac{k_T}{k_B},\qquad k_B=\frac{U_{B(\text{I})}}{U_{B(\text{II})}}}
$k_B$ 称\textbf{基准变比}。若选择相邻段基准电压之比恰好等于该变压器的变比，
则 $k_{T*}=1$，标幺等值电路中不必串理想变压器；但此时各段基准电压不等于网络额定电压，
计算结果仍须换算成有名值才能看出电压水平。

工程简化：规定各级电压都用\textbf{平均额定电压} $U_{av}$ 作基准电压，
各变压器的额定电压和实际变比都取平均额定电压之比，于是所有 $k_{T*}=1$，
标幺等值电路成为纯阻抗网络。这就是短路实用计算（第 6 章）中“不再出现理想变压器”的根据。

\subsection{标幺制的特点}

\begin{table}[htbp]
  \centering
  \caption{标幺制的优缺点}
  \small
  \begin{tabular}{@{}p{0.06\linewidth}p{0.42\linewidth}p{0.40\linewidth}@{}}
    \toprule
     & 优点 & 说明 \\
    \midrule
    数值 & 各元件参数与运行变量都在同一数量级（$0\sim$几） & 便于判断结果是否合理、易发现计算错误 \\
    公式 & 三相与单相公式形式统一；线电压标幺值与相电压标幺值相等 &
           $\sqrt{3}$ 与 $3$ 在标幺制中消失，$S_*=U_*I_*$ \\
    归算 & 用平均额定电压作基准时变压器标幺变比为 $1$ & 多级电压网络不必逐级折算 \\
    \midrule
     & 缺点 & \\
    \midrule
    概念 & 无量纲，物理概念不如有名值清晰 & 结果须乘基准值还原为有名值 \\
    可比 & 基准不同则标幺值不可直接比较 & 必须说明基准，换基准须重算 \\
    \bottomrule
  \end{tabular}
\end{table}

换算回有名值：$U=U_*U_B$，$I=I_*I_B=I_*S_B/(\sqrt{3}U_B)$，$S=S_*S_B$，
$Z=(R_*+jX_*)U_B^{2}/S_B$。

\begin{table}[htbp]
  \centering
  \caption{常用元件标幺电抗算式（统一基准 $S_B$、$U_B$）}
  \small
  \begin{tabular}{@{}lll@{}}
    \toprule
    元件 & 原始数据 & 标幺电抗 \\
    \midrule
    发电机 & $X_{*(N)}$、$S_N$、$U_N$ & $X_{*(B)}=X_{*(N)}\dfrac{S_B}{S_N}\left(\dfrac{U_N}{U_B}\right)^{2}$ \\
    变压器 & $U_s\%$、$S_N$、$U_N$ & $X_{*(B)}=\dfrac{U_s\%}{100}\cdot\dfrac{S_B}{S_N}\left(\dfrac{U_N}{U_B}\right)^{2}$ \\
    电抗器 & $X_R\%$、$U_{RN}$、$I_{RN}$ & $X_{*(B)}=\dfrac{X_R\%}{100}\cdot\dfrac{U_{RN}}{\sqrt{3}I_{RN}}\cdot\dfrac{S_B}{U_B^{2}}$ \\
    线路 & $x_1$（$\Omega$/km）、$l$ & $X_{*(B)}=x_1l\dfrac{S_B}{U_B^{2}}$ \\
    \bottomrule
  \end{tabular}
\end{table}

\corepoint{标幺制计算的三条纪律：（1） 全系统只能有一个 $S_B$，电压基准按变比（或平均额定电压）
逐段确定；（2） 换算公式里 $U_N/U_B$ 是\textbf{平方}关系，漏平方是最常见错误；
（3） 电抗器是唯一用“额定电流”为基准的元件，不能套发电机的换算公式。}
  % 第 2 章
% ============================================================
% 第 3 章  同步发电机的基本方程
% （教材：何仰赞、温增银《电力系统分析》上册 第 3 章）
% 结构对齐 3.1—3.6；例题、习题不录
% ============================================================

\chapter{同步发电机的基本方程}

同步发电机是电力系统中最重要的元件，其运行特性对系统运行状态起决定性影响。
短路电流的暂态过程、机电暂态（稳定）分析都要用到同步电机的暂态模型。
本章从理想同步电机内部各电磁量的关系出发，建立较精确而完整的数学模型。

\section{基本前提}

\subsection{理想同步电机}

为便于分析，采用以下简化假设：
\begin{enumerate}
  \item 忽略磁路饱和、磁滞、涡流的影响，认为电机铁芯部分的导磁系数为常数（磁路线性）；
  \item 转子在结构上对纵轴（$d$ 轴）和横轴（$q$ 轴）分别对称；
  \item 定子 $a$、$b$、$c$ 三相绕组的空间位置互差 $120^\circ$ 电角度，结构完全相同，
        均在气隙中产生正弦分布的磁动势；
  \item 电机空载、转子恒速旋转时，转子绕组的磁动势在定子绕组中感应的空载电势是时间的
        正弦函数；
  \item 定子和转子的槽与通风沟不影响电感，即认为电机具有光滑表面。
\end{enumerate}
符合上述条件的电机称为\textbf{理想同步电机}。

具有阻尼绕组的凸极同步电机共有 \textbf{6 个有磁耦合关系的线圈}：定子侧静止的三个相绕组
$a$、$b$、$c$；转子侧一个励磁绕组 $f$ 和代替阻尼绕组的等值绕组 $D$、$Q$。
三个转子绕组随转子一起旋转，$f$ 与 $D$ 位于纵轴向，$Q$ 位于横轴向。
没有装设阻尼绕组的隐极电机，其实心转子的阻尼作用也可用等值阻尼绕组代表。

\begin{table}[htbp]
  \centering
  \caption{同步电机的 6 个绕组与变量}
  \small
  \begin{tabular}{@{}llll@{}}
    \toprule
    绕组 & 位置 & 端电压 & 说明 \\
    \midrule
    $a$、$b$、$c$ & 定子，静止，互差 $120^\circ$ & $u_a,u_b,u_c$ & 相绕组，正方向由中性点指向端点 \\
    $f$ & 转子纵轴 $d$ & $u_f$ & 励磁绕组，外加直流励磁电压 \\
    $D$ & 转子纵轴 $d$ & $0$ & 纵轴等值阻尼绕组 \\
    $Q$ & 转子横轴 $q$ & $0$ & 横轴等值阻尼绕组 \\
    \bottomrule
  \end{tabular}
\end{table}

\subsection{假定正向的选取}

正方向的规定直接决定方程中各项的符号，是本章最容易出错的地方。约定如下：
\begin{itemize}
  \item 转子\textbf{横轴 $q$ 轴落后于纵轴 $d$ 轴 $90^\circ$} 电角度，转子按一定方向旋转；
  \item 各绕组轴线的正方向就是该绕组\textbf{磁链的正方向}；
  \item 对本绕组产生正向磁链的电流为该绕组的\textbf{正电流}；
  \item 定子电流正方向为\textbf{由绕组中性点流向端点}，各相感应电势正方向与相电流相同，
        向外电路送出正向相电流时的机端相电压为正（发电机惯例）；
  \item 转子各绕组感应电势正方向与本绕组电流正方向相同；向励磁绕组提供正向励磁电流的
        外加励磁电压为正；两个阻尼回路的外加电压均为零。
\end{itemize}

按此正方向，定子电流的正向分量在 $d$ 轴方向起\textbf{去磁}作用，
这一条决定了后面磁链方程中定子电流项带负号。

\section{同步发电机的原始方程}

\subsection{电势方程和磁链方程}

定子和转子各回路的电势方程可写成矩阵形式
\keyformula{\begin{bmatrix}\boldsymbol u_{abc}\\ \boldsymbol u_{fDQ}\end{bmatrix}
  =-\begin{bmatrix}\boldsymbol\psi_{abc}\\ \boldsymbol\psi_{fDQ}\end{bmatrix}^{\!\!\bullet}
  -\begin{bmatrix}\boldsymbol r_s & \boldsymbol 0\\ \boldsymbol 0 & \boldsymbol r_R\end{bmatrix}
  \begin{bmatrix}\boldsymbol i_{abc}\\ \boldsymbol i_{fDQ}\end{bmatrix}}
即 $\boldsymbol u=-\dot{\boldsymbol\psi}-\boldsymbol R\boldsymbol i$，其中“$\,\dot{}\,$”表示对时间求导，
$\boldsymbol r_s=\operatorname{diag}(r,r,r)$ 为定子每相电阻，$\boldsymbol r_R=\operatorname{diag}(r_f,r_D,r_Q)$
为转子各绕组电阻。定子侧取负号是发电机惯例（电流流出）的必然结果。

磁链方程为
\begin{equation}
  \begin{bmatrix}\boldsymbol\psi_{abc}\\ \boldsymbol\psi_{fDQ}\end{bmatrix}
  =\begin{bmatrix}\boldsymbol L_{SS} & \boldsymbol L_{SR}\\
          \boldsymbol L_{RS} & \boldsymbol L_{RR}\end{bmatrix}
  \begin{bmatrix}\boldsymbol i_{abc}\\ \boldsymbol i_{fDQ}\end{bmatrix},
\end{equation}
其中 $\boldsymbol L_{SS}$ 为定子绕组自感与相互感（$3\times3$），$\boldsymbol L_{RR}$ 为转子绕组
自感与相互感，$\boldsymbol L_{SR}=\boldsymbol L_{RS}^{\mathsf T}$ 为定转子间互感。

式（电势方程）与式（磁链方程）共 \textbf{12 个方程式}，含 6 个绕组的磁链、电流、电压共
\textbf{18 个运行变量}。一般把各绕组电压作给定量，剩 12 个待求量（6 个磁链、6 个电流）；
各绕组电阻、自感与互感为已知参数。

\subsection{电感系数}

转子旋转时定、转子绕组相对位置不断变化，凸极机中部分磁通路径的磁导也随转子位置周期性
变化，故许多自感、互感系数是转子位置角 $\alpha$ 的函数。

以 $a$ 相为例：$a$ 相电流产生的正弦分布磁势 $F_a=w_a i_a$ 可分解为 $d$ 轴分量
$F_a\cos\alpha$ 与 $q$ 轴分量 $F_a\sin\alpha$，沿两轴产生气隙磁通
$\phi_{ad}=\lambda_{ad}F_a\cos\alpha$、$\phi_{aq}=\lambda_{aq}F_a\sin\alpha$，
另有漏磁通 $\phi_{\sigma a}=\lambda_{sa}F_a$。于是
\begin{equation}
  \psi_{aa}=w_a^2 i_a\left(\lambda_{sa}+\lambda_{ad}\cos^{2}\alpha+\lambda_{aq}\sin^{2}\alpha\right),
\end{equation}
\begin{equation}
  L_{aa}=\frac{\psi_{aa}}{i_a}=l_0+l_2\cos 2\alpha,
\end{equation}
\begin{equation}
  l_0=w_a^2\left[\lambda_{sa}+\tfrac{1}{2}\left(\lambda_{ad}+\lambda_{aq}\right)\right],
  \qquad
  l_2=\tfrac{1}{2}w_a^2\left(\lambda_{ad}-\lambda_{aq}\right).
\end{equation}
可见\textbf{定子自感系数是转子位置角的周期函数，周期为 $\pi$}（$\alpha=0$ 即 $d$ 轴对准该相
轴线时自感最大）。由三相对称性同理可得
\begin{align}
  L_{bb}&=l_0+l_2\cos 2(\alpha-120^\circ),&
  L_{cc}&=l_0+l_2\cos 2(\alpha+120^\circ),\\
  M_{ab}&=M_{ba}=-m_0+m_2\cos 2(\alpha+30^\circ),&\dots
\end{align}
定子与转子绕组间的互感则随 $\alpha$ 作\textbf{周期为 $2\pi$ 的正弦变化}（因为转子绕组轴线
与定子绕组轴线的夹角变化一整周才复原）。

\begin{table}[htbp]
  \centering
  \caption{原始方程中三类电感系数的变化规律}
  \small
  \begin{tabular}{@{}lll@{}}
    \toprule
    电感系数 & 变化规律 & 变化原因 \\
    \midrule
    定子自感 $L_{aa}$ 等 & $l_0+l_2\cos2\alpha$，周期 $\pi$ & 凸极转子使气隙磁导随位置变化 \\
    定子互感 $M_{ab}$ 等 & $-m_0+m_2\cos2(\alpha+30^\circ)$，周期 $\pi$ & 同上，且三相互差 $120^\circ$ \\
    定转子互感 & 随 $\alpha$ 正弦变化，周期 $2\pi$ & 两绕组轴线夹角变化一周才复原 \\
    转子绕组自感、互感 & 常数 & 转子绕组间无相对运动 \\
    \bottomrule
  \end{tabular}
\end{table}

因此，$abc$ 坐标下的基本方程是\textbf{含周期性变系数的一组微分方程}，直接求解极为困难。
解决办法是坐标变换：把定子三相变量变换到随转子一起旋转的 $d$、$q$、$0$ 坐标系，
使电感系数变为常数。

\corepoint{“变系数”不是电机本身的非线性，而是\textbf{坐标系选得不巧}。
坐标变换不改变物理本质，只把时变电感换成常电感。这是理解派克变换的唯一要点，
也是后面所有暂态分析（运算电抗、时间常数）能够成立的前提。}

\section{$dq0$ 坐标系的同步电机方程}

\subsection{坐标变换和 $dq0$ 系统}

取 $d$ 轴与 $a$ 相绕组轴线之间的电角度为 $\alpha$，令
\keyformula{
  \boldsymbol P=\frac{2}{3}
  \begin{bmatrix}
    \cos\alpha & \cos(\alpha-120^\circ) & \cos(\alpha+120^\circ)\\
    \sin\alpha & \sin(\alpha-120^\circ) & \sin(\alpha+120^\circ)\\
    \tfrac{1}{2} & \tfrac{1}{2} & \tfrac{1}{2}
  \end{bmatrix}}
则 $\boldsymbol i_{dq0}=\boldsymbol P\boldsymbol i_{abc}$，即
\begin{align}
  i_d&=\tfrac{2}{3}\left[i_a\cos\alpha+i_b\cos(\alpha-120^\circ)+i_c\cos(\alpha+120^\circ)\right],\\
  i_q&=\tfrac{2}{3}\left[i_a\sin\alpha+i_b\sin(\alpha-120^\circ)+i_c\sin(\alpha+120^\circ)\right],\\
  i_0&=\tfrac{1}{3}\left(i_a+i_b+i_c\right).
\end{align}
变换矩阵 $\boldsymbol P$ 非奇异，逆阵为
\begin{equation}
  \boldsymbol P^{-1}=
  \begin{bmatrix}
    \cos\alpha & \sin\alpha & 1\\
    \cos(\alpha-120^\circ) & \sin(\alpha-120^\circ) & 1\\
    \cos(\alpha+120^\circ) & \sin(\alpha+120^\circ) & 1
  \end{bmatrix},
  \qquad \boldsymbol i_{abc}=\boldsymbol P^{-1}\boldsymbol i_{dq0}.
\end{equation}

\textbf{零轴分量的物理意义}：由于定子三相绕组完全对称、空间互差 $120^\circ$，
三相零轴电流在气隙中的合成磁势为零，不产生与转子绕组交链的磁通，
只产生与定子绕组交链的磁通，其值与转子位置无关。$i_0$ 的表达式与对称分量法中零序电流
相似，区别是这里用瞬时值、那里用相量。

\textbf{等效绕组的物理图像}：$i_d$、$i_q$ 可看成定子两个等效绕组 $dd$、$qq$ 中的电流。
这两个绕组不像 $a$、$b$、$c$ 那样空间静止，而是\textbf{随转子一起旋转}，
其磁势对转子相对静止，所遇磁路磁阻恒定，相应的电感系数变为常数——这正是变换的目的。

\textbf{频率对应关系}（教材表 3-1 的结论）：设三相电流为对称的 $i_a=I\cos\theta$ 等，
$\theta=\theta_0+\omega' t$，$d$ 轴角速度 $\omega$，则 $i_d=I\cos(\alpha-\theta)$、
$i_q=I\sin(\alpha-\theta)$、$i_0=0$。于是
\begin{table}[htbp]
  \centering
  \caption{$abc$ 系统与 $dq0$ 系统变量的频率对应}
  \small
  \begin{tabular}{@{}lll@{}}
    \toprule
    定子电流角频率 $\omega'$ & $dq0$ 分量 & 含义 \\
    \midrule
    $\omega'=0$（直流） & 频率为 $\omega$ 的交流 & 直流 $\to$ 基频 \\
    $\omega'=\omega$（基频交流） & 常量（直流） & 对称基频交流 $\to$ 直流 \\
    $\omega'=2\omega$（倍频） & 频率为 $\omega$ 的交流 & 倍频 $\to$ 基频 \\
    \bottomrule
  \end{tabular}
\end{table}
即：三相系统的对称\textbf{基频交流}经派克变换后成为 $d$、$q$ 轴分量中的\textbf{直流}；
而三相系统的\textbf{直流与倍频}分量变换后成为 $d$、$q$ 轴中的\textbf{基频交流}。
变换可逆，故反之亦然。这一条是短路电流中“周期分量对应直流励磁磁链、
非周期分量（直流）对应基频”之类结论的数学依据。

\subsection{$dq0$ 系统的电势方程}

只对定子各量施行变换。定子电势方程 $\boldsymbol u_{abc}=-\dot{\boldsymbol\psi}_{abc}-\boldsymbol r_s\boldsymbol i_{abc}$，
左乘 $\boldsymbol P$，并利用 $\dot{\boldsymbol\psi}_{abc}=\boldsymbol P^{-1}\dot{\boldsymbol\psi}_{dq0}
+\dot{\boldsymbol P}^{-1}\boldsymbol\psi_{dq0}$，得
\begin{equation}
  \boldsymbol u_{dq0}=-\dot{\boldsymbol\psi}_{dq0}-\boldsymbol S-\boldsymbol r_s\boldsymbol i_{dq0},
  \qquad \boldsymbol S=\boldsymbol P\dot{\boldsymbol P}^{-1}\boldsymbol\psi_{dq0}.
\end{equation}
把 $\boldsymbol S$ 展开（$\dot\alpha=\omega$）即得
\keyformula{
  \begin{aligned}
    u_d&=-\dot\psi_d-\omega\psi_q-r i_d,\\
    u_q&=-\dot\psi_q+\omega\psi_d-r i_q,\\
    u_0&=-\dot\psi_0-r i_0 .
  \end{aligned}}
后两项值得注意：$\omega\psi_q$、$\omega\psi_d$ 是\textbf{旋转电势}（切割电势），
来源于坐标轴在空间旋转；$\dot\psi$ 是\textbf{变压器电势}。在稳态（$\dot\psi=0$）下只剩旋转电势，
这正是“同步电机稳态电压由励磁磁链乘转速决定”的数学表现。

\subsection{$dq0$ 系统的磁链方程和电感系数}

变换后 $d$、$q$、$0$ 三轴互相解耦，电感系数全部为常数：
\begin{align}
  \psi_d&=X_{ad}\left(i_f+i_D\right)-X_d i_d,\\
  \psi_q&=X_{aq}i_Q-X_q i_q,\\
  \psi_0&=-X_0 i_0,\\
  \psi_f&=X_{ff}i_f+X_{ad}i_D-X_{ad}i_d,\\
  \psi_D&=X_{ad}i_f+X_{DD}i_D-X_{ad}i_d,\\
  \psi_Q&=X_{aq}i_Q-X_{aq}i_q,
\end{align}
其中定子电流项前的负号来自所取正方向（正向定子电流的去磁作用）。电抗关系为
\begin{equation}
  X_d=X_{ad}+X_\sigma,\qquad X_q=X_{aq}+X_\sigma,
\end{equation}
$X_{ad}$、$X_{aq}$ 为纵轴、横轴电枢反应电抗，$X_\sigma$ 为定子漏抗；且
$X_d>X_q$（凸极机），隐极机 $X_d\approx X_q$。

\subsection{功率公式}

由三相瞬时功率 $p=u_ai_a+u_bi_b+u_ci_c$ 代入逆变换得
\keyformula{p=\frac{3}{2}\left(u_di_d+u_qi_q\right)+3u_0i_0}
系数 $3/2$、$3$ 来自变换矩阵的 $2/3$ 因子。若采用恒功率变换（系数取 $\sqrt{2/3}$），
系数变为 $1$，两种写法结果等价，但\textbf{同一题中不能混用}。
对称运行时 $i_0=u_0=0$，故
\begin{equation}
  p=\frac{3}{2}\left(u_di_d+u_qi_q\right),\qquad
  T_e=\frac{p}{\omega}\ \ (\text{标幺制下}\ \omega=1\ \text{时}\ T_e=p).
\end{equation}

\corepoint{派克变换的三个易错点：（1） 变换矩阵的 $2/3$ 系数决定了功率公式必须带 $3/2$；
（2） $q$ 轴分量用 $\sin$（本书约定），若某书用 $-\sin$，则 $i_q$ 差一符号，
后面的 $X_q$、$T_e$ 符号随之改变，\textbf{不可跨书混用公式}；
（3） $i_0$ 与对称分量法的零序电流形式相同但性质不同：这里是瞬时值，那里是相量。}

\section{同步电机的常用标幺制}

采用标幺制可使各量数值归一，并使 $d$、$q$ 轴方程的形式进一步简化。基值选取分两部分：

\subsection{定子侧基值}

选三相容量 $S_B$ 与线电压 $U_B$，则
\begin{equation}
  I_B=\frac{S_B}{\sqrt{3}U_B},\qquad Z_B=\frac{U_B^{2}}{S_B},\qquad
  \psi_B=\frac{U_B}{\omega_B},\qquad \omega_B=2\pi f_N,
\end{equation}
电阻、电抗、磁链、时间分别以 $Z_B$、$\psi_B$、$1/\omega_B$ 为基值（时间基值取 $t_B=1/\omega_B$，
这样 $\omega$ 的标幺值为 $1$）。

\subsection{转子侧基值与“互感可逆”}

转子各绕组（$f$、$D$、$Q$）的基值电流、基值电压不能任意选，须满足
\begin{quote}
  使定子绕组与转子绕组之间的\textbf{互感系数标幺值相等}（即互感可逆），
  从而使 $X_{ad}$、$X_{aq}$ 在定子侧与转子侧只有一个标幺值。
\end{quote}
具体做法是先选定转子基值电流 $I_{fB}$（常取空载额定电压所需的励磁电流作为参考），
再由 $X_{ad*}=1$ 之类的条件反定电压基值 $U_{fB}=Z_{fB}I_{fB}$。
这样处理后，$d$、$q$ 轴磁链方程中标幺化的互感系数都等于 $X_{ad*}$、$X_{aq*}$，
方程形式与有名值方程完全一致，只是各量加了星号。

\textbf{实用效果}：标幺制下 $X_d$、$X_q$、$X_d'$、$X_d''$、$X_q''$ 都是同一数量级的纯数
（如 $X_d=1.0\sim2.5$），时间常数以秒计，转矩与功率标幺值相等（$\omega_*=1$）。

\section{基本方程的拉氏运算形式}

\subsection{有阻尼绕组同步电机的运算方程}

假定转子转速恒定为额定转速，对派克方程作拉氏变换（$p=\mathrm{d}/\mathrm{d}t\to s$），
把微分方程化为象函数的代数方程，然后求解。可得
\begin{equation}
  \psi_d(s)=X_d(s)\,i_d(s)+G_f(s)u_f(s)+\cdots
\end{equation}
其中
\keyformula{X_d(p)=X_\sigma+\frac{1}{\dfrac{1}{X_{ad}}+\dfrac{1}{X_{f\sigma}+r_f/p}
  +\dfrac{1}{X_{D\sigma}+r_D/p}},\qquad
  X_q(p)=X_\sigma+\frac{1}{\dfrac{1}{X_{aq}}+\dfrac{1}{X_{Q\sigma}+r_Q/p}}}
$X_d(p)$、$X_q(p)$ 称为同步电机的\textbf{纵轴、横轴运算电抗}；此外还引出传递函数
$G_f(p)$、$G_D(p)$、$G_Q(p)$，反映励磁电压、阻尼绕组磁链对定子量的作用。

\subsection{运算电抗的等值电路}

运算电抗有直观的等值电路解释：
\begin{itemize}
  \item \textbf{纵轴}：$X_\sigma$ 串联“$X_{ad}$、$X_{f\sigma}+r_f/p$、$X_{D\sigma}+r_D/p$
        \textbf{三支路并联}”；
  \item \textbf{横轴}：$X_\sigma$ 串联“$X_{aq}$ 与 $X_{Q\sigma}+r_Q/p$ 两支路并联”。
\end{itemize}
物理意义很清楚：$p\to\infty$（突变瞬间）时 $r/p\to0$，各阻尼、励磁支路相当于短路，
并联支路被短接，得
\begin{equation}
  X_d(0)=X_\sigma+\left(\frac{1}{X_{ad}}+\frac{1}{X_{f\sigma}}+\frac{1}{X_{D\sigma}}\right)^{-1}
  =X_d'',
\end{equation}
即\textbf{次暂态电抗}；$p\to0$（稳态）时 $r/p\to\infty$，各转子支路相当于开路，
得 $X_d(\infty)=X_{ad}+X_\sigma=X_d$，即\textbf{同步电抗}。
只有一个转子绕组时对应 $X_d'$（暂态电抗）。这一“极限即参数”的对应关系，
是第 5 章短路电流中各阶段电流取值的根本依据。

\subsection{无阻尼绕组同步电机的运算方程}

无阻尼绕组（或隐极机不带阻尼绕组）时，$d$ 轴只剩励磁支路、$q$ 轴无转子绕组：
\begin{equation}
  X_d(p)=X_\sigma+\frac{1}{\dfrac{1}{X_{ad}}+\dfrac{1}{X_{f\sigma}+r_f/p}},\qquad
  X_q(p)=X_q=X_{aq}+X_\sigma .
\end{equation}
此时 $q$ 轴无暂态过程（$X_q$ 恒定），只有 $d$ 轴存在暂态电抗 $X_d'$ 与时间常数 $T_d'$。
第 5 章对“无阻尼绕组电机突然三相短路”的分析就用这一组方程。

\section{同步电机的对称稳态运行}

\subsection{基本方程的实用化}

稳态对称运行时：$\dot\psi_d=\dot\psi_q=0$，$i_0=0$，转子转速为同步速（$\omega_*=1$），
定子电阻压降远小于电抗压降，通常可忽略（$r=0$）。$dq0$ 电势方程退化为
\keyformula{u_d=-\psi_q=X_q i_q,\qquad u_q=\psi_d=X_{ad}i_f-X_d i_d=E_q-X_d i_d}
其中
\begin{equation}
  E_q=X_{ad}i_f=\psi_{fd}
\end{equation}
为励磁电流在定子绕组中产生的互感磁链（有用磁链）及其感应电势，即通常所指的\textbf{空载电势}。

\subsection{稳态运行的电势方程式、相量图和等值电路}

\textbf{隐极机}（$X_d=X_q$）：由 $u_d=X_qi_q$、$u_q=E_q-X_di_d$ 合成相量式
\begin{equation}
  \dot E_q=\dot U+jX_d\dot I .
\end{equation}
等值电路为：电势源 $\dot E_q$ 串 $jX_d$ 后接机端。$\dot E_q$ 与 $\dot U$ 的相位差即\textbf{功角} $\delta$。

\textbf{凸极机}（$X_d\ne X_q$）：因 $X_d\ne X_q$，$d$、$q$ 轴压降不等，须引入辅助电势
\begin{equation}
  \dot E_Q=\dot U+jX_q\dot I ,
\end{equation}
$E_Q$ 落在 $q$ 轴上。再由
\begin{equation}
  E_q=E_Q-(X_d-X_q)i_d
\end{equation}
求空载电势。故凸极机的相量图须先定 $q$ 轴方向（$\dot E_Q$ 方向即 $q$ 轴），
再分解 $d$、$q$ 分量。

\begin{table}[htbp]
  \centering
  \caption{稳态方程与等值电路的对照}
  \small
  \begin{tabular}{@{}p{0.12\linewidth}p{0.34\linewidth}p{0.40\linewidth}@{}}
    \toprule
    机型 & 方程 & 说明 \\
    \midrule
    隐极机 & $\dot E_q=\dot U+jX_d\dot I$ & 一个电势、一个电抗，可直接求 $\delta$ \\
    凸极机 & $\dot E_Q=\dot U+jX_q\dot I$，\newline $E_q=E_Q-(X_d-X_q)i_d$ &
             先由 $\dot E_Q$ 定 $q$ 轴，再分解；$X_d\ne X_q$ 造成附加项 \\
    标幺形式 & $\omega_*=1$ 时 $u_d=X_qi_q$，$u_q=E_q-X_di_d$ & 转矩与功率数值相等 \\
    \bottomrule
  \end{tabular}
\end{table}

\corepoint{稳态分析的三步：（1） 由 $\dot U$、$\dot I$ 及 $X_q$ 定 $q$ 轴（凸极机）或直接用 $X_d$
（隐极机）；（2） 分解 $i_d$、$i_q$；（3） 由 $E_q=u_q+X_di_d$ 反求空载电势，$\delta$ 为 $\dot E_q$
超前 $\dot U$ 的角度。功角 $\delta$ 是第 15、16 章全部稳定分析的基准量，
其物理含义（转子磁极轴线领先定子合成磁极轴线的电角度）必须在本章就建立起来。}
  % 第 3 章
% ============================================================
% 第 4 章  电力网络的数学模型
% （教材：何仰赞、温增银《电力系统分析》上册 第 4 章）
% 结构对齐 4.1—4.4；例题、习题不录
% ============================================================

\chapter{电力网络的数学模型}

第 2 章给了元件参数与等值电路，第 3 章给了发电机方程。本章把二者按网络的拓扑关系组装起来，
得到\textbf{网络方程}——用节点电压（或支路电流）描述整个系统状态的代数方程组。
这一章是短路计算（第 6、8 章）、潮流计算（第 11 章）和稳定计算（第 17、18 章）共同的
“装配车间”。

\section{节点导纳矩阵}

\subsection{节点方程}

设网络有 $n$ 个节点（不含参考节点），各节点注入电流为 $\dot I_i$（由外部注入网络的电流，
发电机取正、负荷取负），节点电压为 $\dot U_i$。由基尔霍夫电流定律，对每个节点写出
“注入电流 $=$ 流出该节点的各支路电流之和”，整理即得节点方程
\keyformula{
  \begin{bmatrix}\dot I_1\\ \dot I_2\\ \vdots\\ \dot I_n\end{bmatrix}
  =\begin{bmatrix}
    Y_{11} & Y_{12} & \cdots & Y_{1n}\\
    Y_{21} & Y_{22} & \cdots & Y_{2n}\\
    \vdots & \vdots & \ddots & \vdots\\
    Y_{n1} & Y_{n2} & \cdots & Y_{nn}
  \end{bmatrix}
  \begin{bmatrix}\dot U_1\\ \dot U_2\\ \vdots\\ \dot U_n\end{bmatrix}}
简记为 $\boldsymbol I=\boldsymbol Y\boldsymbol U$，$\boldsymbol Y$ 称\textbf{节点导纳矩阵}。
它是 $n$ 阶复数方阵，对称（无互感、无受控源时），且高度稀疏。

\textbf{注意参考节点的处理}：大地（或零电位点）取为参考节点，其电压为零，
不出现在方程中；与参考节点相连的对地支路（如线路并联电纳、变压器励磁支路）
只影响相应节点的自导纳。

\subsection{节点导纳矩阵元素的物理意义}

由 $\boldsymbol U=\boldsymbol Y^{-1}\boldsymbol I$ 或直接由节点方程可读出：

\paragraph{对角元素 $Y_{ii}$（自导纳）}
在节点 $i$ 单独注入单位电流、其余节点注入电流为零（即其余节点短路接地）时，
节点 $i$ 的电压值。其值为
\begin{equation}
  Y_{ii}=\sum_{j\in i}y_{ij}+y_{i0},
\end{equation}
即\textbf{与节点 $i$ 相连的所有支路导纳之和}（含节点 $i$ 的对地支路 $y_{i0}$）。

\paragraph{非对角元素 $Y_{ij}$（互导纳）}
在上述条件下节点 $j$ 的电压值，其值为
\begin{equation}
  Y_{ij}=-y_{ij}\qquad(i\ne j),
\end{equation}
即\textbf{节点 $i$、$j$ 之间支路导纳的负值}。若 $i$、$j$ 之间无直接支路，则 $Y_{ij}=0$。

\begin{table}[htbp]
  \centering
  \caption{节点导纳矩阵元素的构成规则}
  \small
  \begin{tabular}{@{}p{0.20\linewidth}p{0.30\linewidth}p{0.34\linewidth}@{}}
    \toprule
    元素 & 构成 & 说明 \\
    \midrule
    $Y_{ii}$ & 与 $i$ 相连全部支路导纳之和（含对地） & 恒不为零，数值最大 \\
    $Y_{ij}$ & $-y_{ij}$ & 只在 $i$、$j$ 有直接支路时非零 \\
    对地支路 $y_{i0}$ & 只加到 $Y_{ii}$ & 不出现在非对角元素中 \\
    无源元件 & 纯虚数（$j\omega$ 或 $1/(j\omega)$） & 实际含电阻时略呈负实部 \\
    \bottomrule
  \end{tabular}
\end{table}

\corepoint{$Y_{ij}$ 带负号，$Y_{ii}$ 是“和”不是“差”，这两个约定几乎所有教材一致；
但“自导纳含不含对地支路”各书写法有别——本书\textbf{含}。做题时先确认这一条，
否则对角元素会差一个并联电纳。}

\subsection{节点导纳矩阵的修改}

实际计算中常需在原网络基础上增删元件而不重新组装矩阵，此时只需修改少数元素：

\begin{table}[htbp]
  \centering
  \caption{网络结构变化时节点导纳矩阵的修改}
  \small
  \begin{tabular}{@{}p{0.22\linewidth}p{0.62\linewidth}@{}}
    \toprule
    操作 & 矩阵元素的修改 \\
    \midrule
    在 $i$、$j$ 间追加导纳为 $y$ 的支路 &
      $Y_{ii}\mathrel{+}=y$，$Y_{jj}\mathrel{+}=y$，$Y_{ij}=Y_{ji}\mathrel{-}=y$ \\
    切除 $i$、$j$ 间导纳为 $y$ 的支路 &
      $Y_{ii}\mathrel{-}=y$，$Y_{jj}\mathrel{-}=y$，$Y_{ij}=Y_{ji}\mathrel{+}=y$ \\
    在节点 $i$ 增加对地支路 $y$ & 仅 $Y_{ii}\mathrel{+}=y$ \\
    切除节点 $i$ 的对地支路 $y$ & 仅 $Y_{ii}\mathrel{-}=y$ \\
    节点 $i$ 的电压变比改变（变压器非标准变比 $k_*$） &
      按 $\dfrac{1}{Z_T}\begin{bmatrix}1&-k_*\\-k_*&k_*^{2}\end{bmatrix}$
      重算与 $i$、$j$ 有关的四个元素（见第 2 章 $\Gamma$ 型等值电路） \\
    两节点合并（如开关合上） & 把两行两列相加，删去一行一列 \\
    \bottomrule
  \end{tabular}
\end{table}

$k_*$ 改变引起的变化最能说明问题：非标准变比变压器的四个导纳元素
$Y_{ii}=y_T$、$Y_{ij}=Y_{ji}=-k_*y_T$、$Y_{jj}=k_*^{2}y_T$ 中，
只有 $k_*=1$ 时才退化为普通串联支路的“$+y$、$-y$、$+y$”形式。
因此\textbf{有载调压变压器的分接头一动，节点导纳矩阵的四个元素同时变}——
这也是潮流计算中“变压器分接头是控制变量”的矩阵解释。

\subsection{支路间存在互感时的节点导纳矩阵}

当网络中存在互感支路（如平行架设的双回线、同杆并架线路）时，
“支路导纳相加”的规则失效，因为支路电流不再只由本支路电压决定。此时用
\begin{equation}
  \boldsymbol Y=\boldsymbol A\,\boldsymbol Z_b^{-1}\boldsymbol A^{\mathsf T},
\end{equation}
其中 $\boldsymbol A$ 为\textbf{节点-支路关联矩阵}（$n\times b$，
支路 $k$ 由 $i$ 指向 $j$ 时 $A_{ik}=1$、$A_{jk}=-1$，其余为零），
$\boldsymbol Z_b$ 为\textbf{支路阻抗矩阵}（$b\times b$，对角为各支路自阻抗，
非对角为互感支路间的互阻抗）。
\keyformula{\boldsymbol I=\boldsymbol A\boldsymbol Z_b^{-1}\boldsymbol A^{\mathsf T}\boldsymbol U
  \;\Longrightarrow\;\boldsymbol Y=\boldsymbol A\boldsymbol Z_b^{-1}\boldsymbol A^{\mathsf T}}
这一式子把“拓扑”（$\boldsymbol A$）与“参数”（$\boldsymbol Z_b$）分开，
有互感时先对 $\boldsymbol Z_b$ 求逆再投影到节点空间，结果自动满足耦合关系；
无互感时 $\boldsymbol Z_b$ 对角，$\boldsymbol Y$ 退化为上面“相加”的规则。
第 7 章处理零序网络中的平行双回线时，用的就是这个式子。

\section{网络方程的解法}

\subsection{用高斯消去法求解网络方程}

$\boldsymbol Y\boldsymbol U=\boldsymbol I$ 是复数线性方程组。直接求解用高斯消去法
（按列消去），把 $\boldsymbol Y$ 化为上三角后回代。消去过程中矩阵元素按
\begin{equation}
  Y_{ij}^{(k)}=Y_{ij}^{(k-1)}-\frac{Y_{ik}^{(k-1)}Y_{kj}^{(k-1)}}{Y_{kk}^{(k-1)}},
  \qquad i,j>k,
\end{equation}
更新，同时右端项按
\begin{equation}
  I_i^{(k)}=I_i^{(k-1)}-\frac{Y_{ik}^{(k-1)}}{Y_{kk}^{(k-1)}}I_k^{(k-1)}
\end{equation}
更新。$\boldsymbol Y$ 对称时只需存上三角，运算量约 $n^3/3$。

由于 $\boldsymbol Y$ 高度稀疏，实际计算采用\textbf{稀疏技术}（只存非零元素、只对非零元素运算）
和\textbf{节点编号优化}（见 4.4 节），使运算量从 $O(n^3)$ 降至近似 $O(n^{1.5})$ 量级。

\subsection{用高斯消去法简化网络}

消去法还有一层\textbf{网络含义}：消去某个节点，等价于按“星-网变换”把该节点从网络中
抹去，得到保留节点的等值网络。

对节点 $k$ 消去后，保留节点之间的等值导纳为
\keyformula{Y_{ij}'=Y_{ij}-\frac{Y_{ik}Y_{kj}}{Y_{kk}}\qquad(i,j\ne k)}
这正是星形（以 $k$ 为中心）化为网形的公式；若 $k$ 只有两条支路，上式就退化为
两阻抗串联。若 $k$ 是“浮游节点”（只与两个节点相连、无注入），消去它相当于把两条
支路合并。

这一操作在短路计算中极有用：把不关心的中间节点逐个消去，最后得到“电源点到短路点”的
等值阻抗，即第 6 章的\textbf{转移阻抗}。

\corepoint{消去法的两种读法：从代数看是高斯消去，从网络看是星-网变换（$Y$-$\Delta$ 变换）。
两者是同一件事。理解这一点，第 6 章“用节点阻抗矩阵算短路电流”和第 4.3 节
“支路追加法”才不是两套孤立的算法。}

\section{节点阻抗矩阵}

\subsection{节点阻抗矩阵元素的物理意义}

节点阻抗矩阵定义为节点导纳矩阵的逆
\begin{equation}
  \boldsymbol Z=\boldsymbol Y^{-1},\qquad \boldsymbol U=\boldsymbol Z\boldsymbol I .
\end{equation}
由 $\boldsymbol U=\boldsymbol Z\boldsymbol I$，令仅在节点 $i$ 注入单位电流（其余节点注入为零），则
\begin{itemize}
  \item $Z_{ii}$：节点 $i$ 的电压，称节点 $i$ 的\textbf{自阻抗}（输入阻抗）；
        它就是\textbf{从节点 $i$ 看进网络、其余节点开路时的等值阻抗}。
  \item $Z_{ji}$（$j\ne i$）：节点 $j$ 的电压，称节点 $j$ 与 $i$ 之间的\textbf{转移阻抗}。
\end{itemize}
\textbf{关键区别}：$Y_{ij}$ 只与 $i$、$j$ 间的直接支路有关（局部量），
而 $Z_{ij}$ 与\textbf{整个网络}有关（全局量）——任一支路变化都会改变所有 $Z$ 元素。
因此 $\boldsymbol Z$ 是满阵，不宜直接存储，只在需要少数元素时用。

\begin{table}[htbp]
  \centering
  \caption{节点导纳矩阵与节点阻抗矩阵的对比}
  \small
  \begin{tabular}{@{}p{0.14\linewidth}p{0.30\linewidth}p{0.34\linewidth}@{}}
    \toprule
    比较项 & $\boldsymbol Y$ & $\boldsymbol Z=\boldsymbol Y^{-1}$ \\
    \midrule
    稀疏性 & 高度稀疏 & 满阵 \\
    元素含义 & 局部（支路导纳） & 全局（整网的等值阻抗） \\
    修改一条支路 & 只改 4 个元素 & 全部元素都变 \\
    主要用途 & 潮流、稳定计算（迭代求解） & 短路计算（$I_f=E/Z_{ff}$，只取少数元素） \\
    求法 & 直接组装 & 支路追加法、或对 $\boldsymbol Y$ 求逆（部分列） \\
    \bottomrule
  \end{tabular}
\end{table}

\subsection{用支路追加法形成节点阻抗矩阵}

支路追加法从“只有一个接地支路的空网络”开始，逐条追加支路，每追加一条就把 $\boldsymbol Z$
扩充或修改。它避免了求逆，且物理意义直观。分三种情形：

\paragraph{（1）从已有节点 $i$ 引出新节点 $j$（追加树支）}
新节点 $j$ 经阻抗 $z$ 接到 $i$，则
\begin{equation}
  Z_{jj}=Z_{ii}+z,\qquad Z_{jk}=Z_{kj}=Z_{ik}\ (k\le i).
\end{equation}
即新行、新列“抄”第 $i$ 行、第 $i$ 列，对角加 $z$。

\paragraph{（2）在原有节点 $i$、$j$ 间追加连支（形成回路）}
若追加的支路阻抗为 $z$、且 $i$、$j$ 间已存在通路，则回路阻抗为
$z_{\Sigma}=Z_{ii}+Z_{jj}-2Z_{ij}+z$，全部元素按
\begin{equation}
  Z_{kl}'=Z_{kl}-\frac{\left(Z_{ki}-Z_{kj}\right)\left(Z_{il}-Z_{jl}\right)}{z_{\Sigma}}
\end{equation}
修正。这是“追加一条连支 $=$ 在回路中并入一个阻抗”的结果。

\paragraph{（3）追加接地支路（从节点 $i$ 到地）}
\begin{equation}
  Z_{kl}'=Z_{kl}-\frac{Z_{ki}Z_{il}}{Z_{ii}+z}.
\end{equation}

三种情形中，只有第（1）种会增加矩阵阶数。追加顺序不同，中间结果不同，但最终 $\boldsymbol Z$ 相同。

\subsection{用线性方程直接解法对导纳矩阵求逆}

实际只需 $\boldsymbol Z$ 的少数元素（如短路点的 $Z_{ff}$）时，不必求全逆：
用因子表法对 $\boldsymbol Y$ 作三角分解 $\boldsymbol Y=\boldsymbol L\boldsymbol D\boldsymbol L^{\mathsf T}$
（对称阵），则求 $\boldsymbol Z$ 的第 $j$ 列等价于解
$\boldsymbol L\boldsymbol D\boldsymbol L^{\mathsf T}\boldsymbol z_j=\boldsymbol e_j$
（$\boldsymbol e_j$ 为单位向量）。由于 $\boldsymbol e_j$ 稀疏，前代、回代都很快，
算一列的成本约为一次完整求解的成本。

\corepoint{短路计算只需要 $Z_{ff}$ 和少量 $Z_{if}$，所以\textbf{不要}为了“求逆”去算满阵——
这是计算量相差几个数量级的区别。教材强调“直接解法求逆的列”，正是这个意思。}

\section{节点编号顺序的优化}

高斯消去时，消去一个节点会在其“邻居”之间产生新的非零元素（注入元，fill-in）。
注入元的多少强烈依赖\textbf{消去顺序}，因此节点编号顺序直接影响计算量和存储量。

\begin{itemize}
  \item 直观准则：\textbf{先消去出线数（连接支路数）最少的节点}。消去一个度为 $d$ 的节点，
        最多新增 $d(d-1)/2$ 条等值支路，度越小注入元越少。
  \item 工程做法：按“静态最小度”排序，或按消去树（elimination tree）的深度分层编号；
        实际程序中常用近似最小度（AMD）算法。
  \item 另一种思路：把网络\textbf{按地理分区}编号——区域内节点连续编号，区域间只留少量联络线，
        这样消去区内节点时注入元被限制在区内。
\end{itemize}

编号优化只改变计算过程的稀疏性，\textbf{不改变}任何电气结果（$\boldsymbol Y$、$\boldsymbol Z$
的元素值按节点对应关系重排而已）。

\begin{table}[htbp]
  \centering
  \caption{本章四个工具的分工}
  \small
  \begin{tabular}{@{}p{0.18\linewidth}p{0.30\linewidth}p{0.40\linewidth}@{}}
    \toprule
    工具 & 解决什么问题 & 用在哪一章 \\
    \midrule
    $\boldsymbol Y$ 矩阵 & 组装网络、迭代求解 & 潮流（11）、稳定（17、18） \\
    $\boldsymbol Z$ 矩阵 & 快速给出任一点短路的等值阻抗 & 短路（6、8） \\
    高斯消去 & 解方程、简化网络（星-网变换） & 全网计算 \\
    编号优化 & 降低注入元、提高速度 & 大规模网络计算 \\
    \bottomrule
  \end{tabular}
\end{table}

\corepoint{本章的考试落点通常是三件事：（1） 给网络图写 $\boldsymbol Y$（注意对地支路只进对角、
互导纳取负）；（2） 给 $\boldsymbol Y$ 与注入电流求节点电压（或反过来由 $\boldsymbol U$ 求注入电流）；
（3） 说明 $Z_{ii}$、$Z_{ij}$ 的物理意义并会用“单点注入单位电流”解释。
把这三件事的物理图像记牢，比背公式有用。}
  % 第 4 章
% ============================================================
% 第 5 章  电力系统三相短路的暂态过程
% （教材：何仰赞、温增银《电力系统分析》上册 第 5 章）
% 结构对齐 5.1—5.6；例题、习题不录
% 本文件用 \chapter 承接教材第 5 章
% ============================================================

\chapter{电力系统三相短路的暂态过程}

短路是电力系统最严重的故障。三相短路是对称短路，可以用一相电路分析，是全部短路计算的
基础；而各种不对称短路（第 7、8 章）最终也都要回到“正序等效网络 $+$ 附加电抗”的形式，
其正序部分与三相短路同构。因此本章既是暂态分析的第一课，也是短路计算的总入口。

\section{短路的一般概念}

\subsection{短路的原因、类型及后果}

\textbf{原因}：绝缘损坏（老化、污秽、受潮、机械损伤）、误操作（带负荷拉刀闸、带地线合闸）、
鸟兽跨接、线路断线倒杆、雷击过电压等。

\textbf{类型}（按相数与接地情况）：

\begin{table}[htbp]
  \centering
  \caption{短路的类型与特点}
  \small
  \begin{tabular}{@{}llll@{}}
    \toprule
    类型 & 符号 & 是否对称 & 出现概率与说明 \\
    \midrule
    三相短路 & $f^{(3)}$ & 对称 & 概率最低（约 $5\%$），但后果最严重 \\
    两相短路 & $f^{(2)}$ & 不对称 & 约 $10\%\sim15\%$，无零序分量 \\
    单相接地短路 & $f^{(1)}$ & 不对称 & 概率最高（约 $65\%\sim70\%$），有零序分量 \\
    两相接地短路 & $f^{(1,1)}$ & 不对称 & 约 $10\%\sim20\%$，有零序分量 \\
    \bottomrule
  \end{tabular}
\end{table}

\textbf{后果}：
\begin{enumerate}
  \item \textbf{电动力效应}：短路电流可达额定电流的十几到几十倍，导体间电动力与电流平方成正比，
        可使设备变形、损坏；
  \item \textbf{热效应}：巨大电流在导体和设备中产生大量热量，短时即可使绝缘烧毁；
  \item \textbf{电压骤降}：短路点附近电压大幅下降，影响用户用电，严重时导致电动机失速、
        系统失去同步；
  \item \textbf{破坏稳定}：短路切除不及时会引发系统振荡甚至崩溃；
  \item \textbf{不对称短路的附加危害}：零序电流产生零序磁通，对邻近通信线路产生干扰。
\end{enumerate}

\subsection{短路计算的目的}

\begin{itemize}
  \item 选择电气设备：按短路电流校验\textbf{动稳定}（电动力）与\textbf{热稳定}（发热），
        并据此选定断路器开断能力；
  \item 继电保护的设计与整定：确定保护范围、整定值与灵敏度；
  \item 确定运行方式：比较不同接线与开机方式下的短路电流水平，选择限制短路电流的措施
        （如变压器中性点接地方式、加装电抗器、分段运行）；
  \item 为稳定计算提供初值：短路期间的电磁功率、切除时刻的功角都依赖短路电流。
\end{itemize}

\corepoint{短路计算的两条主线贯穿全书：\textbf{手算}路线（等值电路 $+$ 化简 $+$ 查计算曲线，
第 6 章）与\textbf{计算机}路线（节点阻抗矩阵 $+$ 序网方程，第 4、8 章）。
两条路线的物理模型完全相同，区别只在“化简到多小”和“谁来解方程”。}

\section{恒定电势源电路的三相短路}

\subsection{短路的暂态过程}

设无限大功率电源（内阻抗为零、端电压恒定）经 $R$、$L$ 供电，$t=0$ 时在出口发生三相短路。
短路前稳态电流
\begin{equation}
  i=\frac{E_m}{\sqrt{R^{2}+(\omega L)^{2}}}\sin(\omega t+\alpha-\varphi)
  =I_m\sin(\omega t+\alpha-\varphi),
\end{equation}
短路后电路成为无源 $R$-$L$ 回路，回路方程为
\begin{equation}
  L\frac{\mathrm{d}i}{\mathrm{d}t}+Ri=E_m\sin(\omega t+\alpha).
\end{equation}
其解为\textbf{周期分量（强制分量）}与\textbf{非周期分量（自由分量）}之和
\keyformula{i=I_{pm}\sin(\omega t+\alpha-\varphi')+C\,e^{-t/T_a},
  \qquad I_{pm}=\frac{E_m}{\sqrt{R^{2}+(\omega L)^{2}}},\quad T_a=\frac{L}{R}}
积分常数由\textbf{短路瞬间电流不能突变}（电感中电流连续）确定：
\begin{equation}
  C=I_m\sin(\alpha-\varphi)-I_{pm}\sin(\alpha-\varphi').
\end{equation}
即自由分量的初值等于“短路前瞬间电流”与“短路后周期分量在 $t=0$ 的值”之差。

\textbf{结论（写清条件）}：
\begin{enumerate}
  \item 短路电流最大可能值出现在\textbf{短路前空载}（$I_m=0$）且
        $\alpha-\varphi'=\pm90^\circ$（即电压过零瞬间短路）时，
        此时自由分量初值等于周期分量幅值，半周后出现最大瞬时值；
  \item 若短路瞬间恰好 $C=0$，则不出现非周期分量，电流立即进入稳态——这是“最有利”的合闸角；
  \item 非周期分量按 $T_a=L/R$ 衰减，$R$ 越大衰减越快；
  \item 周期分量幅值由回路阻抗决定，与合闸角无关。
\end{enumerate}

\subsection{短路冲击电流}

短路电流最大瞬时值称\textbf{短路冲击电流} $i_{sh}$，出现在短路后约半个周期（$t=0.01$~s，
即 $50$~Hz 的半周）：
\keyformula{i_{sh}=I_{pm}+I_{pm}e^{-0.01/T_a}=\left(1+e^{-0.01/T_a}\right)I_{pm}
  =K_{sh}I_{pm}=K_{sh}\sqrt{2}I''}
其中 $I''$ 为短路电流周期分量的\textbf{起始有效值}（次暂态电流），
$K_{sh}=1+e^{-0.01/T_a}$ 称\textbf{冲击系数}，$1<K_{sh}<2$。
工程上：发电机出口短路取 $K_{sh}=1.90$（$i_{sh}=2.69I''$）；
一般高压电网取 $K_{sh}=1.80$（$i_{sh}=2.55I''$）。

\subsection{短路电流的最大有效值}

以 $i_{sh}$ 出现的那个周期为考察对象，该周期内电流有效值称\textbf{短路电流最大有效值}
（冲击电流有效值）：
\keyformula{I_{sh}=\sqrt{I''^{2}+i_{sh}^{2}}\approx I''\sqrt{1+2\left(K_{sh}-1\right)^{2}}}
$K_{sh}=1.8$ 时 $I_{sh}=1.52I''$；$K_{sh}=1.9$ 时 $I_{sh}=1.62I''$。
$I_{sh}$ 用于校验电气设备的\textbf{动稳定}（与制造厂给出的动稳定电流峰值比较）。

\subsection{短路功率}

短路功率（短路容量）定义为短路点处三相视在功率
\keyformula{S_k=\sqrt{3}U_{av}I''}
$U_{av}$ 为短路点所在电压级的平均额定电压（kV），$I''$ 单位 kA，$S_k$ 单位 MV$\cdot$A。
其物理含义是“短路点看进去的等值阻抗所对应的容量”，$S_k=U_{av}^{2}/X_{\Sigma}$
（标幺制下 $S_{k*}=1/X_{\Sigma*}$）。断路器开断能力即按 $S_k$ 选择。

\begin{table}[htbp]
  \centering
  \caption{短路电流的三个特征量及其用途}
  \small
  \begin{tabular}{@{}p{0.16\linewidth}p{0.28\linewidth}p{0.42\linewidth}@{}}
    \toprule
    量 & 定义 & 用途 \\
    \midrule
    $I''$ & 周期分量起始有效值（次暂态电流） & 短路容量、继电保护整定、热稳定校验 \\
    $i_{sh}$ & 最大瞬时值，$t\approx0.01$~s & 动稳定校验（峰值） \\
    $I_{sh}$ & 含非周期分量的最大有效值 & 动稳定校验（有效值） \\
    \bottomrule
  \end{tabular}
\end{table}

\corepoint{$i_{sh}$ 与 $I_{sh}$ 都建立在“最不利合闸角 $+$ 空载”这一\textbf{条件}上，
是偏保守的校验值；而 $I''$ 与合闸角无关，只由电网阻抗决定。
三者中只有 $I''$ 是“电路算出来的量”，另两个是“由 $I''$ 派生的校验量”。}

\section{同步电机突然三相短路的物理分析}

\subsection{突然短路暂态过程的特点}

与恒定电势源电路相比，同步电机突然短路有三个新特点：
\begin{enumerate}
  \item \textbf{电源内阻抗不能忽略且随时间变化}：短路瞬间电枢反应磁通被转子绕组“挡住”，
        只能走漏磁路径，等值电抗小（$X_d''$）；随时间推移磁通逐渐穿入转子，
        电抗增大到 $X_d'$ 直至 $X_d$；
  \item \textbf{励磁磁链不能突变}：转子励磁绕组和阻尼绕组都是闭合回路，电阻很小，
        短路瞬间磁链近似守恒，从而产生一系列自由电流（定子非周期分量、转子自由分量）；
  \item \textbf{定子电流中含有倍频分量}：凸极机 $X_d\ne X_q$，转子纵、横轴磁阻不同，
        转子磁极旋转时定子绕组看到的磁阻周期性变化，产生 $2\omega$ 分量。
\end{enumerate}

\subsection{超导体闭合回路磁链守恒原则}

若闭合回路的电阻为零，则由 $\oint\boldsymbol E\cdot\mathrm{d}\boldsymbol l=-\mathrm{d}\psi/\mathrm{d}t$
与 $\oint\boldsymbol E\cdot\mathrm{d}\boldsymbol l=Ri=0$ 得
\keyformula{\psi=\text{常数}\qquad(R=0)}
即\textbf{超导体闭合回路的磁链不能突变}。发电机转子绕组电阻很小，在短路初期的几个周期内
可近似按此处理。这是整个“突然短路物理分析”的唯一出发点：
\begin{itemize}
  \item 励磁绕组磁链 $\psi_f$ 短路瞬间不变 $\Rightarrow$ 产生励磁绕组自由直流分量，
        以抵消定子电流对它的去磁作用；
  \item 阻尼绕组磁链 $\psi_D$、$\psi_Q$ 不变 $\Rightarrow$ 产生阻尼绕组自由分量（衰减很快）；
  \item 定子各相绕组磁链 $\psi_a$、$\psi_b$、$\psi_c$ 在短路瞬间不变 $\Rightarrow$ 定子电流出现
        非周期（直流）分量与倍频分量。
\end{itemize}

\subsection{无阻尼绕组同步电机突然三相短路的物理分析}

短路瞬间：定子电流突增，电枢反应磁通试图穿入转子，但被励磁绕组的感应电流所抵消，
磁通只能经励磁绕组漏磁路径闭合。此时从定子看进去的等值电抗为
\begin{equation}
  X_d'=X_\sigma+\frac{X_{ad}X_{f\sigma}}{X_{ad}+X_{f\sigma}},
\end{equation}
即“$X_{ad}$ 与励磁绕组漏抗 $X_{f\sigma}$ 并联后再串定子漏抗”，显然 $X_d'<X_d$。
对应地定义\textbf{暂态电势}
\begin{equation}
  E_q'=U_q+X_d'I_d ,
\end{equation}
它正比于励磁绕组磁链 $\psi_f$，故\textbf{短路瞬间 $E_q'$ 不变}——这使短路电流的初值
可以用短路前的运行状态直接算出。这就是“暂态电势”这个概念的全部价值。

\section{无阻尼绕组同步电机三相短路电流计算}

\subsection{暂态电势和暂态电抗}

由 $X_d'$、$E_q'$ 的定义可得，短路后定子电流的\textbf{周期分量}（强制分量）为
\keyformula{I'=\frac{E_q'}{X_d'+X_e}}
$X_e$ 为短路点至机端的外部电抗（含变压器、线路）。因为 $E_q'$ 在短路瞬间等于短路前值，
故 $I'$ 可由短路前运行状态一次算出，无需迭代。
稳态短路电流则为
\begin{equation}
  I_\infty=\frac{E_q}{X_d+X_e},
\end{equation}
$E_q$、$X_d$ 对应转子磁链完全穿入转子的稳态。

\begin{table}[htbp]
  \centering
  \caption{三个电抗、三个电势与三个电流阶段的对应（无阻尼绕组）}
  \small
  \begin{tabular}{@{}llll@{}}
    \toprule
    阶段 & 电抗 & 电势 & 电流 \\
    \midrule
    短路瞬间（初值） & $X_d'$ & $E_q'$（不变） & $I'=E_q'/(X_d'+X_e)$ \\
    暂态过程 & $X_d'\to X_d$ & $E_q'\to E_q$ & 周期分量由 $I'$ 衰减到 $I_\infty$ \\
    稳态 & $X_d$ & $E_q$ & $I_\infty=E_q/(X_d+X_e)$ \\
    \bottomrule
  \end{tabular}
\end{table}

\subsection{不计衰减时的短路电流算式}

不计各自由分量衰减时，定子电流可分解为三部分：
\begin{itemize}
  \item \textbf{基频周期分量}：幅值 $\sqrt{2}I'$，由 $E_q'$、$X_d'$ 决定，是“强制分量”；
  \item \textbf{非周期（直流）分量}：由定子绕组磁链守恒产生，初值使 $t=0$ 时总电流连续；
  \item \textbf{倍频分量}：凸极机 $X_d\ne X_q$ 引起，频率 $2\omega$，幅值较小。
\end{itemize}
三者之和在 $t=0$ 时必须等于短路前瞬间的电流值，这是确定自由分量初值的唯一条件。
工程实用计算中通常只关心基频周期分量的起始有效值 $I''$（即 $I'$ 在不计衰减的近似下），
以及由它派生的 $i_{sh}$、$I_{sh}$。

\subsection{自由电流的衰减}

各自由分量按不同时间常数衰减，形成短路电流的“三段式”特征：

\begin{table}[htbp]
  \centering
  \caption{短路电流各分量的时间常数}
  \small
  \begin{tabular}{@{}lll@{}}
    \toprule
    分量 & 时间常数 & 典型值 \\
    \midrule
    定子非周期分量 & $T_a\approx X''/(\omega r)$（近似 $L/R$） & $0.05\sim0.2$~s \\
    励磁绕组自由分量（暂态） & $T_d'=T_{d0}'\,X_d'/X_d$，$T_{d0}'=X_{ff}/(\omega r_f)$ & $0.5\sim2$~s \\
    阻尼绕组自由分量（次暂态） & $T_d''=T_{d0}''\,X_d''/X_d$ & $0.02\sim0.05$~s \\
    \bottomrule
  \end{tabular}
\end{table}

可见：阻尼绕组分量衰减最快（几个周期即消失），励磁绕组分量衰减较慢（数秒），
定子直流分量居中。三个时间常数相差一个数量级，正是“次暂态—暂态—稳态”分段处理的物理依据。

\subsection{负载状态下的突然短路}

负载短路与空载短路的区别在于\textbf{短路瞬间 $E_q'$ 的取值}：
负载时 $E_q'$ 由短路前的 $U$、$I$ 算出（$E_q'=U_q+X_d'I_d$），一般大于空载值；
同时定子非周期分量的初值由短路前负荷电流决定。实用计算中常取
$E_q'\approx U_0$ 或 $E''\approx1$（标幺），这是偏保守的简化。

\section{有阻尼绕组同步电机的突然三相短路}

\subsection{突然短路的物理过程}

有阻尼绕组时，短路瞬间阻尼绕组（$D$、$Q$）也满足磁链守恒，进一步“挡住”电枢反应磁通，
磁通只能走\textbf{定子漏磁、励磁绕组漏磁、阻尼绕组漏磁}三条并联路径。此时等值电抗为
\begin{equation}
  X_d''=X_\sigma+\left(\frac{1}{X_{ad}}+\frac{1}{X_{f\sigma}}+\frac{1}{X_{D\sigma}}\right)^{-1},
\end{equation}
显然 $X_d''<X_d'<X_d$。$q$ 轴同理有
\begin{equation}
  X_q''=X_\sigma+\frac{X_{aq}X_{Q\sigma}}{X_{aq}+X_{Q\sigma}}.
\end{equation}
过程分为三个阶段：\textbf{次暂态}（阻尼绕组起作用，$0\sim0.05$~s）、
\textbf{暂态}（阻尼绕组已衰减，励磁绕组起作用，$0.05$~s$\sim$数秒）、
\textbf{稳态}（转子磁链完全建立）。

\subsection{次暂态电势和次暂态电抗}

定义\textbf{次暂态电势}
\begin{equation}
  E''\approx U+\mathrm{j}X''\dot I ,
\end{equation}
其纵轴分量 $E_q''=U_q+X_d''I_d$，横轴分量 $E_d''=U_d-X_q''I_q$。
短路瞬间 $E''$ 不变（对应阻尼绕组磁链不变），于是
\keyformula{I''=\frac{E''}{X''+X_e}}
这就是短路电流周期分量起始有效值的\textbf{实用计算式}，也是第 6 章全部实用算法的基石。
工程上常取 $E''\approx1.0\sim1.1$（标幺），即“短路前电压约为额定值”。

\subsection{有阻尼绕组电机的短路电流}

计及各分量衰减，定子短路电流的完整表达式（一相）可写成
\begin{equation}
  i_a=\underbrace{\sqrt{2}I''e^{-t/T_d''}+\sqrt{2}I'e^{-t/T_d'}+\sqrt{2}I_\infty}_{\text{基频周期分量}}
  +\underbrace{\sqrt{2}I''e^{-t/T_a}\cos\alpha_0}_{\text{非周期分量}}
  +\underbrace{\cdots\cos 2\omega t}_{\text{倍频分量}} ,
\end{equation}
其中 $\alpha_0$ 为短路瞬间电压的相位角。三项周期分量的“接力”关系为
$I''>I'>I_\infty$，依次由 $T_d''$、$T_d'$ 控制衰减。

\begin{table}[htbp]
  \centering
  \caption{三个电抗、三个电势、三个电流的完整对应}
  \small
  \begin{tabular}{@{}llll@{}}
    \toprule
    阶段 & 电抗 & 电势 & 电流 \\
    \midrule
    次暂态（$0\sim0.05$~s） & $X_d''$（$X_q''$） & $E''$（不变） & $I''=E''/(X''+X_e)$ \\
    暂态（$0.05$~s$\sim$数秒） & $X_d'$ & $E_q'$ & $I'=E_q'/(X_d'+X_e)$ \\
    稳态 & $X_d$ & $E_q$ & $I_\infty=E_q/(X_d+X_e)$ \\
    \bottomrule
  \end{tabular}
\end{table}

\subsection{自由电流的衰减}

有阻尼绕组时多了一个时间常数 $T_d''$（阻尼绕组），它与 $T_a$、$T_d'$ 共同决定短路电流波形：
\begin{itemize}
  \item $t<0.05$~s：$I''$ 起主导，$i_{sh}$ 就在这个区间出现；
  \item $0.05\sim2$~s：$I'$ 段，继电保护的主要动作区间；
  \item $t>2$~s：接近 $I_\infty$，即稳态短路电流。
\end{itemize}

\section{强行励磁对短路暂态过程的影响}

\textbf{强行励磁}（继电强行励磁）指系统电压因短路而大幅下降时，励磁调节装置快速把励磁电压
升到顶值（可达额定励磁电压的 $1.8\sim2$ 倍甚至更高），以维持机端电压。

对短路暂态过程的影响：
\begin{enumerate}
  \item 提高 $E_q$，使短路电流周期分量的“稳态值”$I_\infty$ 上升，
        衰减过程被“托住”，短路电流在切除时刻可能明显大于不加强行励磁的情形；
  \item 提高发电机电磁功率，改善\textbf{暂态稳定}（第 17 章）——
        这是强行励磁最重要的作用；
  \item 对短路电流计算的影响是\textbf{偏保守}的：若不考虑强行励磁，算得的短路电流偏小，
        校验设备时须留裕度；但在稳定计算中必须计入，否则结果偏于乐观。
\end{enumerate}

\corepoint{本章的主线是一句话：\textbf{磁链守恒 $\Rightarrow$ 电势不变 $\Rightarrow$ 电流可算}。
短路瞬间哪个绕组的磁链守恒，就从定子看进去得到哪个电抗（$X''$、$X'$、$X_d$），
相应地就有哪个电势（$E''$、$E_q'$、$E_q$）不变，从而直接算出对应阶段的电流。
把这条链理顺，5.4、5.5 两节的十几个公式其实是三组同构的结果。}
  % 第 5 章
% ============================================================
% 第 6 章  电力系统三相短路电流的实用计算
% （教材：何仰赞、温增银《电力系统分析》上册 第 6 章）
% 结构对齐 6.1—6.4；例题、习题不录
% ============================================================

\chapter{电力系统三相短路电流的实用计算}

第 5 章解决了“短路电流随时间怎样变化”，本章解决“在给定系统中怎样算出来”。
实用计算有两条腿：\textbf{节点阻抗矩阵法}（计算机路线，精确、通用）与
\textbf{转移阻抗 $+$ 计算曲线法}（手算路线，查表、快速）。
两者的共同前提是：用次暂态电势 $E''$ 与次暂态电抗 $X_d''$ 作发电机的等值电路，
并把网络化为纯电抗网络（略电阻、略线路电容、略变压器励磁支路）。

\section{短路电流计算的基本原理和方法}

\subsection{电力系统节点方程的建立}

短路计算关心的是\textbf{起始次暂态电流}，因此各发电机用 $E''$、$X_d''$ 表示，
负荷一般用恒定阻抗（或按“略去”处理）。选定基准值后作等值网络，
按第 4 章的方法组装节点导纳矩阵 $\boldsymbol Y$，得节点方程
\begin{equation}
  \boldsymbol Y\dot{\boldsymbol U}=\dot{\boldsymbol I},
  \qquad \dot I_i=\frac{\dot E_i''}{jX_{di}''}\quad(\text{电源节点}),
  \qquad \dot I_i=0\quad(\text{无源节点}).
\end{equation}
求解前先在发电机节点“注入”各电源的短路电流，得到\textbf{短路前}的节点电压
$\dot U_i(0)$（也可直接取 $\dot U_i(0)\approx1\angle0^\circ$）。

\begin{table}[htbp]
  \centering
  \caption{短路计算中各类元件的处理}
  \small
  \begin{tabular}{@{}p{0.18\linewidth}p{0.34\linewidth}p{0.36\linewidth}@{}}
    \toprule
    元件 & 处理方式 & 理由 \\
    \midrule
    发电机 & 电势 $\dot E''$ 串 $jX_d''$ & 短路瞬间阻尼绕组磁链不变（第 5 章） \\
    变压器 & 只有电抗，励磁支路略去 & 励磁电流小 \\
    线路 & 只有电抗，并联电纳略去 & 短路计算中容性电流次要 \\
    负荷 & 恒定阻抗，或按“略去”处理 & 负荷电流远小于短路电流 \\
    无限大电源 & 内电抗为零，端电压恒定 & 系统容量远大于所研究部分 \\
    \bottomrule
  \end{tabular}
\end{table}

\subsection{利用节点阻抗矩阵计算短路电流}

设在节点 $f$ 发生短路，用一条阻值为 $z_f$ 的故障支路接地（$z_f=0$ 为金属性短路）。
短路后节点 $f$ 的电压为该支路两端电压，短路电流即流过 $z_f$ 的电流。
由叠加原理（把故障看作在 $f$ 点叠加一个 $-U_f(0)/z_f$ 的注入电流），可得
\keyformula{\dot I_f=-\frac{\dot U_f(0)}{Z_{ff}+z_f}}
金属性短路（$z_f=0$）时 $\dot I_f=-\dot U_f(0)/Z_{ff}$，其模为 $U_f(0)/Z_{ff}$。
各节点电压随之为
\begin{equation}
  \dot U_i=\dot U_i(0)-Z_{if}\dot I_f
  \qquad(i=1,2,\dots,n).
\end{equation}
仅需 $Z_{ff}$（输入阻抗）与一列 $Z_{if}$（转移阻抗），不必求全逆——这正是第 4.3 节
“只求一列”的动机。

若短路点不是原网络节点，可在该处增设一个节点（追加树支），再按上式算。

\subsection{利用电势源对短路点的转移阻抗计算短路电流}

把网络化简为“以短路点 $f$ 为中心、各电源点为顶点的星形电路”，则第 $i$ 个电源
对 $f$ 的支路电抗 $X_{if}$ 称\textbf{转移阻抗}，其电流为
\keyformula{\dot I_i=\frac{\dot E_i}{Z_{if}},\qquad \dot I_f=\sum_{i=1}^{m}\dot I_i}
\textbf{转移阻抗的物理意义}：除第 $i$ 个电源外其余电源电势为零时，
$f$ 点短路电流与 $E_i$ 之比；等价地，$Z_{if}$ 是“$i$ 点注入单位电流时 $f$ 点的电压”。

转移阻抗与输入阻抗的关系（同一网络、全电源同电势时两种算法结果必须相等）：
\begin{equation}
  \frac{1}{Z_{ff}}=\sum_{i=1}^{m}\frac{1}{Z_{if}} .
\end{equation}
这一恒等式是校验网络化简是否正确的常用手段。

\begin{table}[htbp]
  \centering
  \caption{两种算路方法的比较}
  \small
  \begin{tabular}{@{}p{0.16\linewidth}p{0.32\linewidth}p{0.34\linewidth}@{}}
    \toprule
    比较项 & 节点阻抗矩阵法 & 转移阻抗法 \\
    \midrule
    适用 & 复杂网络、多点短路、需各节点电压 & 手算、电源少、需分电源电流 \\
    关键量 & $Z_{ff}$ 与一列 $Z_{if}$ & 各 $Z_{if}$（星形化） \\
    优点 & 一次求逆可算所有节点短路 & 物理图像清晰、便于查计算曲线 \\
    缺点 & 需计算机 & 网络大时化简繁琐 \\
    \bottomrule
  \end{tabular}
\end{table}

\corepoint{“转移阻抗”与“输入阻抗”不是一回事：$Z_{ff}$ 是 $f$ 点看进去的等值阻抗，
$Z_{if}$ 是第 $i$ 个电源到 $f$ 的等值阻抗。二者只在单电源网络中数值相等。
多电源时若误用 $Z_{ff}$ 去算某一台机的电流，就会把其他电源的贡献也算进去。}

\section{起始次暂态电流和冲击电流的实用计算}

\subsection{起始次暂态电流}

实用计算中取各发电机 $\dot E''\approx1.0\sim1.1$（标幺），工程上常直接取 $1$，
于是短路点电流的标幺值为
\keyformula{I''_{*}=\frac{1}{X_{\Sigma *}},\qquad I''=I''_{*}I_B=\frac{I''_{*}S_B}{\sqrt{3}U_{av}}}
$X_{\Sigma*}$ 为从短路点看进去的等值电抗（标幺），$U_{av}$ 取短路点所在电压级的平均额定电压。

\textbf{关于负荷的处理}：负荷在短路瞬间可能“变成电源”向短路点供电（异步电动机的反馈电流），
判断依据是\textbf{残余电压}：
\begin{itemize}
  \item 若短路后该负荷节点残余电压 $U>0.8U_N$，负荷不供给短路电流，可略去；
  \item 若 $U<0.8U_N$，异步电动机可能倒送电流，应按 $E''\approx0.8$、$X''\approx0.2$ 的
        电动机模型计入（其冲击系数取 $1$）。
\end{itemize}

\subsection{冲击电流的实用计算}

\begin{equation}
  i_{sh}=K_{sh}\sqrt{2}I'' ,
\end{equation}
冲击系数按短路点位置与电源类型选取：

\begin{table}[htbp]
  \centering
  \caption{冲击系数的取法}
  \small
  \begin{tabular}{@{}lll@{}}
    \toprule
    短路点位置 & $K_{sh}$ & $i_{sh}$ \\
    \midrule
    发电机出口或发电机电压母线 & $1.90$ & $2.69I''$ \\
    高压电网（远离发电机） & $1.80$ & $2.55I''$ \\
    由异步电动机供给的分量 & $1.00$ & $\sqrt{2}I''$ \\
    \bottomrule
  \end{tabular}
\end{table}

继电保护整定和热稳定校验用 $I''$；动稳定校验用 $i_{sh}$ 与 $I_{sh}=1.52I''$（或 $1.62I''$）。

\section{短路电流计算曲线及其应用}

\subsection{计算曲线的概念}

由第 5 章，短路电流周期分量是许多参数的复杂函数：各发电机的电抗与时间常数、
反映短路前运行状态的电势初值、励磁系统参数、短路点离机端的距离、时间 $t$。
在发电机（含励磁系统）参数和运行初态给定后，它只是“电气距离”和“时间”的函数。
把\textbf{归算到发电机额定容量}的外接电抗标幺值与发电机纵轴次暂态电抗标幺值之和定义为
\textbf{计算电抗}
\keyformula{X_{js}=X_d''+X_e}
则短路电流周期分量的标幺值可写成
\begin{equation}
  I_{p*}=f\left(X_{js},t\right),
\end{equation}
反映这一函数关系的一族曲线（或数字表）即\textbf{计算曲线}。

\subsection{计算曲线的制作条件}

计算曲线不是严格的理论结果，而是在一组\textbf{典型条件}下对多台样机逐台计算再取平均得到的。
用曲线前必须知道它是在什么条件下做出来的，否则“查表”就失去依据：

\begin{table}[htbp]
  \centering
  \caption{计算曲线的制作条件}
  \small
  \begin{tabular}{@{}p{0.20\linewidth}p{0.68\linewidth}@{}}
    \toprule
    项目 & 条件 \\
    \midrule
    典型接线 & 发电机经变压器接到高压母线，$50\%$ 负荷接于发电厂高压母线，
               其余负荷经线路送到短路点以外 \\
    短路前状态 & 发电机额定满载运行 \\
    负荷模型 & 恒定阻抗 $Z_{LD}=U^{2}/\left(S_{LD}(\cos\varphi+j\sin\varphi)\right)$，
               取 $U=1$、$\cos\varphi=0.9$ \\
    励磁 & 均配强行励磁，顶值电压为额定励磁电压的 $1.8$ 倍；
           励磁系统等值时间常数 $T_e$：汽轮发电机 $0.25$~s，水轮发电机 $0.02$~s \\
    样机 & 汽轮发电机取 $12\sim200$~MW 的 18 种型号；水轮发电机取 $12.5\sim225$~MW 的 17 种型号；
           对给定 $X_{js}$ 与 $t$ 取各机计算结果的算术平均值 \\
    适用范围 & 只做到 $X_{js}=3.45$；$X_{js}\ge3.45$ 时认为周期分量幅值不随时间变化，
               直接取 $I_{p*}=1/X_{js}$ \\
    \bottomrule
  \end{tabular}
\end{table}

制作曲线的网络只含\textbf{一台}发电机，且计算电抗与负荷支路无关。因此实际系统用曲线前，
必须先把网络化为“以短路点为中心、以各电源点为顶点的星形电路”。

\subsection{计算曲线的应用}

\paragraph{（1）绘制等值网络}
取 $S_B$、$U_B=U_{av}$；发电机电抗用 $X_d''$；略去各元件电阻、线路电容和变压器励磁支路；
无限大电源内电抗取零；略去负荷。

\paragraph{（2）网络变换、合并电源}
按“短路电流变化规律相近”的原则合并：与短路点电气距离相近的\textbf{同类型}发电机可合并；
远离短路点的同类型电厂可合并；\textbf{直接接于短路点的发电机应单独考虑}；
无限大电源必须单独计算（其电流不衰减）。合并后得 $g$ 组等值机，求出各组的转移电抗
$X_{if}$ 与无限大电源的 $X_{fs}$。

\paragraph{（3）归算计算电抗}
\begin{equation}
  X_{jsi}=X_{if}\frac{S_{Ni}}{S_B}\qquad(i=1,2,\dots,g),
\end{equation}
$S_{Ni}$ 为第 $i$ 台等值机所代表的那部分发电机额定容量之和。

\paragraph{（4）查曲线}
由 $X_{jsi}$ 查相应类型（汽轮/水轮）的计算曲线，得指定时刻 $t$ 的标幺值 $I_{pti*}$。

\paragraph{（5）无限大电源的分量}
\begin{equation}
  I_{ps*}=\frac{1}{X_{fs}} .
\end{equation}

\paragraph{（6）换算为有名值}
\begin{equation}
  I_{ti}=I_{pti*}\frac{S_{Ni}}{\sqrt{3}U_{av}},\qquad
  I_{s}=I_{ps*}\frac{S_B}{\sqrt{3}U_{av}},\qquad
  I_t=\sum_{i=1}^{g}I_{ti}+I_s ,
\end{equation}
其中 $U_{av}$ 取短路处电压级的平均额定电压，$S_{Ni}/(\sqrt{3}U_{av})$ 即归算到短路处电压级的
第 $i$ 台等值机额定电流。

\begin{table}[htbp]
  \centering
  \caption{手算路线的三个“归算”易混点}
  \small
  \begin{tabular}{@{}p{0.24\linewidth}p{0.30\linewidth}p{0.38\linewidth}@{}}
    \toprule
    量 & 归算到什么容量 & 用途 \\
    \midrule
    转移电抗 $X_{if}$ & 统一基准 $S_B$ & 描述网络拓扑，用于合并化简 \\
    计算电抗 $X_{jsi}$ & 各电源自己的 $S_{Ni}$ & 查计算曲线（曲线的横坐标） \\
    电流标幺值 $I_{pti*}$ & 各电源自己的 $S_{Ni}$ & 乘 $S_{Ni}/(\sqrt{3}U_{av})$ 得有名字 \\
    \bottomrule
  \end{tabular}
\end{table}

\corepoint{计算曲线的横坐标是“按本机容量归算的”计算电抗，而网络化简得到的转移电抗是
“按 $S_B$ 归算的”。二者差一个 $S_{Ni}/S_B$ 的因子，漏掉这一步是最常见的错误。
另外 $X_{js}\ge3.45$ 时曲线终止，此时直接取 $1/X_{js}$，不要再外推查表。}

\section{短路电流周期分量的近似计算}

工程估算中可假定短路电路接于\textbf{内阻抗为零的恒电势电源}，于是周期分量幅值不随时间变化，
只有非周期分量衰减。取各电源电势标幺值为 $1$，则
\keyformula{I_{p*}=\frac{1}{X_{ff*}},\qquad I_p=I_{p*}I_B=\frac{I_{p*}S_B}{\sqrt{3}U_{av}},\qquad
  S_k=\frac{S_B}{X_{ff*}}}
$X_{ff*}$ 为短路点的输入电抗（标幺）。这样算出的短路电流（或短路功率）比实际偏大，
但差别随短路点距离增大而迅速减小——因为短路点越远，“电源电压恒定”的假设越接近实际，
尤其当发电机装有自动励磁调节器时更是如此。故这一简化可用来估计短路电流
（或短路功率）的\textbf{最大可能值}。

\paragraph{未知系统参数的估算}
若只知某母线处的短路容量 $S_s$（或有名值电流 $I_s$），可把该系统等效为经电抗 $X_s$
接于该母线的无限大电源：
\begin{equation}
  X_{s*}=\frac{S_B}{S_s}\qquad\left(\text{或由}\ S_s=\sqrt{3}U_{av}I_s\ \text{反算}\right).
\end{equation}
若连短路容量也不知道，可由该处变电所断路器的\textbf{额定切断容量}反推：
认为断路器后三相短路时其切断容量恰好被充分利用，即 $S_s$ 取断路器额定切断容量。

\begin{table}[htbp]
  \centering
  \caption{本章四种算法的适用场合}
  \small
  \begin{tabular}{@{}p{0.22\linewidth}p{0.30\linewidth}p{0.36\linewidth}@{}}
    \toprule
    方法 & 精度 & 适用 \\
    \midrule
    节点阻抗矩阵法 & 高（取决于 $E''$ 取值） & 计算机计算、需全网点电压 \\
    转移阻抗 $+$ 计算曲线 & 中（受制作条件限制） & 手算、需分电源、需任意时刻电流 \\
    近似计算（$E=1$） & 偏低精度、偏保守 & 估算最大短路电流、初选设备 \\
    按短路容量反推系统电抗 & 粗 & 缺少系统详细数据时 \\
    \bottomrule
  \end{tabular}
\end{table}

\corepoint{本章的核心是一条计算链：
\textbf{选基准 $\to$ 作纯电抗等值网络 $\to$ 化简到短路点 $\to$ 求 $X_{\Sigma}$（或 $X_{js}$）
$\to$ 得 $I''$ $\to$ 派生 $i_{sh}$、$I_{sh}$、$S_k$}。
链条上每一环都只做一件事；出错往往不在物理，而在“归算到哪个容量”“用哪一电压级的
$U_{av}$”这两个纯约定问题上。}
  % 第 6 章
% ============================================================
% 第 7 章  电力系统各元件的序阻抗和等值电路
% （教材：何仰赞、温增银《电力系统分析》上册 第 7 章）
% 结构对齐 7.1—7.7；例题、习题不录
% ============================================================

\chapter{电力系统各元件的序阻抗和等值电路}

不对称短路不能用一相电路直接分析，但如果把不对称的三相量分解为三组对称分量，
则对每一组对称分量，网络仍是三相对称的，可以用单相等值电路处理。
本章做两件事：把\textbf{对称分量法}装到短路计算上；给出各元件在\textbf{正序、负序、零序}
三种情况下的等值电路与参数——这是第 8 章全部不对称故障计算的物质基础。

\section{对称分量法在不对称短路计算中的应用}

\subsection{不对称三相量的分解}

任一组不对称三相相量（电压或电流）都可唯一分解为三组对称分量：
\textbf{正序}（$a\to b\to c$ 依次滞后 $120^\circ$）、\textbf{负序}（相序相反）、
\textbf{零序}（三相同相）。引入算子
\begin{equation}
  a=e^{j120^\circ}=-\tfrac{1}{2}+j\tfrac{\sqrt{3}}{2},\qquad
  a^{2}=e^{j240^\circ}=e^{-j120^\circ},\qquad 1+a+a^{2}=0,
\end{equation}
则
\keyformula{
  \begin{bmatrix}\dot F_a\\ \dot F_b\\ \dot F_c\end{bmatrix}
  =\begin{bmatrix}1&1&1\\ a^{2}&a&1\\ a&a^{2}&1\end{bmatrix}
  \begin{bmatrix}\dot F_{a1}\\ \dot F_{a2}\\ \dot F_{a0}\end{bmatrix},
  \qquad
  \begin{bmatrix}\dot F_{a1}\\ \dot F_{a2}\\ \dot F_{a0}\end{bmatrix}
  =\frac{1}{3}\begin{bmatrix}1&a&a^{2}\\ 1&a^{2}&a\\ 1&1&1\end{bmatrix}
  \begin{bmatrix}\dot F_a\\ \dot F_b\\ \dot F_c\end{bmatrix}}
即 $\dot F_{a1}=\frac{1}{3}(\dot F_a+a\dot F_b+a^{2}\dot F_c)$，
$\dot F_{a2}=\frac{1}{3}(\dot F_a+a^{2}\dot F_b+a\dot F_c)$，
$\dot F_{a0}=\frac{1}{3}(\dot F_a+\dot F_b+\dot F_c)$。
零序分量三相相等，故\textbf{零序电流必须有通路}才能存在（中性点接地、三角形回路）。

\subsection{序阻抗的概念}

对线性、三相对称的元件（或网络），三序分量彼此\textbf{解耦}：正序电压只产生正序电流，
负序、零序亦然。故可对每一序单独定义阻抗：
\begin{equation}
  \dot U_{1}=Z_{1}\dot I_{1},\qquad \dot U_{2}=Z_{2}\dot I_{2},\qquad \dot U_{0}=Z_{0}\dot I_{0}.
\end{equation}
$Z_1$、$Z_2$、$Z_0$ 分别称\textbf{正序、负序、零序阻抗}。
\begin{itemize}
  \item 静止元件（线路、变压器、电抗器）：$Z_1=Z_2$（正、负序所见的磁路完全相同，
        只是相序相反）；
  \item 旋转元件（同步电机、异步电动机）：$Z_1\ne Z_2$（负序磁场与转子转向相反，
        相对转速为 $2n_1$，转子绕组对负序磁场的作用类似短路绕组，故 $Z_2$ 远小于 $Z_1$）；
  \item 零序阻抗与前三者都可能不同，且强烈依赖\textbf{接线方式与接地情况}。
\end{itemize}

\subsection{对称分量法在不对称短路计算中的应用}

做法分三步：
\begin{enumerate}
  \item \textbf{分解}：把故障处的三相不对称边界条件（电压、电流）用对称分量表示；
  \item \textbf{分序}：对正、负、零序分别作出等值网络（序网），在故障端口列出序网方程；
  \item \textbf{合成}：由边界条件与序网方程联立求出各序分量，再合成为相量，
        求各相电流、电压。
\end{enumerate}

\begin{table}[htbp]
  \centering
  \caption{三序网络的一般结构}
  \small
  \begin{tabular}{@{}p{0.10\linewidth}p{0.30\linewidth}p{0.48\linewidth}@{}}
    \toprule
    序网 & 电源 & 特点 \\
    \midrule
    正序 & 有（各发电机 $\dot E''$） & 与三相短路网络同构；故障点为正序电流的“出口” \\
    负序 & 无（发电机负序电势为零） & 无源网络；故障点可视为负序电流的“注入点” \\
    零序 & 无 & 只有零序电流通路存在时才成立；网络结构与接地方式强相关 \\
    \bottomrule
  \end{tabular}
\end{table}

\corepoint{三序网络只有\textbf{正序}有电源，负序、零序都是无源网络——这是全部不对称短路
计算的结构性事实。因此正序网中的故障端口电势就是 $E_{\Sigma}$（或 $\dot U_f(0)$），
负序、零序网在端口上只表现为“从故障点看进去的等值阻抗”。}

\section{同步发电机的负序和零序电抗}

\textbf{负序电抗 $X_2$}：负序电流产生的磁场与转子反向旋转，相对速度为 $2n_1$，
在转子绕组中感应出 $2f$ 的电流，其作用使负序磁通被“挤”到漏磁路径，
故 $X_2$ 接近 $X_d''$、$X_q''$ 的某种平均值：
\begin{equation}
  X_2\approx\frac{1}{2}\left(X_d''+X_q''\right)
\end{equation}
实用计算中常取 $X_2\approx1.22X_d''$（汽轮发电机）或按制造厂数据。
对无阻尼绕组电机，$X_2\approx\sqrt{X_d'X_q}$ 一类的近似式也可用，但精度较低。

\textbf{零序电抗 $X_0$}：零序电流三相同相，合成磁势为零（只在定子绕组漏磁路径中闭合），
故 $X_0$ 很小，且与绕组节距、接线有关，一般
\begin{equation}
  X_0=(0.15\sim0.6)X_d'' ,
\end{equation}
具体值须查制造厂数据。零序电流不能在发电机内部流通时（如中性点不接地），
$X_0$ 无意义，零序网在该处断开。

\begin{table}[htbp]
  \centering
  \caption{同步电机三序电抗的数量关系}
  \small
  \begin{tabular}{@{}lll@{}}
    \toprule
    电抗 & 相对大小 & 物理原因 \\
    \midrule
    $X_d$ & 最大 & 稳态，电枢反应磁通完全穿入转子 \\
    $X_d'$ & 中 & 励磁绕组挡住磁通，走漏磁路径 \\
    $X_d''$ & 小 & 阻尼绕组也挡住磁通，三路漏磁并联 \\
    $X_2$ & 与 $X_d''$、$X_q''$ 同量级 & 负序磁场相对转子 $2n_1$，转子绕组近似短路 \\
    $X_0$ & 最小 & 零序合成磁势为零，只有定子漏磁 \\
    \bottomrule
  \end{tabular}
\end{table}

\section{变压器的零序等值电路及其参数}

变压器的正序、负序等值电路与参数和正常运行时完全相同（静止元件 $Z_1=Z_2$），
\textbf{零序则完全不同}，原因有二：
\begin{enumerate}
  \item 零序电流三相同相，能否流通取决于绕组接线方式（Y 不接地则无通路，$\Delta$ 内可环流，
        Y$_0$ 经中性点入地）；
  \item 零序磁通路径与正序不同——三柱式铁芯中零序磁通不能经铁芯闭合，
        须经油、油箱壁等结构件闭合，磁阻大、损耗大，故零序励磁阻抗远小于正序励磁阻抗。
\end{enumerate}

\subsection{普通变压器的零序等值电路及其参数}

零序等值电路的形式（三支路：一次漏抗、励磁支路、二次漏抗）与正序相同，
但有三点区别：
\begin{itemize}
  \item 各绕组漏抗 $X_1$、$X_2$ 与正序相同（漏磁路径不受相序影响）；
  \item 零序励磁电抗 $X_{\mu0}$ 远小于正序励磁电抗：三柱式铁芯 $X_{\mu0}\approx0.3\sim1.0$
        （标幺，以绕组容量为基准），五柱式或壳式可接近正序值；
  \item 零序电流能否流过某个绕组，取决于该绕组的接线与中性点接地情况。
\end{itemize}

\subsection{变压器零序等值电路与外电路的联接}

把零序等值电路“挂”到序网上时，用一个等效开关表示各侧的通断：

\begin{table}[htbp]
  \centering
  \caption{变压器各侧接线方式对零序通断的规定}
  \small
  \begin{tabular}{@{}lll@{}}
    \toprule
    绕组接线 & 零序电流能否流出（与外电路相接） & 等值处理 \\
    \midrule
    Y 中性点接地（Y$_0$） & 能（经中性点入地） & 该侧与外电路接通 \\
    Y 中性点不接地 & 不能（无通路） & 该侧开路 \\
    $\Delta$（三角形） & 不能流出，但可在内部环流 & 该侧对零序相当于\textbf{短路接地}（端口电压为零） \\
    Y$_0$/$\Delta$ 组合 & 零序电流在 $\Delta$ 内环流，不流出 & $\Delta$ 侧短接，Y$_0$ 侧接通 \\
    \bottomrule
  \end{tabular}
\end{table}

因此常用“三句话”判断：\textbf{Y$_0$ 通、Y 断、$\Delta$ 短}（对零序而言）。

\subsection{中性点有接地阻抗时变压器的零序等值电路}

若中性点经阻抗 $Z_n$ 接地，零序电流 $3\dot I_0$ 流经 $Z_n$，在中性点产生的电位升高为
$3\dot I_0Z_n$。折算到每相，等值电路中串入
\begin{equation}
  Z_n'=3Z_n .
\end{equation}
\keyformula{Z_{n}'=3Z_{n}}
“三倍”来自零序电流三相之和为 $3I_0$；这个 $3$ 倍关系是不对称短路计算中出现频率最高的
系数之一，与后面“零序网中线路阻抗取 $3Z_m$”同源。

\subsection{自耦变压器的零序等值电路及其参数}

自耦变压器的零序等值电路与普通三绕组变压器形式相同，但有三点特殊：
\begin{enumerate}
  \item 高、中压绕组间有\textbf{直接电气联系}，零序电流可以直接从高压侧传到中压侧，
        因此零序通断不再只看“是否 Y$_0$”，还要看高、中压侧是否都接地；
  \item 第三绕组（低压）总是接成 $\Delta$，为零序电流提供环流通路，同时消除三次谐波；
  \item 若高、中压侧中性点均接地，则零序电流可在两侧流通，等值电路须按三绕组方式
        （三支路星形）处理，且要考虑中性点接地阻抗的 $3Z_n$（分别在高、中压侧）。
\end{enumerate}

\corepoint{变压器零序等值电路有三件事必须同时交代：\textbf{接线方式}（Y$_0$／Y／$\Delta$）、
\textbf{中性点是否接地及接地阻抗}、\textbf{铁芯结构}（决定 $X_{\mu0}$ 大小）。
缺任何一条，零序网就画不对。考试中“画出零序网络”失分，多半是漏了 $\Delta$ 侧短接。}

\section{架空输电线路的零序阻抗及其等值电路}

\subsection{单导线-大地回路的自阻抗和互阻抗}

零序电流经导线流向故障点、经大地返回，故必须考虑大地回路。引入\textbf{大地等值深度}
\begin{equation}
  D_e=660\sqrt{\frac{\rho}{f}}\quad(\text{m}),
\end{equation}
$\rho$ 为大地电阻率（$\Omega\cdot$m），$f$ 为频率。则单导线-大地回路的自阻抗
\begin{equation}
  Z_s=\left(r_a+r_g\right)+j0.1445\lg\frac{D_e}{D_s}\quad(\Omega/\text{km}),
\end{equation}
两根导线-大地回路间的互阻抗
\begin{equation}
  Z_m=r_g+j0.1445\lg\frac{D_e}{D}\quad(\Omega/\text{km}),
\end{equation}
其中 $r_a$ 为导线电阻，$D_s=0.779r$ 为导线自几何均距，$D$ 为两导线间距，
$r_g$ 为大地电阻
\begin{equation}
  r_g=\pi^{2}f\times10^{-4}\ \Omega/\text{km}\approx0.05\ \Omega/\text{km}\ (f=50\ \text{Hz}).
\end{equation}

\subsection{三相输电线路的零序阻抗}

零序时三相电流同相，与 $a$ 相交链的磁通由本相电流与 $b$、$c$ 相电流共同产生，
且三者同相叠加（不像正序那样相互抵消），故
\keyformula{Z_0=Z_s+2Z_m
  =\left(r_a+3r_g\right)+j0.1445\lg\frac{D_e^{3}}{D_sD_{eq}^{2}}\quad(\Omega/\text{km})}
与正序阻抗相比：
\begin{itemize}
  \item 零序磁通经大地返回，磁路磁阻大，故 $Z_0>Z_1$，工程上\textbf{单回架空线}
        $Z_0\approx(3.0\sim3.5)Z_1$；
  \item 有良导体架空地线或电缆时，零序电流可经地线返回，$Z_0$ 明显下降；
  \item 零序电抗与正序电抗之比记作 $k=Z_0/Z_1$，是不对称短路计算中的重要参数。
\end{itemize}

\subsection{平行架设的双回输电线路的零序阻抗及等值电路}

双回线并行时，两回线之间不仅有各相间的互感，还有\textbf{回路间}的零序互感 $Z_{m0}$。
零序网中不能把两回线简单并联，须用两端口方程
\begin{equation}
  \begin{bmatrix}\dot U_{0\text{I}}\\ \dot U_{0\text{II}}\end{bmatrix}
  =\begin{bmatrix}Z_0 & Z_{m0}\\ Z_{m0} & Z_0\end{bmatrix}
  \begin{bmatrix}\dot I_{0\text{I}}\\ \dot I_{0\text{II}}\end{bmatrix}
\end{equation}
处理，或用第 4.1.4 节的 $\boldsymbol Y=\boldsymbol A\boldsymbol Z_b^{-1}\boldsymbol A^{\mathsf T}$
组装。当两回线零序电流方向相同时（同向并联），等值阻抗为
$Z_0+Z_{m0}$ 的一半形式；方向相反时（一回线检修接地），互感使等值阻抗增大。

\subsection{架空地线对输电线路零序阻抗及等值电路的影响}

架空地线（避雷线）经杆塔接地，为零序电流提供并联通路，且与三相导线间有互感。
其效果是\textbf{降低}线路的零序阻抗（相当于在零序网中并联一条阻抗支路），
减小系数约为 $10\%\sim15\%$（钢绞线）或更多（良导体地线）。
精确计算时把地线作为独立支路，用互感耦合的三支路方程处理；工程估算可按经验系数折减。

\begin{table}[htbp]
  \centering
  \caption{线路零序阻抗的估算与影响因素}
  \small
  \begin{tabular}{@{}p{0.24\linewidth}p{0.26\linewidth}p{0.40\linewidth}@{}}
    \toprule
    情形 & $Z_0/Z_1$ & 说明 \\
    \midrule
    单回架空线（无地线） & $3.0\sim3.5$ & 常用估算值 \\
    有架空地线 & 约降 $10\%\sim15\%$ & 地线提供并联返回通路 \\
    电缆线路 & 远小于 $3.5$，可接近 $1$ & 铅包、铠装提供低阻返回通路 \\
    平行双回线 & 与两回线电流方向有关 & 互感耦合，不能简单并联 \\
    \bottomrule
  \end{tabular}
\end{table}

\section{架空输电线路的零序电纳}

（选学内容）零序电纳的计算与正序类似，但零序时三相电荷同相，电场分布不同，
且架空地线的镜像会影响电场。无架空地线时零序电容按“三相导线 $+$ 大地 $+$ 镜像”系统计算；
有地线时须把地线及其镜像一并计入。工程上零序电纳对短路电流的影响较小，
一般在零序网中按经验值取用；只有在需要精确计算零序电流分布时才详细计算。

\section{综合负荷的序阻抗}

综合负荷以异步电动机为主。
\begin{itemize}
  \item \textbf{正序阻抗}：按额定电压、额定功率因数下的等值阻抗计，标幺值约 $1.2$；
  \item \textbf{负序阻抗}：负序磁场相对转子以 $2n_1$ 旋转，转子绕组（笼型）近似短接，
        故负序阻抗远小于正序阻抗，标幺值约 $0.2\sim0.35$；
  \item \textbf{零序阻抗}：取决于负荷的接线与接地方式；若负荷接于 $\Delta$ 或中性点不接地，
        零序网在该处断开（$Z_0\to\infty$）；若中性点接地，则零序阻抗与负序阻抗同量级。
\end{itemize}
短路实用计算中，一般\textbf{略去负荷}（除短路点附近的大电动机外），
这使短路电流略偏大，属于偏保守的处理。

\section{电力系统各序网络的制订}

\subsection{正序网络}

正序网络与三相短路时的网络完全一样：
\begin{itemize}
  \item 发电机用 $\dot E''$ 串 $jX_d''$（或按精度要求用 $\dot E_q'$ 串 $jX_d'$）；
  \item 变压器、线路用正序参数，负荷可略去或按恒定阻抗；
  \item 故障点作为正序网的端口，端口电势为 $\dot U_f(0)$（短路前该点电压）。
\end{itemize}

\subsection{负序网络}

负序网络与正序网络\textbf{结构相同，但所有电源电势为零}（发电机负序电势为零），
发电机用 $jX_2$ 代替 $jX_d''$，其余元件用负序参数（静止元件 $Z_2=Z_1$）。
故负序网是\textbf{无源网络}，从故障点看进去的等值阻抗记为 $Z_{2\Sigma}$。

\subsection{零序网络}

零序网络的制订最需要小心，步骤是：
\begin{enumerate}
  \item 从故障点出发，凡\textbf{零序电流能流通}的路径才画入；
  \item 变压器按“Y$_0$ 通、Y 断、$\Delta$ 短”处理（$\Delta$ 侧在零序网中接地）；
  \item 中性点接地阻抗按 $3Z_n$ 串入；
  \item 线路用 $Z_0$，平行双回线计及零序互感，有架空地线时计入地线支路；
  \item 发电机中性点不接地时，零序网在发电机处断开（$X_0$ 不出现）。
\end{enumerate}
零序网也是无源网络，从故障点看进去的等值阻抗记为 $Z_{0\Sigma}$。
若故障点与电源之间没有零序通路（如全部变压器均为 Y/$\Delta$ 且中性点不接地），
则 $Z_{0\Sigma}\to\infty$，零序电流为零，该故障不产生零序分量。

\begin{table}[htbp]
  \centering
  \caption{三序网络制订要点对照}
  \small
  \begin{tabular}{@{}p{0.09\linewidth}p{0.16\linewidth}p{0.26\linewidth}p{0.34\linewidth}@{}}
    \toprule
    序网 & 电源 & 发电机 & 关键处理 \\
    \midrule
    正序 & 有 & $\dot E''$ 串 $jX_d''$ & 与三相短路同构，端口电势为 $\dot U_f(0)$ \\
    负序 & 无 & $jX_2$ & 结构同正序，全部电势置零 \\
    零序 & 无 & $jX_0$（中性点接地时） & 逐条检查零序通路；$\Delta$ 短接、Y 断开、$3Z_n$ \\
    \bottomrule
  \end{tabular}
\end{table}

\corepoint{画零序网的唯一标准是“零序电流有没有路走”。判断顺序：
从故障点往外走 $\to$ 遇到变压器看接线 $\to$ 看中性点是否接地 $\to$ 看线路有无地线。
三序网画对之后，第 8 章的各种故障公式只是“在三个序网的端口上接不同的边界条件”，
不需要再记新的网络知识。}
  % 第 7 章
% ============================================================
% 第 8 章  电力系统不对称故障的分析和计算
% （教材：何仰赞、温增银《电力系统分析》上册 第 8 章）
% 结构对齐 8.1—8.5；例题、习题不录
% ============================================================

\chapter{电力系统不对称故障的分析和计算}

有了第 7 章的三序网络，不对称故障的计算就变成一套固定套路：
\textbf{在故障端口列出边界条件 $\to$ 化为序分量 $\to$ 接入三序网 $\to$ 解出各序电流电压
$\to$ 合成相量}。本章把四种简单不对称短路、非全相断线都用这一套路走一遍，
最后归纳为“正序等效定则”与“计算通式”，使不同故障之间的差别只剩一个附加阻抗 $Z_\Delta$。

\section{简单不对称短路的分析}

设故障发生前故障点电压为 $\dot U_f(0)$，三序网从故障点看进去的等值阻抗分别为
$Z_{1\Sigma}$、$Z_{2\Sigma}$、$Z_{0\Sigma}$；以 $a$ 相为特殊相（故障相或非故障相的基准相）。

\subsection{单相（$a$ 相）接地短路}

\textbf{边界条件}：$b$、$c$ 相无电流，$a$ 相电压为零
\begin{equation}
  \dot I_b=\dot I_c=0,\qquad \dot U_a=0 .
\end{equation}
化为序分量：
\begin{equation}
  \dot I_{a1}=\dot I_{a2}=\dot I_{a0}=\tfrac{1}{3}\dot I_a,
  \qquad \dot U_{a1}+\dot U_{a2}+\dot U_{a0}=0 .
\end{equation}
即三序电流相等、三序电压之和为零——在序网中相当于\textbf{三个序网串联}。
代入序网方程
$\dot U_{a1}=\dot U_f(0)-Z_{1\Sigma}\dot I_{a1}$、$\dot U_{a2}=-Z_{2\Sigma}\dot I_{a2}$、
$\dot U_{a0}=-Z_{0\Sigma}\dot I_{a0}$，得
\keyformula{\dot I_{a1}=\frac{\dot U_f(0)}{Z_{1\Sigma}+Z_{2\Sigma}+Z_{0\Sigma}},
  \qquad \dot I_f=\dot I_a=3\dot I_{a1}}
故障相电流为三倍正序电流，这是单相接地最显著的特征。

\subsection{两相（$b$ 相和 $c$ 相）短路}

\textbf{边界条件}：$a$ 相无电流，$b$、$c$ 相电流相反、电压相等
\begin{equation}
  \dot I_a=0,\qquad \dot I_b=-\dot I_c,\qquad \dot U_b=\dot U_c .
\end{equation}
化为序分量：
\begin{equation}
  \dot I_{a0}=0,\qquad \dot I_{a1}=-\dot I_{a2},\qquad \dot U_{a1}=\dot U_{a2}.
\end{equation}
零序无通路（$I_0=0$），正、负序网\textbf{并联}。于是
\keyformula{\dot I_{a1}=\frac{\dot U_f(0)}{Z_{1\Sigma}+Z_{2\Sigma}},\qquad
  \dot I_b=-j\sqrt{3}\dot I_{a1},\qquad I_f=\sqrt{3}\left|\dot I_{a1}\right|}
两相短路电流为正序电流的 $\sqrt{3}$ 倍，且比同一点三相短路电流小
（$I^{(2)}=\frac{\sqrt{3}Z_{1\Sigma}}{Z_{1\Sigma}+Z_{2\Sigma}}I^{(3)}$，
当 $Z_{1\Sigma}\approx Z_{2\Sigma}$ 时 $I^{(2)}\approx0.866I^{(3)}$）。

\subsection{两相（$b$ 相和 $c$ 相）短路接地}

\textbf{边界条件}：$a$ 相无电流，$b$、$c$ 相电压为零
\begin{equation}
  \dot I_a=0,\qquad \dot U_b=\dot U_c=0 .
\end{equation}
化为序分量：
\begin{equation}
  \dot I_{a1}+\dot I_{a2}+\dot I_{a0}=0,\qquad
  \dot U_{a1}=\dot U_{a2}=\dot U_{a0}.
\end{equation}
三序网\textbf{并联}。于是
\keyformula{\dot I_{a1}=\frac{\dot U_f(0)}{Z_{1\Sigma}+Z_{2\Sigma}\parallel Z_{0\Sigma}},
  \qquad Z_{2\Sigma}\parallel Z_{0\Sigma}=\frac{Z_{2\Sigma}Z_{0\Sigma}}{Z_{2\Sigma}+Z_{0\Sigma}}}
两故障相电流为
\begin{equation}
  \dot I_b=\dot I_{a0}+a^{2}\dot I_{a1}+a\dot I_{a2},\qquad
  \dot I_c=\dot I_{a0}+a\dot I_{a1}+a^{2}\dot I_{a2},
\end{equation}
其模相等，$I_f=\sqrt{3}\sqrt{1-\dfrac{Z_{2\Sigma}Z_{0\Sigma}}{\left(Z_{2\Sigma}+Z_{0\Sigma}\right)^{2}}}\left|\dot I_{a1}\right|$。
流入大地的电流为 $3\dot I_{a0}$。

\begin{table}[htbp]
  \centering
  \caption{四种简单不对称短路的边界条件与序网联接方式}
  \small
  \begin{tabular}{@{}llll@{}}
    \toprule
    故障类型 & 边界条件 & 序网联接 & 故障电流 \\
    \midrule
    三相短路 & $\dot U_a=\dot U_b=\dot U_c=0$ & 只有正序网 & $I_f=I_{a1}$ \\
    两相短路 & $\dot I_a=0$，$\dot I_b=-\dot I_c$，$\dot U_b=\dot U_c$ & 正、负序并联 & $I_f=\sqrt{3}I_{a1}$ \\
    单相接地 & $\dot I_b=\dot I_c=0$，$\dot U_a=0$ & 三序串联 & $I_f=3I_{a1}$ \\
    两相接地 & $\dot I_a=0$，$\dot U_b=\dot U_c=0$ & 三序并联 & $I_f=mI_{a1}$ \\
    \bottomrule
  \end{tabular}
\end{table}

\subsection{正序等效定则}

把上述四式写成同一形式：
\keyformula{\dot I_{a1}=\frac{\dot U_f(0)}{Z_{1\Sigma}+Z_{\Delta}},\qquad I_f=m\left|\dot I_{a1}\right|}
即：\textbf{任何不对称短路，其正序电流等于“在正序网故障点串入一个附加阻抗 $Z_\Delta$ 后
发生三相短路”的正序电流}。$Z_\Delta$ 与 $m$ 只由故障类型决定：

\begin{table}[htbp]
  \centering
  \caption{正序等效定则中的附加阻抗与电流倍数}
  \label{tab:zDelta}
  \small
  \begin{tabular}{@{}lll@{}}
    \toprule
    故障类型 & 附加阻抗 $Z_\Delta$ & 倍数 $m$ \\
    \midrule
    三相短路 & $0$ & $1$ \\
    两相短路 & $Z_{2\Sigma}$ & $\sqrt{3}$ \\
    单相接地短路 & $Z_{2\Sigma}+Z_{0\Sigma}$ & $3$ \\
    两相接地短路 & $Z_{2\Sigma}\parallel Z_{0\Sigma}$ &
      $\sqrt{3}\sqrt{1-\dfrac{Z_{2\Sigma}Z_{0\Sigma}}{\left(Z_{2\Sigma}+Z_{0\Sigma}\right)^{2}}}$ \\
    \bottomrule
  \end{tabular}
\end{table}

正序等效定则的价值在于：\textbf{三相短路的全部算法（节点阻抗法、转移阻抗法、计算曲线）
都可以直接用于不对称短路}，只需把正序网的故障点串入 $Z_\Delta$，最后乘以 $m$。

\corepoint{三个必须记牢的结构性事实：（1） 单相接地是\textbf{三序串联}（电流相等、电压和为零），
故 $Z_\Delta$ 为三者之和，故障电流最大；（2） 两相短路\textbf{无零序}，故 $Z_\Delta$ 只含负序；
（3） 两相接地是\textbf{三序并联}，故 $Z_\Delta$ 为负序与零序的并联值，介于前两者之间。
“串/并”这两个字记住了，公式可以现场推出来。}

\subsection{非故障处的电流和电压的计算}

求出故障点各序分量后，网络中任一点 $i$ 的各序电压、电流由序网的电压方程给出：
\begin{align}
  \dot U_{i1}&=\dot U_i(0)-Z_{i f,1}\dot I_1,\\
  \dot U_{i2}&=-Z_{i f,2}\dot I_2,\\
  \dot U_{i0}&=-Z_{i f,0}\dot I_0,
\end{align}
其中 $Z_{if,k}$ 为第 $k$ 序网中 $i$ 与故障点 $f$ 之间的转移阻抗。再把三序分量合成即得
各相电压、电流。工程上关心两件事：
\begin{itemize}
  \item \textbf{非故障相电压升高}：单相接地时非故障相电压可能升到线电压甚至更高
        （中性点不接地系统中可达 $1.7\sim1.9$ 倍相电压），是绝缘配合的依据；
  \item \textbf{零序电流的分布}：决定接地保护的灵敏度与对通信线路的干扰。
\end{itemize}

\section{电压和电流对称分量经变压器后的相位变换}

变压器绕组接线（Y,y；Y,d；D,y；D,d）会使正序、负序分量经过变压器时产生\textbf{相位移动}，
而零序分量则可能根本不能通过。以最常见的 \textbf{Y,d11} 接线为例：
\begin{itemize}
  \item \textbf{正序分量}：由 Y 侧到 $\Delta$ 侧相位\textbf{滞后 $30^\circ$}（由 $\Delta$ 侧到 Y 侧则超前 $30^\circ$）；
  \item \textbf{负序分量}：移相方向与正序\textbf{相反}，即由 Y 侧到 $\Delta$ 侧超前 $30^\circ$；
  \item \textbf{零序分量}：$\Delta$ 侧无零序电流流出通路，零序分量\textbf{不能通过}该变压器。
\end{itemize}
因此，在两侧分别取基准相计算时，须按接线组别引入移相系数
$e^{\pm j30^\circ}$；若全网统一在一侧计算（把另一侧的量折算过来），则移相在折算中自动抵消，
不必另行处理——这是实用计算中常用的简化。

\begin{table}[htbp]
  \centering
  \caption{不同接线变压器的序分量通过情况}
  \small
  \begin{tabular}{@{}llll@{}}
    \toprule
    接线 & 正序 & 负序 & 零序 \\
    \midrule
    Y,y（中性点均接地） & 通，无移相 & 通，无移相 & 通（须两侧均接地） \\
    Y,d11 & 通，移相 $-30^\circ$ & 通，移相 $+30^\circ$ & 不通 \\
    D,y11 & 通，移相 $+30^\circ$ & 通，移相 $-30^\circ$ & 不通 \\
    Y,$\Delta$（Y 侧中性点不接地） & 通 & 通 & 不通 \\
    \bottomrule
  \end{tabular}
\end{table}

\corepoint{相位变换只影响“两侧分别计算再对接”的情形，不影响故障点本身的电流大小。
考试中若题目要求“求 $\Delta$ 侧各相电流”，就必须移相；若只问“故障点电流”，可完全忽略。
判断依据是：所要求的量在变压器的哪一侧。}

\section{非全相断线的分析计算}

非全相断线（一相或两相断开）是\textbf{纵向故障}：故障处不是“相对地短路”，
而是“导体断开”，故其边界条件写在\textbf{断口两端}。

\subsection{单相（$a$ 相）断线}

断口处 $a$ 相电流为零，$b$、$c$ 相仍连通（断口电压为零）：
\begin{equation}
  \dot I_a=0,\qquad \Delta\dot U_b=\Delta\dot U_c=0,
\end{equation}
其中 $\Delta\dot U$ 为断口两侧的电压差。化为序分量：
\begin{equation}
  \dot I_{a1}+\dot I_{a2}+\dot I_{a0}=0,\qquad
  \Delta\dot U_{a1}=\Delta\dot U_{a2}=\Delta\dot U_{a0},
\end{equation}
即\textbf{三序网在断口处并联}——与两相接地短路的序网联接方式相同，只是“端口”从横向
（相对地）变为纵向（串联在支路中）。故可用同一套正序等效定则，只是把
$Z_\Delta$ 换成“断口两端看进去的序阻抗”。

\subsection{两相（$b$ 相和 $c$ 相）断开}

\begin{equation}
  \dot I_b=\dot I_c=0,\qquad \Delta\dot U_a=0,
\end{equation}
化为
\begin{equation}
  \dot I_{a1}=\dot I_{a2}=\dot I_{a0},\qquad
  \Delta\dot U_{a1}+\Delta\dot U_{a2}+\Delta\dot U_{a0}=0,
\end{equation}
即\textbf{三序网串联}——与单相接地短路对应。

\begin{table}[htbp]
  \centering
  \caption{横向故障（短路）与纵向故障（断线）的对应关系}
  \small
  \begin{tabular}{@{}lll@{}}
    \toprule
    短路（横向） & 断线（纵向） & 序网联接 \\
    \midrule
    单相接地 & 两相断开 & 三序串联 \\
    两相短路 & （无对应简单情形） & 正、负序并联 \\
    两相接地 & 单相断线 & 三序并联 \\
    \bottomrule
  \end{tabular}
\end{table}

断线的危害不在于过电流（断线处电流反而减小），而在于：
\begin{enumerate}
  \item 断口两侧电压不对称，负序、零序电流可能引起保护误动；
  \item 单相断线时非故障相电压升高、负荷电流重新分配，可能造成设备过负荷；
  \item 断线常与接地故障同时发生（断线落地），形成复杂故障。
\end{enumerate}

\section{应用节点阻抗矩阵计算不对称故障}

\subsection{各序网络的电压方程式}

设故障点为节点 $f$，三序网的节点阻抗矩阵分别为 $\boldsymbol Z_1$、$\boldsymbol Z_2$、$\boldsymbol Z_0$。
故障点各序电压为
\begin{align}
  \dot U_{f1}&=\dot U_f(0)-Z_{ff1}\dot I_1,\\
  \dot U_{f2}&=-Z_{ff2}\dot I_2,\\
  \dot U_{f0}&=-Z_{ff0}\dot I_0,
\end{align}
其中 $Z_{ffk}$ 为第 $k$ 序网中故障点的自阻抗（输入阻抗）。
把这三式与故障边界条件联立，即得各序电流。对四种横向故障，
其结果可统一写成
\keyformula{\dot I_1=\frac{\dot U_f(0)}{Z_{ff1}+Z_{\Delta}},\qquad
  Z_{\Delta}=\begin{cases}
    0, & \text{三相短路}\\
    Z_{ff2}, & \text{两相短路}\\
    Z_{ff2}+Z_{ff0}, & \text{单相接地}\\
    Z_{ff2}\parallel Z_{ff0}, & \text{两相接地}
  \end{cases}}
这正是正序等效定则在节点阻抗矩阵形式下的表述：$Z_{1\Sigma}$ 换成 $Z_{ff1}$。

\subsection{横向不对称故障}

横向故障（短路）的故障端口在“节点与地”之间：
\begin{itemize}
  \item 先由第 4 章方法求出 $\boldsymbol Z_1$、$\boldsymbol Z_2$、$\boldsymbol Z_0$ 中所需的一列；
  \item 按故障类型取 $Z_\Delta$，由上式求 $\dot I_1$，再由边界条件求 $\dot I_2$、$\dot I_0$；
  \item 由 $\dot U_{ik}=\dot U_i(0)-Z_{ifk}\dot I_k$ 求任一点各序电压，合成即得相电压、相电流。
\end{itemize}

\subsection{纵向不对称故障}

纵向故障（断线）的故障端口在“支路两端”之间。设断线支路连接节点 $p$、$q$，
则断口电压为 $\Delta\dot U=\dot U_p-\dot U_q$，端口阻抗要取“$p$、$q$ 之间的转移阻抗”：
\begin{equation}
  Z_{\Delta}^{(\text{纵})}=Z_{pp}+Z_{qq}-2Z_{pq},
\end{equation}
再按断线类型（单相断线并联、两相断开串联）取用。这与第 4 章“支路追加法”中
连支回路阻抗 $z_\Sigma=Z_{ii}+Z_{jj}-2Z_{ij}+z$ 是同一个式子（令 $z=0$）。

\subsection{简单不对称故障的计算通式}

把四种横向故障与两种纵向故障统一起来，可得计算通式
\begin{equation}
  \dot I_1=\frac{\dot U_f(0)}{Z_{11}+Z_{\Delta}},\qquad
  \dot I_k=c_k\dot I_1\ (k=1,2,0),\qquad
  \dot I_{\text{fault}}=m\dot I_1 ,
\end{equation}
其中系数 $c_k$ 由边界条件决定（如单相接地 $c_1=c_2=c_0=1$；
两相短路 $c_2=-1$、$c_0=0$；两相接地 $c_2=-Z_0/(Z_2+Z_0)$、$c_0=-Z_2/(Z_2+Z_0)$）。
计算通式的好处是：\textbf{程序只需一份代码}，故障类型作为输入参数（$Z_\Delta$ 与 $c_k$）即可。

\section{复杂故障的计算方法}

（选学内容）复杂故障指\textbf{同时或相继发生两个及以上故障}，如一点断线加另一点接地、
双重接地等。分析方法的一般思路：
\begin{enumerate}
  \item 把每个故障都视为一个“端口”，横向故障为并联端口、纵向故障为串联端口；
  \item 对每个端口写出边界条件，化为序分量后得到端口电压、电流的关系；
  \item 联立各端口的序网方程（用节点阻抗矩阵表示），解出各端口的序电流；
  \item 再求全网各点电压、电流。
\end{enumerate}
端口数增多时方程阶数增加，手算困难，故实际用计算机求解。

\textbf{不对称故障的通用边界条件}可写成一个矩阵形式：对第 $j$ 个端口，
\begin{equation}
  \boldsymbol M_j\dot{\boldsymbol U}_{j,(120)}=\boldsymbol N_j\dot{\boldsymbol I}_{j,(120)},
\end{equation}
其中 $\boldsymbol M_j$、$\boldsymbol N_j$ 由故障类型决定。这一写法把“串/并联”统一为矩阵约束，
是编制通用故障计算程序的基础。

\corepoint{本章的复习线索只有两条：
（1） \textbf{边界条件 $\to$ 序网串并联}（单相接地串联、两相短路正负序并联、两相接地三序并联、
单相断线并联、两相断开串联）；
（2） \textbf{正序等效定则}把一切故障折算成“三相短路 $+$ 附加阻抗 $Z_\Delta$”。
把表 \ref{tab:zDelta} 与“串/并”对应关系背熟，剩下都是代数运算。}
  % 第 8 章
% ============================================================
% 第 9 章  电力系统的负荷
% （教材：何仰赞、温增银《电力系统分析》下册 第 9 章）
% 结构对齐 9.1—9.3；例题、习题不录
% ============================================================

\chapter{电力系统的负荷}

前面八章把发电机、变压器、线路都变成了等值电路，剩下的最后一个元件是负荷。
负荷是系统的“目的”，也是分析中最难精确描述的部分：它由成千上万台设备组成，
随生产、作息、气候而变，且对电压、频率都有依赖。
本章给出三个层次的内容：\textbf{负荷是什么}（组成）、\textbf{负荷怎样随时间变}（曲线与指标）、
\textbf{负荷怎样随电压频率变}（特性与模型）。

\section{负荷的组成}

系统中所有电力用户的用电设备所消耗的电功率总和，就是电力系统的负荷，
亦称\textbf{综合用电负荷}——它把不同地区、不同性质的所有用户负荷总加起来得到。

\subsection{用电设备的类型与用户分类}

主要用电设备有：异步电动机、同步电动机、电热装置、整流装置、照明设备。
按用户性质可分为工业负荷、农业负荷、交通运输业负荷和人民生活用电负荷。
在不同性质的用户中，各类设备消耗功率的比重不同；在工业内部，不同行业的比重也不同。

\begin{table}[htbp]
  \centering
  \caption{几个工业部门用电设备比重的统计（单位：\%）}
  \small
  \begin{tabular}{@{}lcccccc@{}}
    \toprule
    用电设备 & 综合性中小工业 & 纺织工业 & 化学工业（化肥、焦化） & 化学工业（电化） & 大型机械加工 & 钢铁工业 \\
    \midrule
    异步电动机 & 79.1 & 99.8 & 56.0 & 13.0 & 82.5 & 20.0 \\
    同步电动机 & 3.2 & — & 44.0 & — & 1.3 & 10.0 \\
    电热装置 & 17.7 & 0.2 & — & — & 15.0 & 70.0 \\
    整流装置 & — & — & — & 87.0 & 1.2 & — \\
    \midrule
    合计 & 100.0 & 100.0 & 100.0 & 100.0 & 100.0 & 100.0 \\
    \bottomrule
  \end{tabular}
\end{table}

这张表说明一件重要的事：\textbf{钢铁、电化等行业的负荷以电热、整流为主，功率因数高、
对电压不敏感；纺织、机械行业以异步电动机为主，功率因数低、无功占比大}。
负荷构成决定了无功需求，也决定了第 12 章调压措施的取舍。

\subsection{负荷的功率因数与无功构成}

异步电动机在总负荷中占比大（多数行业 $50\%\sim90\%$），其功率因数随负载率变化剧烈：
\begin{itemize}
  \item 额定负载时 $\cos\varphi\approx0.85\sim0.90$；
  \item 半载时 $\cos\varphi\approx0.7\sim0.8$；
  \item 空载时 $\cos\varphi$ 只有 $0.2\sim0.3$（几乎全是励磁无功）。
\end{itemize}
因此“大马拉小车”（电动机长期轻载）是系统无功需求偏大、功率因数偏低的重要原因。
这也是无功补偿（第 12 章）与“电动机经济运行”的直接依据。

系统的自然功率因数一般在 $0.7\sim0.85$。若不做补偿，无功负荷将大量占用发电机容量与
输电能力：输送同样的有功，功率因数从 $0.85$ 降到 $0.7$，电流增大约 $21\%$，
线路损耗增大约 $47\%$。这就是“无功就地平衡”原则的经济含义。

\subsection{三级负荷概念}

\begin{table}[htbp]
  \centering
  \caption{综合用电负荷、供电负荷与发电负荷}
  \small
  \begin{tabular}{@{}p{0.20\linewidth}p{0.34\linewidth}p{0.34\linewidth}@{}}
    \toprule
    名称 & 定义 & 关系 \\
    \midrule
    综合用电负荷 & 所有用户用电设备消耗的功率之和 & 基准量 \\
    供电负荷 & 综合用电负荷 $+$ 电力网的功率损耗 & 各发电厂应供给的功率 \\
    发电负荷 & 供电负荷 $+$ 发电厂厂用电 & 各发电厂应发出的功率 \\
    \bottomrule
  \end{tabular}
\end{table}

三者的差额就是“网损”与“厂用电”，它们正是第 14 章经济运行的考核对象
（网损率、厂用电率）。

\section{负荷曲线}

负荷随时间变化的规律用\textbf{负荷曲线}描述。常用的有日负荷曲线和年负荷曲线。

\subsection{日负荷曲线}

日负荷曲线描述一天 24 小时负荷的变化。曲线最大值称\textbf{日最大负荷} $P_{\max}$（峰荷），
最小值称\textbf{日最小负荷} $P_{\min}$（谷荷）。为便于计算，常把连续曲线画成阶梯形。

由日负荷曲线可算一日总耗电量与日平均负荷：
\keyformula{W_d=\int_{0}^{24}P\,\mathrm{d}t,\qquad P_{av}=\frac{W_d}{24}}
衡量曲线起伏特性的两个系数：
\begin{align}
  \text{负荷率}\quad k_m&=\frac{P_{av}}{P_{\max}},\\
  \text{最小负荷系数}\quad \alpha&=\frac{P_{\min}}{P_{\max}}.
\end{align}
峰谷差 $\Delta P_m=P_{\max}-P_{\min}$ 是调峰容量的依据。
这两个系数不仅用于日负荷曲线，也用于其他时段的负荷曲线。

不同性质用户的日负荷曲线形状差别很大：
\begin{itemize}
  \item 三班制连续生产的重工业（如钢铁）：曲线平坦，$\alpha$ 可达 $0.85$；
  \item 一班制生产的轻工业（如食品）：变化幅度大，$\alpha$ 只有 $0.13$ 左右；
  \item 非排灌季节的农业：农村加工用电每天仅约 12 小时；
  \item 市政生活：有明显照明高峰。
\end{itemize}
由于各用户的最大、最小负荷不在同一时刻出现，
\textbf{系统的最大负荷总小于各用户最大负荷之和，系统的最小负荷总大于各用户最小负荷之和}
——这就是“负荷同时率”现象的来源，也是系统调峰比单用户调峰更容易的原因。

日负荷曲线是安排\textbf{日发电计划}和确定\textbf{系统运行方式}的重要依据。

\subsection{峰谷差与调峰}

日负荷曲线最受关注的特征是\textbf{峰谷差}
\begin{equation}
  \Delta P_m=P_{\max}-P_{\min}=P_{\max}\left(1-\alpha\right).
\end{equation}
峰谷差决定了系统需要多少“可调节容量”：峰荷时要有足够的机组顶峰，
谷荷时要有足够的机组压负荷（或停机）。若调节能力不足，只能靠“开停机”应对，
既不经济也影响设备寿命。

\begin{table}[htbp]
  \centering
  \caption{典型用户的负荷曲线特征}
  \small
  \begin{tabular}{@{}p{0.24\linewidth}p{0.18\linewidth}p{0.46\linewidth}@{}}
    \toprule
    用户类型 & 最小负荷系数 $\alpha$ & 曲线特征 \\
    \midrule
    三班制重工业（钢铁） & 约 $0.85$ & 平坦，基荷型，宜带基荷 \\
    一班制轻工业（食品） & 约 $0.13$ & 峰谷差大，昼开夜停 \\
    农业（非排灌季节） & — & 每天约 12 小时用电，季节性明显 \\
    市政生活 & — & 有明显照明晚高峰 \\
    \bottomrule
  \end{tabular}
\end{table}

系统峰谷差与用户构成有关：工业比重大则曲线平坦，生活用电比重大则峰谷差大。
随着新能源（光伏）大规模接入，\textbf{“鸭子曲线”}现象（午间光伏出力大、
净负荷下凹，傍晚光伏退出、负荷上升）使峰谷差进一步扩大，对调峰能力提出更高要求。
这也是第 13 章“调频与调峰”、第 14 章“经济运行”要解决的问题的现实背景。

\subsection{年最大负荷曲线与年持续负荷曲线}

\paragraph{年最大负荷曲线}
描述一年内每月（或每日）最大有功功率负荷的变化，主要用来安排\textbf{发电设备的检修计划}，
也为制订发电机组或发电厂的扩建、新建计划提供依据。

\paragraph{年持续负荷曲线}
按一年中系统负荷的数值大小及其持续小时数顺序排列绘制。例如全年 8760 小时中，
有 $t_1$ 小时负荷为 $P_1$（$=P_{\max}$），$t_2$ 小时为 $P_2$，……，于是得到阶梯下降的曲线。
由它可以算出全年耗电量
\begin{equation}
  W=\int_{0}^{8760}P\,\mathrm{d}t .
\end{equation}
若负荷始终等于 $P_{\max}$，经过 $T_{\max}$ 小时所消耗的电能恰好等于全年实际耗电量，则
\keyformula{T_{\max}=\frac{W}{P_{\max}}=\frac{1}{P_{\max}}\int_{0}^{8760}P\,\mathrm{d}t}
$T_{\max}$ 称\textbf{最大负荷利用小时数}。它是衡量负荷特性的综合性指标：
$T_{\max}$ 越大，说明负荷越平稳、设备利用率越高。年持续负荷曲线常用于安排发电计划
和进行可靠性估算。

\begin{table}[htbp]
  \centering
  \caption{负荷曲线的种类与用途}
  \small
  \begin{tabular}{@{}p{0.22\linewidth}p{0.24\linewidth}p{0.42\linewidth}@{}}
    \toprule
    曲线 & 描述对象 & 主要用途 \\
    \midrule
    日负荷曲线 & 一天 24 小时 & 安排日发电计划、确定运行方式、调峰 \\
    年最大负荷曲线 & 每月（日）最大值 & 安排检修计划、规划扩建 \\
    年持续负荷曲线 & 负荷大小与持续小时数 & 年耗电量、$T_{\max}$、可靠性估算 \\
    \bottomrule
  \end{tabular}
\end{table}

\corepoint{三个指标不要混：\textbf{负荷率} $k_m=P_{av}/P_{\max}$ 反映“平均用得多满”；
\textbf{最小负荷系数} $\alpha=P_{\min}/P_{\max}$ 反映“低谷有多低”（调峰难度）；
\textbf{最大负荷利用小时数} $T_{\max}=W/P_{\max}$ 反映“全年等效满负荷运行多少小时”。
$k_m$ 与 $T_{\max}$ 的关系是 $T_{\max}=8760\,k_m$。}

\section{负荷特性与负荷模型}

\subsection{负荷的静态特性}

负荷特性指负荷功率随电压、频率变化而变化的规律。
\textbf{静态特性}指电压、频率缓慢变化时的规律；\textbf{动态特性}指快速变化（如短路、电动机
失速）时的规律。本课程主要用静态特性。

\paragraph{电压静态特性（二次多项式，ZIP 模型）}
\begin{equation}
  P=P_N\left[a_p\left(\frac{U}{U_N}\right)^{2}+b_p\left(\frac{U}{U_N}\right)+c_p\right],
\end{equation}
\begin{equation}
  Q=Q_N\left[a_q\left(\frac{U}{U_N}\right)^{2}+b_q\left(\frac{U}{U_N}\right)+c_q\right],
\end{equation}
系数由实测特性用最小二乘法拟合，且满足
\begin{equation}
  a_p+b_p+c_p=1,\qquad a_q+b_q+c_q=1 .
\end{equation}
三项各有明确物理含义：
\begin{itemize}
  \item 与 $U^{2}$ 成正比的第一项：代表\textbf{恒定阻抗}消耗的功率；
  \item 与 $U$ 成正比的第二项：代表\textbf{恒定电流}负荷对应的功率；
  \item 常数项：代表\textbf{恒定功率}分量。
\end{itemize}
故该模型又称 \textbf{ZIP}（恒阻抗—恒电流—恒功率）模型。

\paragraph{线性化模型}
当电压、频率在额定值附近小幅度变化时，可线性化：
\begin{align}
  P&=P_N\left(1+k_{pv}\delta U\right),&
  Q&=Q_N\left(1+k_{qv}\delta U\right),\\
  P&=P_N\left(1+k_{pf}\delta f\right),&
  Q&=Q_N\left(1+k_{qf}\delta f\right),
\end{align}
其中
\begin{equation}
  \delta U=\frac{U-U_N}{U_N},\qquad \delta f=\frac{f-f_N}{f_N}.
\end{equation}
$k_{pv}$、$k_{qv}$ 称\textbf{负荷的电压调节效应系数}，$k_{pf}$、$k_{qf}$ 称\textbf{频率调节效应系数}。
同时考虑电压与频率变化时可用
\begin{align}
  P&=P_N\left(1+k_{pv}\delta U\right)\left(1+k_{pf}\delta f\right),\\
  Q&=Q_N\left(1+k_{qv}\delta U\right)\left(1+k_{qf}\delta f\right).
\end{align}
这些系数是第 12 章（电压与无功）和第 13 章（频率与有功）中“负荷自身调节作用”的量化依据。

\subsection{负荷的动态模型}

反映动态特性的模型一般由微分方程与代数方程组成，随所研究问题不同而取不同形式。
建立综合负荷动态模型包含两件事：\textbf{制订模型结构}与\textbf{确定模型参数}。
由于综合负荷所代表的设备数量极大、分布极广、种类繁多，工作状态带有随机性与时变性
（甚至跃变性），配电网结构也可能变化，因此“既准确又实用”的负荷模型至今仍未很好解决。

\subsection{综合负荷的等值电路}

电力系统分析中元件都用等值电路代表，负荷也不例外。最常采用的综合负荷等值电路有：
\begin{itemize}
  \item \textbf{含源等值阻抗（或导纳）支路}；
  \item \textbf{恒定阻抗（或导纳）支路}；
  \item \textbf{异步电动机等值电路}（阻抗值随转差而变）；
  \item 上述电路的不同组合。
\end{itemize}

\begin{table}[htbp]
  \centering
  \caption{不同计算中负荷模型的取用}
  \small
  \begin{tabular}{@{}p{0.18\linewidth}p{0.34\linewidth}p{0.38\linewidth}@{}}
    \toprule
    计算类型 & 负荷的处理 & 理由 \\
    \midrule
    潮流计算 & 恒定功率 $S_{LD}=P_{LD}+jQ_{LD}$（必要时用线性化静态特性） & 负荷功率由用户决定，与网络电压弱相关 \\
    短路计算 & 含源阻抗支路，或恒定阻抗支路 & 短路瞬间电动机可倒送电流（第 6 章） \\
    稳定计算 & 恒定阻抗，或恒定阻抗与异步电动机的组合 & 需反映电压、频率变化时的功率响应 \\
    \bottomrule
  \end{tabular}
\end{table}

采用“恒定阻抗 $+$ 异步电动机”组合模型时，还须补充一个反映异步电动机机械运动状态的
\textbf{转子运动方程}（第 17 章），否则无法描述电动机失速与恢复的过程。

\corepoint{同一套系统，在潮流、短路、稳定三类计算中负荷模型不同，原因不是“简化程度不同”，
而是\textbf{所关心的物理过程不同}：潮流关心功率分布（负荷按功率给定）、
短路关心瞬间电流（负荷可视为阻抗或电源）、稳定关心机电过程（负荷须有动态）。
选模型前先问“我要算的是哪个时间尺度上的什么量”。}
  % 第 9 章
% ============================================================
% 第 10 章  电力传输的基本概念
% （教材：何仰赞、温增银《电力系统分析》下册 第 10 章）
% 结构对齐 10.1—10.4；例题、习题不录
% ============================================================

\chapter{电力传输的基本概念}

本章从“功率怎样在网络上传输”这个角度重新看元件，得到三个贯穿全书的结论：
\textbf{电压幅值差由无功决定、电压相角差由有功决定}（10.1）；
\textbf{输电线路存在功率极限，极限由电压与电抗之比决定}（10.2、10.4）；
\textbf{长线的电压分布由波过程决定，自然功率是“理想运行点”}（10.3）。
第 11～14 章的潮流、调压、调频、经济运行，全部建立在这几个概念上。

\section{网络元件的电压降落和功率损耗}

\subsection{网络元件的电压降落}

设元件一相等值电路为 $Z=R+jX$，首末端相电压 $\dot U_1$、$\dot U_2$，电流 $\dot I$。则
\begin{equation}
  \dot U_1-\dot U_2=(R+jX)\dot I .
\end{equation}
以 $\dot U_2$ 为参考轴，把电压降相量 $(R+jX)\dot I$ 分解为与 $\dot U_2$ 同方向的分量
$\Delta\dot U_2$（\textbf{纵分量}）与垂直的分量 $\delta\dot U_2$（\textbf{横分量}），得
\begin{align}
  \Delta U_2&=RI\cos\varphi_2+XI\sin\varphi_2,\\
  \delta U_2&=XI\cos\varphi_2-RI\sin\varphi_2 .
\end{align}
电力网分析习惯用功率运算。以末端一相功率 $S_2=P_2+jQ_2=U_2I^{*}$ 代替电流，得
\keyformula{\Delta U_2=\frac{P_2R+Q_2X}{U_2},\qquad
  \delta U_2=\frac{P_2X-Q_2R}{U_2}}
于是首端电压
\begin{equation}
  U_1=\sqrt{\left(U_2+\Delta U_2\right)^{2}+\left(\delta U_2\right)^{2}},
  \qquad \delta=\arctan\frac{\delta U_2}{U_2+\Delta U_2}.
\end{equation}
同理，若以首端电压 $\dot U_1$ 为参考、用首端功率 $S_1=P_1+jQ_1$，则
\keyformula{\Delta U_1=\frac{P_1R+Q_1X}{U_1},\qquad
  \delta U_1=\frac{P_1X-Q_1R}{U_1}}
\begin{equation}
  U_2=\sqrt{\left(U_1-\Delta U_1\right)^{2}+\left(\delta U_1\right)^{2}} .
\end{equation}

\textbf{必须注意两点}：
\begin{enumerate}
  \item 用某一点的功率，就必须用\textbf{同一点}的电压——$\Delta U_2$ 用 $U_2$、$\Delta U_1$ 用 $U_1$，
        不能混用；
  \item 上述公式按“电流落后于电压（功率因数角 $\varphi$ 为正）”导出；若电流超前，
        $Q$ 应取负值。本书约定：\textbf{感性无功为正，容性无功为负}。
\end{enumerate}
显然 $\Delta U_1\ne\Delta U_2$、$\delta U_1\ne\delta U_2$——两种分解的结果不同，
只是同一相量的不同投影，不可混用。

\subsection{电压损耗与电压偏移}

\paragraph{电压损耗} 两点间电压\textbf{绝对值}之差：
\begin{equation}
  \Delta U=U_1-U_2 .
\end{equation}
当两端电压相角差 $\delta$ 不大时，$\Delta U$ 与纵分量 $\Delta U_2$ 相差不大，
可近似认为\textbf{电压损耗就等于电压降落的纵分量}。工程上常用额定电压的百分数表示。
从电源点到某负荷点的总电压损耗，等于沿途各串联元件电压损耗的代数和。

\paragraph{电压偏移} 网络中某点实际电压与该处额定电压之差：
\begin{equation}
  \text{电压偏移}(\%)=\frac{U-U_N}{U_N}\times 100\% .
\end{equation}
电压损耗是“元件两端之差”（过程量），电压偏移是“某点对额定值之差”（状态量），
二者是衡量电压质量的两个不同概念；实际电压的高低对用户有影响，而电压的相位对用户没有影响。

\subsection{电压降落公式的分析}

由 $\Delta U=(PR+QX)/U$、$\delta U=(PX-QR)/U$ 可见：
\begin{itemize}
  \item 元件两端电压的\textbf{幅值差}主要由纵分量决定，\textbf{相角差}由横分量决定；
  \item 高压输电线的电抗远大于电阻，取极端 $R=0$，则
        \begin{equation}
          \Delta U=\frac{QX}{U},\qquad \delta U=\frac{PX}{U}.
        \end{equation}
\end{itemize}
\corepoint{纯电抗元件中：电压降落的\textbf{纵分量由传送无功产生}，\textbf{横分量由传送有功产生}。
换句话说——\textbf{元件两端存在电压幅值差是传送无功的条件，存在电压相角差是传送有功的条件}。
感性无功从电压高的一端流向电压低的一端；有功从相位超前的一端流向相位落后的一端。
这两句话是第 11 章（潮流）和第 15 章（功角）的共同起点。}

实际元件都有电阻：电流的有功分量流过电阻会增大纵分量；
电流的感性无功分量流过电阻会使横分量略有减小。

\subsection{网络元件的功率损耗}

\paragraph{串联阻抗上的功率损耗}
\begin{equation}
  \Delta S_L=\Delta P_L+j\Delta Q_L=I^{2}(R+jX)
  =\frac{P^{2}+Q^{2}}{U^{2}}(R+jX),
\end{equation}
其中 $P$、$Q$、$U$ 取同一端的值。可见功率损耗与\textbf{电流平方}（即功率平方）成正比、
与\textbf{电压平方}成反比——这是“提高电压等级可降低网损”的定量依据。

\paragraph{并联导纳上的功率损耗}
线路电纳产生容性无功（按本书约定 $Q$ 取负号）：
\begin{equation}
  \Delta Q_{B1}=-\tfrac{1}{2}BU_1^{2},\qquad \Delta Q_{B2}=-\tfrac{1}{2}BU_2^{2}.
\end{equation}
变压器励磁支路的损耗可直接用空载试验数据：
\begin{equation}
  \Delta S_0=\Delta P_0+jQ_0=\Delta P_0+j\frac{I_0\%}{100}S_N ,
\end{equation}
一般不考虑电压变化对它的影响。35~kV 以下电网的简化计算中常略去变压器励磁功率。

\paragraph{输电效率}
\begin{equation}
  \text{输电效率}=\frac{P_2}{P_1}\times 100\% .
\end{equation}

以上公式由单相电路导出，实际计算中改用电网习惯的\textbf{线电压与三相功率}时形式不变，
只需把单位统一为：阻抗 $\Omega$、导纳 S、电压 kV、功率 MV$\cdot$A。

\section{输电线路的功率特性}

\subsection{两端口方程}

线路两端电压电流满足
\begin{equation}
  \dot U_1=A\dot U_2+B\dot I_2,\qquad \dot I_1=C\dot U_2+D\dot I_2,
  \qquad AD-BC=1 .
\end{equation}
由之可解出用 $\dot U_1$、$\dot U_2$ 表示的两端功率。

\subsection{功率特性与功率极限}

对“发电机经电抗 $X$ 供一负荷”的最简单系统（略电阻），受端功率为
\keyformula{P_2=\frac{U_1U_2}{X}\sin\delta,\qquad
  Q_2=\frac{U_1U_2\cos\delta-U_2^{2}}{X}}
其中 $\delta$ 为首末端电压的相角差。这两式是全书最重要的两个式子：
\begin{itemize}
  \item $P_2$ 随 $\delta$ 按正弦变化，$\delta=90^\circ$ 时达\textbf{功率极限}
        \begin{equation}
          P_{\max}=\frac{U_1U_2}{X};
        \end{equation}
  \item $Q_2$ 在 $\delta=0$ 时最大（$=U_2(U_1-U_2)/X$），随 $\delta$ 增大而减小，
        甚至变负（受端发出无功）。
\end{itemize}

\begin{table}[htbp]
  \centering
  \caption{功率传输的三个基本结论}
  \small
  \begin{tabular}{@{}p{0.26\linewidth}p{0.60\linewidth}@{}}
    \toprule
    结论 & 说明 \\
    \midrule
    $P\propto\sin\delta$，极限 $P_{\max}=U_1U_2/X$ & 有功传输靠相角差；要增大传输能力，
      只能提高电压或减小电抗（串补、分裂导线、增加回路数） \\
    $Q$ 由电压幅值差决定 & 无功从电压高处流向低处；要维持受端电压，须就地补偿无功 \\
    $P$、$Q$ 相互耦合 & 改变 $\delta$ 会改变 $Q$，改变 $U$ 会改变 $P_{\max}$——
      这正是潮流计算必须联立求解的原因 \\
    \bottomrule
  \end{tabular}
\end{table}

\subsection{输送功率与电压水平的关系}

由 $P_2=U_1U_2\sin\delta/X$ 与 $Q_2=(U_1U_2\cos\delta-U_2^{2})/X$ 消去 $\delta$，可得
\begin{equation}
  \left(\frac{P_2X}{U_1U_2}\right)^{2}
  =\sin^{2}\delta
  =1-\left(\frac{Q_2X+U_2^{2}}{U_1U_2}\right)^{2},
\end{equation}
整理即得 $P$-$Q$ 平面上的\textbf{圆方程}：受端功率的轨迹是一族以
$(0,\,-U_2^{2}/X)$ 为圆心、以 $U_1U_2/X$ 为半径的圆。
其工程含义是：
\begin{itemize}
  \item 输送有功越多，受端能维持的电压越低（同一 $U_1$ 下）；
  \item 若要保持受端电压不变，输送功率增大时必须相应增加无功补偿（降低 $Q$，即“发无功”）；
  \item 当 $\delta\to90^\circ$ 时，$P_2\to P_{\max}$、$Q_2\to-U_2^{2}/X$（受端大量发出无功），
        这是功率极限点，也是静态稳定极限。
\end{itemize}
这一结论把“电压”“无功”“有功”“功角”四者放在同一张图上，
是第 12 章调压与第 18 章静态稳定的共同基础。

\section{沿长线的功率传送}

\subsection{长线的稳态方程}

长线方程（第 2 章）为
\begin{equation}
  \dot U=\dot U_2\operatorname{ch}\gamma x+\dot I_2Z_c\operatorname{sh}\gamma x,\qquad
  \dot I=\frac{\dot U_2}{Z_c}\operatorname{sh}\gamma x+\dot I_2\operatorname{ch}\gamma x .
\end{equation}
其中传播常数 $\gamma=\beta+j\alpha$、波阻抗 $Z_c=\sqrt{z_0/y_0}$。

\subsection{输电线路的自然功率}

对\textbf{无损线}（$r_0=g_0=0$），$\gamma=j\alpha=j\omega\sqrt{L_0C_0}$，
$Z_c=R_c=\sqrt{L_0/C_0}$ 为纯电阻。若线路末端所接负荷恰好等于波阻抗
\begin{equation}
  Z_{LD}=Z_c ,
\end{equation}
则由长线方程可得沿线各点
\begin{equation}
  \dot U=\dot U_2e^{j\alpha x},\qquad \dot I=\dot I_2e^{j\alpha x},
\end{equation}
即\textbf{电压、电流的幅值沿线不变，只有相位随距离线性滞后}。
此时线路输送的功率称\textbf{自然功率}
\keyformula{S_{nat}=\frac{U^{2}}{Z_c}\quad\Longrightarrow\quad P_{nat}=\frac{U_N^{2}}{Z_c}}
自然功率是长线的“理想运行点”：
\begin{itemize}
  \item 输送功率 $=$ 自然功率时，线路本身既不吸收也不发出无功，沿线电压平坦；
  \item 输送功率 $<$ 自然功率时，线路呈\textbf{容性}（充电功率大于串联电抗消耗的无功），
        末端电压升高；
  \item 输送功率 $>$ 自然功率时，线路呈\textbf{感性}，末端电压降低，需加串补或并联电容。
\end{itemize}
波阻抗：单导线架空线 $370\sim410\ \Omega$，分裂导线 $270\sim310\ \Omega$，电缆 $30\sim50\ \Omega$。
故 $500$~kV 线路的自然功率约 $900\sim1000$~MW，$220$~kV 约 $120\sim150$~MW，
电缆的自然功率则大得多（因 $Z_c$ 小）。

\begin{table}[htbp]
  \centering
  \caption{各级电压线路的自然功率（数量级）}
  \small
  \begin{tabular}{@{}lccc@{}}
    \toprule
    额定电压/kV & 波阻抗/$\Omega$（分裂导线） & 自然功率/MW & 说明 \\
    \midrule
    110 & $350\sim400$ & 约 $30\sim35$ & 单导线居多 \\
    220 & $300\sim350$ & 约 $120\sim150$ & 单导线或双分裂 \\
    330 & $280\sim310$ & 约 $350\sim390$ & 双分裂 \\
    500 & $250\sim280$ & 约 $900\sim1000$ & 四分裂 \\
    750 & $240\sim260$ & 约 $2100\sim2300$ & 六分裂 \\
    \bottomrule
  \end{tabular}
\end{table}

这张表给出“一条线能送多少”的数量级依据：若实际输送功率远小于自然功率，
线路呈容性、末端电压升高，须投并联电抗器；若远大于自然功率，线路呈感性、
末端电压降低，须加串补或并联电容。超高压线路的设计输送功率一般取自然功率的
$1.0\sim1.5$ 倍。

\subsection{无损线的功率圆图}

固定首端电压 $U_1$ 时，受端功率 $S_2=P_2+jQ_2$ 随 $\dot U_2$ 变化，
在 $P$-$Q$ 平面上构成一族圆（\textbf{功率圆图}）。圆心与半径由 $U_1$、$Z_c$、$\alpha l$ 决定。
圆图的价值在于：把“电压约束”与“功率约束”放在同一张图上，
可直观看出\textbf{输送功率越大，受端电压越低；要维持电压，必须减少输送功率或补偿无功}。

\subsection{沿长线的电压和电流分布}

\paragraph{无损线}
由 $\dot U=\dot U_2\cos\alpha x+jZ_c\dot I_2\sin\alpha x$ 可见，沿线电压、电流为
\textbf{驻波}形式，波长
\begin{equation}
  \lambda=\frac{v}{f}\approx\frac{3\times10^{5}\ \text{km/s}}{50\ \text{Hz}}=6000\ \text{km}.
\end{equation}
\begin{itemize}
  \item \textbf{空载}（$\dot I_2=0$）：$\dot U=\dot U_2\cos\alpha x$，
        末端电压最高，$\dot U_1=\dot U_2\cos\alpha l$，
        即 $U_2=U_1/\cos\alpha l>U_1$——这就是长线的\textbf{电容效应}（空载末端电压升高），
        超高压长线必须用并联电抗器补偿。
  \item \textbf{末端短路}（$\dot U_2=0$）：$\dot U=jZ_c\dot I_2\sin\alpha x$，
        首端电流由线路波阻抗限制，不会无限增大。
\end{itemize}

\paragraph{有损耗线路}
$r_0\ne0$、$g_0\ne0$ 时 $\gamma=\beta+j\alpha$ 有衰减，电压、电流幅值沿线按 $e^{-\beta x}$ 衰减，
驻波不再是等幅的；末端电压升高的程度比无损线小。工程上仍可用“自然功率”概念作定性判断。

\subsection{单端供电系统的功率特性}

对“电源经阻抗 $Z=R+jX$ 供一负荷、受端电压 $U_2$ 恒定”的单端供电系统，
受端功率特性与 10.2 节同形（略电阻时 $P_2=U_1U_2\sin\delta/X$），
但工程上更关心\textbf{首端}（电源侧）的功率与电压关系：
\begin{equation}
  P_1=\frac{U_1^{2}}{Z^{2}}R+\frac{U_1U_2}{Z^{2}}\left(R\cos\delta+X\sin\delta\right),
\end{equation}
其最大值即\textbf{传输功率极限}，出现在 $\delta$ 略小于 $90^\circ$ 处（计及电阻时）。

\begin{table}[htbp]
  \centering
  \caption{本章公式的物理含义一览}
  \small
  \begin{tabular}{@{}p{0.24\linewidth}p{0.30\linewidth}p{0.36\linewidth}@{}}
    \toprule
    公式 & 说明什么 & 用在哪 \\
    \midrule
    $\Delta U=(PR+QX)/U$ & 幅值差由无功主导 & 电压调整（12）、潮流（11） \\
    $\delta U=(PX-QR)/U$ & 相角差由有功主导 & 功角（15、16）、稳定 \\
    $P=U_1U_2\sin\delta/X$ & 有功传输与极限 & 静态稳定（18）、暂态稳定（17） \\
    $Z_c=\sqrt{L_0/C_0}$ & 长线特性阻抗 & 自然功率、过电压 \\
    $P_{nat}=U^{2}/Z_c$ & 长线理想输送能力 & 超高压输电规划 \\
    \bottomrule
  \end{tabular}
\end{table}

\corepoint{本章要建立的是“数量级感”而不是推导能力：一条 $500$~kV 线路的
自然功率约 $1000$~MW、波阻抗约 $300\ \Omega$、波长 $6000$~km；
一条线路的电压降落的纵分量约为 $(QX)/U$，横分量约为 $(PX)/U$。
有了这些数量级，做题时一眼就能看出结果是否离谱，这是教材反复强调的“相对数值概念”。}
  % 第 10 章
% ============================================================
% 第 11 章  电力系统的潮流计算
% （教材：何仰赞、温增银《电力系统分析》下册 第 11 章）
% 结构对齐 11.1—11.5；例题、习题不录
% ============================================================

\chapter{电力系统的潮流计算}

潮流计算要回答：\textbf{给定网络结构与部分节点的功率、电压，求全网各节点的电压和各支路的功率分布}。
它是稳态分析的核心，也是短路、稳定、经济调度、调压计算的公共基础。
手算路线（开式网、两端供电网）给出物理图像，计算机路线（牛顿—拉夫逊、PQ 分解）给出工程算法，
两条线在本章并行展开。

\section{开式网络的电压和功率分布计算}

\subsection{已知供电点电压和负荷节点功率时的计算方法}

开式网（放射式、干线式）中每个负荷只有唯一供电路径，故功率方向确定，
可逐段推算。给定：首端（供电点）电压 $U_A$、各负荷节点功率 $S_i$；求各节点电压与各段功率。

\textbf{计算步骤（三步一迭代）}：
\begin{enumerate}
  \item \textbf{设全网电压为额定值} $U_i=U_N$（第一次近似）。
        因为功率损耗与电压平方成反比，而电压偏差通常不超过 $\pm10\%$，故这一步误差不大。
  \item \textbf{由末端向首端}逐段计算功率损耗与功率分布：
        \begin{equation}
          \Delta S_i=\frac{P_i^{2}+Q_i^{2}}{U_i^{2}}\left(R_i+jX_i\right),
          \qquad S_{i-1}=S_i+\Delta S_i+\Delta S_{Ti},
        \end{equation}
        其中 $\Delta S_{Ti}$ 为该段变压器绕组的功率损耗（若有）。
  \item \textbf{由首端向末端}用已知的首端电压与上一步求得的功率，逐段计算电压降落，
        得各节点电压：
        \begin{equation}
          \Delta U_i=\frac{P_iR_i+Q_iX_i}{U_i},\qquad
          \delta U_i=\frac{P_iX_i-Q_iR_i}{U_i},
        \end{equation}
        $U_{i-1}=\sqrt{(U_i+\Delta U_i)^{2}+(\delta U_i)^{2}}$。
  \item \textbf{校验}：若算得的节点电压与第 1 步假设值相差较大（一般 $>5\%$），
        用新电压回到第 2 步重算。通常迭代 $1\sim2$ 次即满足工程精度。
\end{enumerate}

\begin{table}[htbp]
  \centering
  \caption{开式网手算的两个方向}
  \small
  \begin{tabular}{@{}p{0.18\linewidth}p{0.28\linewidth}p{0.44\linewidth}@{}}
    \toprule
    方向 & 算什么 & 用什么电压 \\
    \midrule
    末端 $\to$ 首端 & 功率损耗、各段功率 & 假设（或上次迭代）的节点电压 \\
    首端 $\to$ 末端 & 电压降落、各节点电压 & 已知首端电压与该段功率 \\
    \bottomrule
  \end{tabular}
\end{table}

\corepoint{“先算功率、后算电压”是手算开式网的基本次序，原因是\textbf{功率损耗只依赖电压的近似值，
而电压降落依赖已求得的功率}。若颠倒次序，就要解非线性方程，失去手算的可能。
这一“功率单向推算 $+$ 电压单向推算”的交替，正是牛顿法“解耦”思想的手算原型。}

\subsection{两级电压的开式电力网计算}

当网络含变压器、且首末两端处于不同电压等级时，有两种处理方式：
\begin{enumerate}
  \item \textbf{归算法}：把变压器二次侧各元件阻抗、电压、功率全部折算到一次侧
        （乘 $k^{2}$、$k$），化成单一电压级网络，算出结果再折算回原侧。
        折算后变压器只表现为串联阻抗 $R_T+jX_T$，处理最简单。
  \item \textbf{分别计算法}：在各自电压级内计算，把变压器的电压降落单独处理。
        此时须注意变压器的\textbf{实际变比}（含分接头）对两侧电压的约束：
        \begin{equation}
          U_2=\frac{U_1-\Delta U_T}{k},
        \end{equation}
        其中 $\Delta U_T=\left(P R_T+Q X_T\right)/U_1$ 为归算到一次侧的变压器电压降落。
\end{enumerate}
实用中多用归算法，因其把“变比”这一非线性环节消去了。
若用标幺制并以平均额定电压为基准，变压器标幺变比为 $1$，归算自动完成（第 2 章）。

\section{简单闭式网络的功率分布计算}

\subsection{两端供电网络的功率分布}

两端供电网（含环网）中功率方向不预先知道，不能像开式网那样逐段推算。
基本结论是：\textbf{功率按各支路共轭阻抗的反比分配}，再加上由两端电压不等引起的\textbf{循环功率}。

设两端电源 $A$、$B$，负荷 $S_i$ 接于节点 $i$，支路阻抗 $Z_{AB}$（$A$ 到 $B$ 的总阻抗），
$Z_{iB}$ 为节点 $i$ 到 $B$ 的阻抗，则 $A$ 端送出的功率为
\keyformula{\dot S_A=\frac{\sum_i\dot S_i\hat Z_{iB}}{\hat Z_{AB}}
  +\frac{\left(\dot U_A-\dot U_B\right)U_N}{\sqrt{3}\,\hat Z_{AB}}}
其中 $\hat Z$ 表示共轭；第二项即\textbf{循环功率}，当 $\dot U_A=\dot U_B$ 时为零。
$B$ 端功率由功率平衡求出：$\dot S_B=\sum_i\dot S_i-\dot S_A$。

\textbf{均一网}（各段 $X/R$ 相同，如全线同型号导线）时，共轭阻抗之比等于长度之比：
\begin{equation}
  \dot S_A=\frac{\sum_i\dot S_iL_{iB}}{L_{AB}}+S_{cir},
\end{equation}
即功率按\textbf{长度反比}分配，手算极为方便。

\begin{table}[htbp]
  \centering
  \caption{闭式网功率分布的两部分}
  \small
  \begin{tabular}{@{}p{0.20\linewidth}p{0.34\linewidth}p{0.36\linewidth}@{}}
    \toprule
    分量 & 来源 & 特点 \\
    \midrule
    供载功率 & 各负荷功率按共轭阻抗反比分配 & 与负荷成正比，方向随负荷而变 \\
    循环功率 & 两端电压不等（或变比不匹配） & 与负荷无关，环流不衰减，增大网损 \\
    \bottomrule
  \end{tabular}
\end{table}

\subsection{闭式电力网中的电压损耗计算}

闭式网算出功率分布后，须先确定\textbf{功率分点}：功率由两侧流入该节点的位置
（有功分点与无功分点可能不同）。以功率分点为界，把环网\textbf{解开成两个开式网}，
再按 11.1 节的方法分别计算电压损耗。全网最大电压损耗出现在电源到功率分点的路径上。

\subsection{含变压器的简单环网的功率分布}

环网中含变压器时，各段电压等级不同，必须归算到同一电压级。
若各变压器的变比与基准变比不匹配，即使两端电压相等也会产生循环功率：
\begin{equation}
  S_{cir}=\frac{U_N\Delta U}{\sqrt{3}\,\hat Z_{\Sigma}},
  \qquad \Delta U=\sum k_iU_{Ni}-U_N\ \text{（变比失配引起的等值电动势）}.
\end{equation}
因此，环网中并列运行的变压器要求\textbf{变比尽量一致}，否则会产生额外的无功环流与网损。

\subsection{环网中的潮流控制}

闭式网的功率自然分布由阻抗决定，未必等于经济分布。要改变潮流，须在环路中引入
\textbf{附加电动势}或改变阻抗：
\begin{itemize}
  \item \textbf{串联电容补偿}：减小某段电抗，使该段分担更多有功（也提高稳定极限）；
  \item \textbf{串联电抗}：增大某段阻抗，迫使功率走另一路径（用于限制短路电流或均衡负载）；
  \item \textbf{移相变压器（相位调节器）}：在环路中引入可控相角差，直接调节有功分布，
        是控制环流最直接的手段；
  \item \textbf{开环运行}：把环网解开，牺牲可靠性换取潮流可控。
\end{itemize}

\corepoint{环网潮流的三个必考点：（1） 循环功率与负荷无关，两端电压差（或变比失配）一存在，
环流就存在，且全部转化为网损；（2） 功率分点是“解开环网”的位置，找错了后续电压计算全错；
（3） 均一网按长度反比分配，是唯一可以不算阻抗的情形。}

\section{复杂电力系统潮流计算的数学模型}

\subsection{潮流计算的定解条件}

对 $n$ 个节点的网络，节点方程 $\boldsymbol Y\dot{\boldsymbol U}=\dot{\boldsymbol I}$ 含 $2n$ 个实数未知量
（各节点电压幅值与相角）。每个节点可给定两个量，共 $2n$ 个条件，问题才可解。
按给定量的不同，节点分为三类：

\begin{table}[htbp]
  \centering
  \caption{潮流计算中的节点分类}
  \small
  \begin{tabular}{@{}p{0.14\linewidth}p{0.22\linewidth}p{0.20\linewidth}p{0.34\linewidth}@{}}
    \toprule
    节点类型 & 给定 & 待求 & 典型节点 \\
    \midrule
    PQ 节点 & $P$、$Q$ & $U$、$\delta$ & 负荷节点、给定有功出力的发电厂（无励磁调节） \\
    PV 节点 & $P$、$U$ & $Q$、$\delta$ & 有自动励磁调节的发电厂、调相机 \\
    平衡节点 & $U$、$\delta$ & $P$、$Q$ & 调频厂（全系统一般只设一个） \\
    \bottomrule
  \end{tabular}
\end{table}

平衡节点的作用是“吸收”全网的功率不平衡量（网损与给定出力的偏差），
其 $P$、$Q$ 由计算结果给出，不预先指定。设 PQ 节点 $m$ 个、PV 节点 $n-m-1$ 个，
则待求量为 $2m+(n-m-1)=n+m-1$ 个，方程数亦相同。

\subsection{潮流计算的约束条件}

数学上的解未必是物理上可接受的运行状态，须满足：
\begin{itemize}
  \item \textbf{电压约束}：$U_{i\min}\le U_i\le U_{i\max}$（一般 $\pm5\%\sim\pm10\%$）；
  \item \textbf{功率约束}：发电机有功、无功出力在其上下限内；
        节点注入功率不超过设备容量；
  \item \textbf{稳定约束}：功角 $\delta$ 须小于稳定极限对应的角度；
  \item \textbf{线路载流约束}：支路功率不超过热稳定极限。
\end{itemize}

\section{牛顿—拉夫逊法潮流计算}

\subsection{牛顿—拉夫逊法的基本原理}

设非线性方程组 $\boldsymbol f(\boldsymbol x)=\boldsymbol 0$，在初值 $\boldsymbol x^{(k)}$ 处线性化：
\begin{equation}
  \boldsymbol f\left(\boldsymbol x^{(k)}\right)+\boldsymbol J^{(k)}\Delta\boldsymbol x^{(k)}=\boldsymbol 0,
  \qquad \boldsymbol J=\frac{\partial\boldsymbol f}{\partial\boldsymbol x},
\end{equation}
\begin{equation}
  \Delta\boldsymbol x^{(k)}=-\left[\boldsymbol J^{(k)}\right]^{-1}\boldsymbol f\left(\boldsymbol x^{(k)}\right),
  \qquad \boldsymbol x^{(k+1)}=\boldsymbol x^{(k)}+\Delta\boldsymbol x^{(k)}.
\end{equation}
$\boldsymbol J$ 称\textbf{雅可比矩阵}。牛顿法的特点是\textbf{平方收敛}（误差按平方递减），
一般 $3\sim5$ 次迭代即达 $10^{-6}$ 精度；缺点是对初值敏感，
初值太差可能不收敛，且每次迭代都要重新形成并分解 $\boldsymbol J$。

\subsection{节点电压用直角坐标表示时的牛顿—拉夫逊法潮流计算}

设 $\dot U_i=e_i+jf_i$，则节点功率
\begin{equation}
  P_i+jQ_i=\dot U_i\sum_{j}\hat Y_{ij}\hat U_j ,
\end{equation}
实部、虚部各给出一个关于 $e$、$f$ 的二次方程。对 PQ 节点列 $\Delta P_i$、$\Delta Q_i$，
对 PV 节点列 $\Delta P_i$ 与 $U_i^{2}=e_i^{2}+f_i^{2}$ 的偏差方程。修正方程为
\begin{equation}
  \begin{bmatrix}\Delta\boldsymbol P\\ \Delta\boldsymbol Q\\ \Delta\boldsymbol U^{2}\end{bmatrix}
  =\begin{bmatrix}
    \boldsymbol H & \boldsymbol N\\
    \boldsymbol K & \boldsymbol L\\
    \boldsymbol R & \boldsymbol S
  \end{bmatrix}
  \begin{bmatrix}\Delta\boldsymbol e\\ \Delta\boldsymbol f\end{bmatrix}.
\end{equation}
直角坐标的优点是功率表达式为二次多项式、雅可比元素计算简单；
缺点是 PV 节点的电压约束要用 $e^{2}+f^{2}$ 的二次式表示，处理略繁。

\subsection{节点电压用极坐标表示时的牛顿—拉夫逊法潮流计算}

设 $\dot U_i=U_i\angle\delta_i$，则节点注入功率为
\keyformula{P_i=U_i\sum_{j=1}^{n}U_j\left(G_{ij}\cos\delta_{ij}+B_{ij}\sin\delta_{ij}\right),
  \qquad
  Q_i=U_i\sum_{j=1}^{n}U_j\left(G_{ij}\sin\delta_{ij}-B_{ij}\cos\delta_{ij}\right)}
其中 $\delta_{ij}=\delta_i-\delta_j$，$G_{ij}$、$B_{ij}$ 为节点导纳矩阵元素的实部、虚部。
修正方程（对 PQ 与 PV 节点列有功、对 PQ 节点列无功）为
\begin{equation}
  \begin{bmatrix}\Delta\boldsymbol P\\ \Delta\boldsymbol Q\end{bmatrix}
  =\begin{bmatrix}\boldsymbol H & \boldsymbol N\\ \boldsymbol M & \boldsymbol L\end{bmatrix}
  \begin{bmatrix}\Delta\boldsymbol\delta\\ \Delta\boldsymbol U/U\end{bmatrix},
\end{equation}
\begin{equation}
  H_{ij}=\frac{\partial P_i}{\partial\delta_j},\quad
  N_{ij}=U_j\frac{\partial P_i}{\partial U_j},\quad
  M_{ij}=\frac{\partial Q_i}{\partial\delta_j},\quad
  L_{ij}=U_j\frac{\partial Q_i}{\partial U_j}.
\end{equation}
极坐标下 PV 节点的 $U$ 已知，故其 $\Delta U$ 与 $\Delta Q$ 不进入方程，方程阶数比直角坐标略低，
是目前工程程序的主流形式。

\begin{table}[htbp]
  \centering
  \caption{两种坐标的牛顿法比较}
  \small
  \begin{tabular}{@{}p{0.14\linewidth}p{0.30\linewidth}p{0.38\linewidth}@{}}
    \toprule
    比较项 & 直角坐标 & 极坐标 \\
    \midrule
    状态量 & $e_i$、$f_i$ & $U_i$、$\delta_i$ \\
    功率表达式 & 二次多项式，无三角函数 & 含 $\sin$、$\cos$，计算量略大 \\
    PV 节点处理 & 用 $U^{2}=e^{2}+f^{2}$ 约束 & 直接令 $U$ 已知，方程更少 \\
    实用地位 & 早期程序 & 主流 \\
    \bottomrule
  \end{tabular}
\end{table}

\section{P-Q 分解法潮流计算}

\subsection{解耦的物理依据}

高压电网中 $R\ll X$，由第 10 章的两个结论：
\begin{itemize}
  \item 有功功率的传输主要取决于\textbf{电压相角差}；
  \item 无功功率的传输主要取决于\textbf{电压幅值差}。
\end{itemize}
故可近似认为 $\Delta P$ 只与 $\Delta\delta$ 有关、$\Delta Q$ 只与 $\Delta U$ 有关，
即雅可比矩阵中的 $\boldsymbol N$、$\boldsymbol M$ 可以略去，方程解耦为两组低阶方程。

\subsection{分解后的修正方程}

\keyformula{\Delta\boldsymbol P=-\boldsymbol B'\Delta\boldsymbol\delta,\qquad
  \Delta\boldsymbol Q=-\boldsymbol B''\frac{\Delta\boldsymbol U}{\boldsymbol U}}
其中 $\boldsymbol B'$ 由导纳矩阵的虚部构成（略去电阻，考虑变压器变比），
$\boldsymbol B''$ 由导纳矩阵的虚部构成（略去电阻与移相，计及并联支路）。
两组方程的系数矩阵都是\textbf{常数矩阵}，只需在迭代前形成并三角分解一次，
以后每次迭代只做前代、回代——这是 PQ 分解法比牛顿法快得多的原因。

\subsection{特点与适用}

\begin{table}[htbp]
  \centering
  \caption{牛顿法与 PQ 分解法的比较}
  \small
  \begin{tabular}{@{}p{0.16\linewidth}p{0.26\linewidth}p{0.34\linewidth}@{}}
    \toprule
    比较项 & 牛顿—拉夫逊法 & PQ 分解法 \\
    \midrule
    收敛性 & 平方收敛，$3\sim5$ 次 & 线性收敛（等比），迭代次数多（$8\sim20$ 次） \\
    每次迭代 & 重新形成、分解 $\boldsymbol J$ & 系数矩阵不变，只做前代回代 \\
    总速度 & 中等 & 快（大规模系统优势明显） \\
    适用 & 各类网络，含低压配网（$R$ 不可略） & 高压输电网（$R\ll X$） \\
    精度 & 严格 & 与牛顿法同解（收敛后相同） \\
    \bottomrule
  \end{tabular}
\end{table}

\begin{table}[htbp]
  \centering
  \caption{潮流计算的两条路线}
  \small
  \begin{tabular}{@{}p{0.14\linewidth}p{0.34\linewidth}p{0.40\linewidth}@{}}
    \toprule
    路线 & 方法 & 特点 \\
    \midrule
    手算 & 开式网逐段推算、闭式网功率分点解开、均一网按长度分配 & 物理图像清晰，
      只能算小系统；是理解潮流的入口 \\
    计算机 & 牛顿法（直角/极坐标）、PQ 分解法、高斯—赛德尔法 & 通用、快速；
      PQ 分解法适合高压大网，牛顿法适合含低压支路的网络 \\
    \bottomrule
  \end{tabular}
\end{table}

\corepoint{本章的答题骨架：（1） 判节点类型（PQ／PV／平衡）并数方程；
（2） 写功率方程（极坐标两个和式）；（3） 写修正方程（牛顿法用 $\Delta P$、$\Delta Q$，
PQ 分解法用 $\boldsymbol B'$、$\boldsymbol B''$）；（4） 说明收敛与约束校验。
手算题则是“设额定电压 $\to$ 末端向首端算功率 $\to$ 首端向末端算电压 $\to$ 迭代校验”，
或“两端供电网先算功率分布、找功率分点、再解成开式网”。}
  % 第 11 章
% ============================================================
% 第 12 章  电力系统的无功功率平衡和电压调整
% （教材：何仰赞、温增银《电力系统分析》下册 第 12 章）
% 结构对齐 12.1—12.4；例题、习题不录
% ============================================================

\chapter{电力系统的无功功率平衡和电压调整}

电压是衡量电能质量的重要指标。质量合格的电压应在\textbf{供电电压偏移、电压波动和闪变、
电网谐波、三相不对称程度}四个方面满足国家标准——本章只讨论第一项（电压偏移）。
核心逻辑是一条因果链：\textbf{无功功率平衡决定电压水平}，故调压的本质是\textbf{无功的平衡与再分配}。

\section{电力系统的无功功率平衡}

保证用户处电压接近额定值，是系统运行调整的基本任务之一。
系统的运行电压水平取决于\textbf{无功功率平衡}：各种无功电源的出力应能满足负荷和网络损耗
在额定电压下对无功的需求，否则电压就偏离额定值。以下假定频率维持额定不变。

\subsection{无功功率负荷和无功功率损耗}

\paragraph{无功负荷}
异步电动机在电力系统负荷（特别是无功负荷）中占比很大，故系统无功负荷的电压特性
主要由异步电动机决定。其消耗的无功功率为
\begin{equation}
  Q_M=Q_m+Q_\sigma=\frac{U^{2}}{X_m}+I^{2}X_\sigma ,
\end{equation}
$Q_m$ 为励磁功率，与电压平方成正比（电压较高时因饱和，稍有偏离二次曲线）；
$Q_\sigma$ 为漏抗中的无功损耗——若负载功率不变，电压降低时转差增大、定子电流增大，
故 $Q_\sigma$ 反而上升。两部分叠加的结果是：\textbf{电压降低时，无功负荷先略降后回升}，
这个非线性特性是电压稳定的物理根源（第 15 章）。

\paragraph{无功功率损耗}
\begin{itemize}
  \item \textbf{变压器}：励磁无功 $\Delta Q_0=I_0\%S_N/100$（与电压平方近似成正比）
        与绕组漏抗中的无功 $\Delta Q_T\approx U_s\%S^{2}/(100S_N)$（与负载平方成正比）。
        大型变压器满载时，无功损耗可达其额定容量的 $10\%$ 以上，不能忽略。
  \item \textbf{输电线路}：串联电抗消耗无功
        $\Delta Q_L=(P^{2}+Q^{2})X/U^{2}$；并联电纳发出容性无功（按本书约定取负）
        $\Delta Q_B=-BU^{2}$。$220$~kV 及以下线路 $\Delta Q_L>\Delta Q_B$，
        $500$~kV 及以上因电纳大，轻载时充电功率可能超过消耗，线路呈净容性。
\end{itemize}

\subsection{无功功率电源}

除发电机外，还有四类无功补偿装置：

\begin{table}[htbp]
  \centering
  \caption{无功电源的类型与特性}
  \small
  \begin{tabular}{@{}p{0.16\linewidth}p{0.26\linewidth}p{0.20\linewidth}p{0.28\linewidth}@{}}
    \toprule
    类型 & 无功输出 & 调节 & 特点 \\
    \midrule
    发电机 & 受 $P$-$Q$ 容量曲线限制 & 连续 & 最经济，但受有功出力与稳定限制 \\
    同步调相机 & 可发可吸 & 连续 & 过激发感性无功、欠激发容性无功；损耗大、投资高 \\
    静电电容器 & $Q_C=U^{2}/X_C$ & 分组投切（不连续） & 只能发感性无功；$Q\propto U^{2}$，电压低时出力骤降 \\
    静止补偿器 SVC & 可发可吸 & 快速、平滑 & 由饱和电抗器／TCR／TSC 与固定电容组合；响应快 \\
    静止无功发生器 SVG & 可发可吸 & 极快、范围宽 & 又称 STATCOM；谐波小、占地小 \\
    \bottomrule
  \end{tabular}
\end{table}

\textbf{静电电容器}的无功输出与电压平方成正比：
\begin{equation}
  Q_C=\frac{U^{2}}{X_C},
\end{equation}
这既是它的优点（电压低时少发，有一定自调节性），也是缺点（系统最需要无功时它出力最少）。
实际装置分组接入或切除，实现补偿功率的不连续调节。

\textbf{静止补偿器}由静电电容器与电抗器（经晶闸管或饱和电抗器控制）组合而成：
\begin{itemize}
  \item 饱和电抗器型：工作范围内电压少许变化即引起电流大幅变化；与电抗器串联的电容用于
        斜率校正，改变电容大小可调节补偿器伏安特性的斜率；
  \item TCR 型（晶闸管控制电抗器）$+$ 固定电容：连续可调，但产生谐波，需装滤波器；
  \item TSC 型（晶闸管投切电容器）：不产生谐波，分级投切；
  \item 静止无功发生器（SVG／STATCOM）：响应更快、运行范围更宽、谐波电流小。
\end{itemize}
静止补偿器能在电压变化时快速、平滑地调节无功，满足\textbf{动态无功补偿}的需要；
TCR、TSC 型还可分相补偿以适应不平衡负荷。

\subsection{无功功率平衡}

无功平衡的条件是：在额定电压下，
\keyformula{\sum Q_{G}+\sum Q_{C}=\sum Q_{LD}+\Delta Q_{\Sigma}+Q_{res}}
即无功电源出力之和等于无功负荷、网络无功损耗与备用容量之和。
\begin{itemize}
  \item 无功备用容量一般为最大无功负荷的 $7\%\sim8\%$；
  \item 有功备用按最大有功负荷的百分数留，无功补偿装置按额定容量计；
  \item 330~kV 及以上超高压电网中，线路充电功率可满足无功需求的 $90\%$ 以上，
        须配置\textbf{并联电抗器}吸收过剩充电功率，同时在较低电压等级的配电网配置
        并联电容补偿——两者并存是正常的。
\end{itemize}
\textbf{配置原则}：按无功平衡需要增添补偿容量，并按\textbf{无功就地平衡}原则分配——
小容量、分散的补偿用静电电容器；大容量、配置在系统中枢点的补偿宜用同步调相机或静止补偿器。
采用补偿措施时应统筹兼顾，选用既能发又能吸无功的设备；有大量超高压线路时，
还应考虑发电机\textbf{进相运行}（功率因数超前）的可能性。

\subsection{无功功率平衡和电压水平的关系}

无功平衡与电压水平是同一件事的两种说法：
\begin{itemize}
  \item 无功充足时，系统可维持在额定电压附近运行；
  \item 无功不足时，电压被迫降低——因为负荷的无功需求随电压降低而减少，
        在新的（较低的）电压下重新达到平衡；
  \item 因此\textbf{电压水平是系统无功充裕程度的标志}，而不是可以任意“调”出来的量。
\end{itemize}
这一条常被误解：以为调压只是“把电压调上去”。实际上若无功电源不足，
强行提高电压会使无功更加不足，最终导致电压崩溃。

\corepoint{“无功不足则电压偏低”要说到位：不是“电压低所以无功不足”，而是
\textbf{无功不足 $\to$ 负荷在低电压下减少无功吸收 $\to$ 在新的低电压下平衡}。
所以调压的根本手段是\textbf{增加无功电源或减少无功损耗}，改变变比只是把已有无功
在空间上重新分配，不能增加无功总量。}

\section{电压调整的基本概念}

\subsection{允许电压偏移}

用电设备的额定电压与电网额定电压一致，但实际电压总偏离额定值。
允许偏移（第 1 章）：35~kV 及以上 $\pm5\%$，10~kV 及以下 $\pm7\%$；
对电压质量要求高的用户另有更严的规定。

\subsection{中枢点的电压管理}

由于负荷点很多、不可能对每个负荷点都直接调压，运行中选定一些\textbf{电压中枢点}
（枢纽变电所母线、重要发电厂母线），只要控制中枢点电压在允许范围内，
就能保证由其供电的所有负荷点电压合格。

\begin{table}[htbp]
  \centering
  \caption{中枢点的三种调压方式}
  \small
  \begin{tabular}{@{}p{0.16\linewidth}p{0.34\linewidth}p{0.40\linewidth}@{}}
    \toprule
    方式 & 做法 & 适用 \\
    \midrule
    逆调压 & 最大负荷时升到 $1.05U_N$，最小负荷时降到 $U_N$ & 负荷变化大、线路长、电压损耗大的网络 \\
    顺调压 & 最大负荷时不低于 $1.025U_N$，最小负荷时不超过 $1.075U_N$ & 负荷变化小、线路短 \\
    常调压 & 任何负荷下都维持在同一水平（如 $1.02\sim1.05U_N$） & 负荷变化很小 \\
    \bottomrule
  \end{tabular}
\end{table}

中枢点电压的允许变化范围，由“各负荷点的允许电压范围 $+$ 中枢点到各负荷点的电压损耗”
反推得到：中枢点电压下限 $=$ 电压最低负荷点的允许下限 $+$ 该路径的电压损耗；
上限 $=$ 电压最高负荷点的允许上限 $+$ 其路径的电压损耗。

\subsection{电压调整的基本原理}

以“发电机—升压变—线路—降压变—负荷”的典型供电链为例，负荷点电压为
\keyformula{U_2=\frac{U_G-\Delta U_{\text{网}}-\Delta U_T}{k}}
其中 $U_G$ 为发电机端电压，$\Delta U_{\text{网}}$ 为线路电压损耗，$\Delta U_T$ 为变压器
归算到高压侧的电压损耗，$k$ 为降压变压器变比。

由该式立刻看出调压的三条途径：
\begin{enumerate}
  \item \textbf{改变 $U_G$}：提高发电机端电压（发电机调压）；
  \item \textbf{减小 $\Delta U$}：减少线路与变压器的电压损耗——改变无功分布
        （无功补偿）、减小电抗（串联电容补偿）、提高电压等级；
  \item \textbf{改变 $k$}：改变变压器变比（分接头、有载调压变压器）。
\end{enumerate}

\section{电压调整的措施}

\subsection{发电机调压}

通过调节励磁电流改变发电机端电压。这是\textbf{最经济合理}的调压方式，不需另增设备。
\begin{itemize}
  \item 对孤立运行的发电厂，若发电机到最远负荷点的总电压损耗在允许范围内，
        仅靠发电机调压（如逆调压）即可满足要求；
  \item 多级变压供电系统中，最大负荷时总电压损耗可达 $35\%$、最小负荷时 $15\%$，
        变化幅度达 $20\%$，单靠发电机调压（一般只有 $\pm5\%$）不够；
  \item 多机并列运行时，各机端电压须协调，否则无功分配不合理，故发电机调压一般只作
        \textbf{辅助性措施}。
\end{itemize}

\subsection{改变变压器变比调压}

双绕组变压器在\textbf{高压绕组}设若干分接头（主接头对应 $U_N$），三绕组变压器一般在
高压和中压绕组设分接头，低压绕组不设。改变变比调压即“适当选择分接头”。

\paragraph{降压变压器分接头的选择}
已知高压侧电压 $U_1$、归算到高压侧的变压器电压损耗 $\Delta U_T$、低压侧要求的电压 $U_2$，
变比 $k=U_{1t}/U_{2N}$，则
\begin{equation}
  U_2=\frac{U_1-\Delta U_T}{k}=\frac{\left(U_1-\Delta U_T\right)U_{2N}}{U_{1t}},
\end{equation}
故所需高压侧分接头电压
\keyformula{U_{1t}=\frac{\left(U_1-\Delta U_T\right)U_{2N}}{U_2}}
按最大、最小负荷分别算出 $U_{1t.\max}$、$U_{1t.\min}$，取其算术平均值
$U_{1t.av}$，选最接近的标准分接头，然后\textbf{校验}最大、最小负荷下低压母线电压是否
都在允许范围内。普通变压器只能在停电时改变分接头，运行中只能用同一个固定分接头。

\paragraph{升压变压器分接头的选择}
方法与降压变压器基本相同，但功率方向相反，电压损耗与高压侧电压的关系由“相减”变为“相加”：
\begin{equation}
  U_{1t}=\frac{\left(U_1+\Delta U_T\right)U_{2N}}{U_2}.
\end{equation}
选择发电厂升压变压器分接头时，须保证最大、最小负荷下发电机端电压都在允许范围内，
一般对发电机母线采用逆调压方式。

\paragraph{有载调压变压器}
当“最大与最小负荷下的电压变化幅度超过分接头可调范围（如 $\pm5\%$）”，
或“调压要求的变化趋势与实际相反”（如逆调压）时，靠固定分接头无法满足，须装设
\textbf{带负荷调压}的变压器：
\begin{itemize}
  \item 有载调压变压器：本身具有调压绕组，可在带负荷条件下切换分接头，调节范围较大
        （可达 $\pm8\times1.25\%$ 等）；
  \item 带有附加调压器的加压调压变压器：用附加绕组串入主回路，产生纵、横向附加电压。
\end{itemize}

\subsection{利用无功功率补偿调压}

在负荷点附近装设无功补偿装置，减少流过线路的无功，从而减小电压损耗。
设简单辐射网首端电压 $U_1$ 恒定、末端负荷 $S=P+jQ$、线路阻抗 $R+jX$。
补偿前末端电压
\begin{equation}
  U_2=U_1-\frac{PR+QX}{U_1},
\end{equation}
装设补偿容量 $Q_C$ 后（补偿后末端电压 $U_{2c}$）
\begin{equation}
  U_{2c}=U_1-\frac{PR+\left(Q-Q_C\right)X}{U_1},
\end{equation}
两式相减得所需补偿容量
\keyformula{Q_C=\frac{U_1\left(U_{2c}-U_2\right)}{X}}
上式略去了补偿装置本身的无功电压特性（电容器的 $Q_C\propto U^{2}$），用于初步估算；
精确计算时应迭代修正，因为补偿后电压升高会使电容器实际出力增大。
若已知的是“补偿后要求达到的电压”，也可写成
$Q_C=Q-\dfrac{U_1\left(U_1-U_{2c}\right)-PR}{X}$，两式等价。

\paragraph{补偿装置的选择}
\begin{itemize}
  \item 需要\textbf{容性无功}（提高电压）时用电容器或调相机；
  \item 需要既能发又能吸（电压波动大、需动态补偿）时用静止补偿器或调相机；
  \item 补偿点应靠近无功负荷，遵循“就地平衡、分层分区”原则，避免无功长距离输送。
\end{itemize}

\subsection{线路串联电容补偿调压}

在线路中串联电容器，用电容的负电抗抵消部分线路电抗，减小电压损耗：
\begin{equation}
  \Delta U=\frac{PR+Q\left(X-X_C\right)}{U},
\end{equation}
若要求末端电压由 $U_2$ 提高到 $U_{2c}$，可解出所需串联容抗
\begin{equation}
  X_C=X-\frac{U_1\left(U_1-U_{2c}\right)-PR}{Q}.
\end{equation}
定义\textbf{补偿度}
\begin{equation}
  K_C=\frac{X_C}{X},
\end{equation}
一般取 $K_C<1$（避免谐振与过补偿）。串联电容补偿的优点是提高电压质量的同时
\textbf{提高输电能力与静态稳定极限}（因 $P_{\max}=U_1U_2/(X-X_C)$ 增大），
缺点是需保护（短路时电容器过电压）、可能引发次同步谐振。

\begin{table}[htbp]
  \centering
  \caption{四种调压措施的比较}
  \small
  \begin{tabular}{@{}p{0.14\linewidth}p{0.22\linewidth}p{0.24\linewidth}p{0.30\linewidth}@{}}
    \toprule
    措施 & 手段 & 效果 & 适用与限制 \\
    \midrule
    发电机调压 & 改变 $U_G$ & 直接、经济 & 范围有限（$\pm5\%$），多机并列需协调 \\
    改变变比 & 选分接头、有载调压 & 不增加无功，只重分配 & 固定分接头须停电切换；有载调压投资高 \\
    无功补偿 & 并联电容、调相机、SVC & 减少无功传输，降低电压损耗 & 需有可控无功电源；电容器 $Q\propto U^{2}$ \\
    串联电容 & 减小线路电抗 & 同时提高电压质量与稳定极限 & 需保护，可能谐振；补偿度受限 \\
    \bottomrule
  \end{tabular}
\end{table}

\corepoint{四类措施的本质区别：\textbf{发电机调压与无功补偿是“增加或就地提供无功”，
改变变比是“把已有无功在空间上重新分配”，串联电容是“减少无功的传输损耗”}。
所以无功总量不足时，只有前两类（以及新增补偿）能真正解决问题；只调分接头，
电压只能“此升彼降”。}

\section{调压措施的应用}

\subsection{各种调压措施的合理应用}

调压方案的选择要按“经济性 $+$ 有效性”顺序考虑：
\begin{enumerate}
  \item 首先充分利用\textbf{发电机调压}（不花钱）；
  \item 电压损耗不大而负荷变化较大时，用\textbf{改变变比}（尤其是逆调压场合用有载调压变压器）；
  \item 无功不足或无功传输距离长时，用\textbf{无功补偿}（就地平衡）；
  \item 远距离输电、需要同时提高稳定极限时，用\textbf{串联电容补偿}；
  \item 以上措施组合使用，并校验在各种运行方式下（最大／最小负荷、检修方式）
        都能满足电压要求。
\end{enumerate}

\subsection{各种措施调压效果的综合分析}

不同措施对电压的影响在“空间”和“时间”上不同：
\begin{itemize}
  \item \textbf{空间上}：无功补偿影响其装设点及其邻近节点；改变变比只影响变压器二次侧
        及其下游；发电机调压影响全网（但衰减快）。
  \item \textbf{时间上}：发电机调压、SVC、有载调压变压器可连续、快速调节；
        固定分接头与分组投切电容器只能分级、慢速调节。
  \item \textbf{耦合性}：多种措施同时作用时效果不是简单叠加——例如提高某点电压会增加
        该点负荷的无功需求，抵消一部分调压效果；无功补偿与变压器变比同时调整时，
        须联合优化，否则可能出现“补偿容量增加但电压不变”的情形。
\end{itemize}

\corepoint{调压计算的答题顺序：（1） 由允许电压偏移定中枢点（或负荷点）的电压上下限；
（2） 算出各元件在最大、最小负荷下的电压损耗；（3） 判断单靠发电机调压是否够；
（4） 不够则选分接头（用 $U_{1t}$ 公式，取平均、校验两端）；（5） 仍不够则加无功补偿
（用 $Q_C$ 公式）；（6） 最后校验所有运行方式。每一步都要“算出数、再校验”，
不能只写公式。}
  % 第 12 章
% ============================================================
% 第 13 章  电力系统的有功功率平衡和频率调整
% （教材：何仰赞、温增银《电力系统分析》下册 第 13 章）
% 结构对齐 13.1—13.4；例题、习题不录
% ============================================================

\chapter{电力系统的有功功率平衡和频率调整}

第 12 章讲“无功—电压”，本章讲“有功—频率”，二者是稳态运行的两根支柱，结构完全平行：
平衡决定水平，调节手段实现平衡。区别在于：频率是\textbf{全网统一}的量，
不存在“频率高低不同”的问题；因此调频必然涉及\textbf{谁来做主调频厂}、\textbf{互联系统怎样分区}
这类组织问题，这是调压所没有的。

\section{频率调整的必要性}

频率偏离额定值对各方面都有害：
\begin{itemize}
  \item \textbf{对用户}：异步电动机转速与频率成正比，频率下降使转速下降，
        影响纺织、造纸等对转速敏感的行业；电动机出力下降、电流增大，可能过热；
        电子设备、钟表计时失准。
  \item \textbf{对发电厂}：汽轮机叶片有固有振动频率，频率偏离会引起叶片共振、疲劳损伤；
        厂用电动机（给水泵、风机）转速下降，导致锅炉出力下降，形成恶性循环。
  \item \textbf{对系统}：频率下降说明有功不足，严重时引发频率崩溃，甚至系统瓦解。
\end{itemize}
我国额定频率 $50$~Hz，正常运行时允许偏移 $\pm0.2\sim\pm0.5$~Hz（大系统取严、小系统取宽）。
频率主要取决于系统的\textbf{有功功率平衡}：发出的有功不足，频率就偏低。

\begin{table}[htbp]
  \centering
  \caption{频率偏移的允许值与危害}
  \small
  \begin{tabular}{@{}lll@{}}
    \toprule
    偏移范围 & 状态 & 主要危害 \\
    \midrule
    $\pm0.2$ Hz 以内 & 正常 & 无 \\
    $\pm0.2\sim\pm0.5$ Hz & 允许 & 对少数敏感用户有影响 \\
    $>0.5$ Hz & 异常 & 电动机过热、汽轮机叶片振动 \\
    $-1\sim-2$ Hz 以上 & 事故 & 厂用电动机失速、频率崩溃 \\
    \bottomrule
  \end{tabular}
\end{table}

\section{电力系统的频率特性}

频率调整之所以可能，是因为“负荷”与“发电机”都对频率有\textbf{自身调节作用}——
这是理解一次调频的关键。二者统称系统的\textbf{有功功率—频率静态特性}。

\subsection{系统负荷的有功功率—频率静态特性}

负荷从电网取用的有功功率随频率变化的关系。频率下降时，电动机转速下降、机械负载功率下降，
取用的有功随之减少。在额定频率附近可线性化：
\keyformula{K_L=\frac{\Delta P_L}{\Delta f}\quad(\text{MW/Hz}),\qquad
  k_L=\frac{K_Lf_N}{P_{LN}}=1\sim3}
$K_L$ 称\textbf{负荷的频率调节效应系数}（负荷的单位调节功率），$k_L$ 为其标幺值。
$k_L$ 的物理含义是：频率变化 $1\%$ 时，负荷有功变化 $k_L\%$。
它由负荷构成决定，不可人为整定，是系统的固有属性。

\subsection{发电机组的有功功率—频率静态特性}

发电机组（或调频器未动作时的原动机调速系统）具有“频率下降 $\to$ 调速器动作 $\to$
进汽（水）量增加 $\to$ 出力增加”的特性，称\textbf{发电机组的静态调节特性}：
\begin{equation}
  \Delta f=-\delta\,f_N\frac{\Delta P_G}{P_{GN}},
\end{equation}
其中 $\delta$ 称\textbf{调差系数}（一般 $4\%\sim5\%$），其倒数为机组的单位调节功率
\keyformula{K_G=-\frac{\Delta P_G}{\Delta f}=\frac{P_{GN}}{\delta f_N}\quad(\text{MW/Hz}),
  \qquad k_G=\frac{1}{\delta}}
$K_G$ 称\textbf{发电机组的功率—频率静态特性系数}（单位调节功率）。
机组的调差系数可以整定（改变调速器特性），故 $K_G$ 是\textbf{可调}的；
但 $\delta$ 不能取得太小（$K_G$ 太大会引起调速系统振荡），也不能取得太大（调频能力弱）。

\subsection{电力系统的有功功率—频率静态特性}

系统同时含发电机组与负荷，二者的调节作用叠加：
\keyformula{K_S=K_G+K_L,\qquad k_S=k_G+k_L}
$K_S$ 称\textbf{系统的单位调节功率}。它表示：负荷增加 $1$~MW 时，
在发电机调速器与负荷自身调节的共同作用下，系统频率下降多少赫兹。

\begin{table}[htbp]
  \centering
  \caption{三个调节系数的比较}
  \small
  \begin{tabular}{@{}p{0.12\linewidth}p{0.24\linewidth}p{0.20\linewidth}p{0.32\linewidth}@{}}
    \toprule
    系数 & 定义 & 可否整定 & 说明 \\
    \midrule
    $K_L$（$k_L$） & $\Delta P_L/\Delta f$ & 不可 & 负荷固有属性，$k_L=1\sim3$ \\
    $K_G$（$k_G$） & $-\Delta P_G/\Delta f$ & 可（改 $\delta$） & $\delta=4\%\sim5\%$，$k_G=20\sim25$ \\
    $K_S$（$k_S$） & $K_G+K_L$ & 部分可 & 系统调频能力的总度量 \\
    \bottomrule
  \end{tabular}
\end{table}

\corepoint{注意符号约定：$K_L$ 定义中负荷功率随频率下降而减少，即“频率降、负荷降”，
故 $K_L>0$；$K_G$ 定义中加了负号，使“频率降、出力增”对应 $K_G>0$。
两者同号相加才表示“共同承担负荷增量”。若符号搞反，一次调频的公式会差一个负号。}

\section{电力系统的频率调整}

\subsection{频率的一次调整}

负荷突然增加 $\Delta P_{L0}$，调速器动作（一次调整），发电机增发、负荷因频率下降而少取，
二者共同平衡这一增量。设频率下降 $|\Delta f|$，则
\begin{equation}
  \Delta P_{L0}=K_G\left|\Delta f\right|+K_L\left|\Delta f\right|=K_S\left|\Delta f\right|,
\end{equation}
故
\keyformula{\Delta f=-\frac{\Delta P_{L0}}{K_S}=-\frac{\Delta P_{L0}}{K_G+K_L}}
一次调整的特点是：
\begin{itemize}
  \item \textbf{有差调节}：频率不可能回到原值（$\Delta f\ne0$），只能减小偏差；
  \item 各机组按各自调差系数\textbf{自动分配}负荷增量（$\delta$ 小的机组多担）；
  \item 调整过程快速（秒级），是频率的“第一道防线”；
  \item $K_S$ 越大，频率偏差越小。增加 $K_G$（减小 $\delta$）可改善，但受稳定性限制。
\end{itemize}

\textbf{机组间的负荷分配}：由 $K_G$ 的定义直接可得，第 $i$ 台机组承担的增量
\begin{equation}
  \Delta P_{Gi}=K_{Gi}\left|\Delta f\right|
  =\frac{P_{GNi}}{\delta_i f_N}\left|\Delta f\right|,
\end{equation}
即
\keyformula{\frac{\Delta P_{Gi}}{\Delta P_{Gj}}=\frac{K_{Gi}}{K_{Gj}}
  =\frac{P_{GNi}\delta_j}{P_{GNj}\delta_i}}
调差系数 $\delta$ 小的机组多担负荷，容量大的机组也多担负荷。
这一分配是\textbf{自动、无需通信}的，由调速器自身的下垂特性决定，
是“分散自治”控制的经典范例。但它也带来一个问题：$K_G$ 大的机组长期过载、
$K_G$ 小的机组长期欠载，故二次调整时须人为修正。

\subsection{频率的二次调整}

一次调整后频率仍有偏差，须由调频器（二次调整）动作，人为增加发电出力：
\begin{equation}
  \Delta P_{L0}-\Delta P_{G0}=K_S\left|\Delta f\right|,
\end{equation}
即
\keyformula{\Delta f=-\frac{\Delta P_{L0}-\Delta P_{G0}}{K_G+K_L}}
当 $\Delta P_{G0}=\Delta P_{L0}$ 时 $\Delta f=0$，称\textbf{无差调节}。
二次调整由调频厂（或 AGC）完成，是频率的“第二道防线”，时间尺度为分钟级。

\begin{table}[htbp]
  \centering
  \caption{一次调整与二次调整的比较}
  \small
  \begin{tabular}{@{}p{0.16\linewidth}p{0.34\linewidth}p{0.36\linewidth}@{}}
    \toprule
    比较项 & 一次调整 & 二次调整 \\
    \midrule
    动作元件 & 调速器 & 调频器（AGC） \\
    时间尺度 & 秒级 & 分钟级 \\
    调节性质 & 有差 & 可无差 \\
    承担者 & 所有机组（按 $\delta$ 分配） & 主调频厂（或分区） \\
    \bottomrule
  \end{tabular}
\end{table}

\subsection{互联系统的频率调整}

互联系统通过联络线相连，频率调整须考虑两个约束：\textbf{本区频率}与\textbf{联络线功率}。
设 A、B 两区经联络线相连，A 区负荷增加 $\Delta P_{LA}$，则
\begin{align}
  \Delta P_{LA}-\Delta P_{GA}&=K_A\Delta f+\Delta P_{AB},\\
  -\Delta P_{GB}&=K_B\Delta f-\Delta P_{AB},
\end{align}
其中 $\Delta P_{AB}$ 为经联络线由 A 流向 B 的功率增量。联立可得
\begin{equation}
  \Delta f=-\frac{\Delta P_{LA}-\Delta P_{GA}-\Delta P_{GB}}{K_A+K_B},
  \qquad
  \Delta P_{AB}=\frac{K_B\left(\Delta P_{LA}-\Delta P_{GA}\right)-K_A\Delta P_{GB}}{K_A+K_B}.
\end{equation}
\textbf{分区调频}（联合电力系统）的要点：
\begin{itemize}
  \item 各区按“\textbf{联络线功率偏差 $+$ 频率偏差}”构成\textbf{区域控制误差}（ACE）
        作为调频信号，使本区只承担自己的负荷变化，不干扰邻区；
  \item 若各区都只按本区 ACE 调整，则全网频率与联络线功率都能恢复；
  \item 若只有一区调频，则该区须承担全网的负荷增量，经济上不合理。
\end{itemize}

\subsection{主调频厂的选择}

主调频厂应满足：
\begin{enumerate}
  \item 有足够的\textbf{调频容量}（一般不小于系统最大负荷的 $2\%\sim5\%$，并留事故备用）；
  \item 出力调整\textbf{速度快}、范围大、增减灵活；
  \item 调频时\textbf{经济性}好（效率高、耗量特性平坦）；
  \item 地理位置适当，调频引起的潮流变化不致使网络过载或电压越限。
\end{enumerate}
一般优先选\textbf{水电厂}（调整快、经济）或高效率大容量火电厂；核电厂一般带基荷，不参与调频。

\subsection{频率调整和电压调整的关系}

\begin{itemize}
  \item \textbf{对象不同}：频率靠有功平衡，电压靠无功平衡，二者相对独立；
  \item \textbf{耦合存在}：有功变化会改变网络电压损耗（$\Delta U=(PR+QX)/U$），
        从而影响电压；无功补偿装置（尤其电容器）的投切也会改变有功损耗；
  \item \textbf{调节速度不同}：频率调整要求更快（秒级—分钟级），电压调整可以慢一些；
  \item \textbf{调节手段不同}：调频靠原动机（调速器、调频器），调压靠励磁、变比、无功补偿。
\end{itemize}
工程上两者分别由“有功—频率”与“无功—电压”两套控制系统（AGC 与 AVC）承担。

\section{有功功率平衡和系统负荷在各类发电厂间的合理分配}

\subsection{有功功率平衡和备用容量}

系统有功平衡要求：所有发电厂发出的有功等于负荷、网损与厂用电之和。为应对负荷波动与事故，
须留\textbf{备用容量}：

\begin{table}[htbp]
  \centering
  \caption{备用容量的分类与取值}
  \small
  \begin{tabular}{@{}p{0.18\linewidth}p{0.30\linewidth}p{0.42\linewidth}@{}}
    \toprule
    备用类型 & 作用 & 一般取值 \\
    \midrule
    负荷备用 & 应对负荷短期波动（调频容量） & 最大负荷的 $2\%\sim5\%$ \\
    事故备用 & 应对机组或线路事故停运 & 最大负荷的 $5\%\sim10\%$（且不小于最大单机容量） \\
    检修备用 & 应对计划检修 & 按检修计划确定 \\
    国民经济备用 & 应对新用户、新增长 & 按发展规划确定 \\
    \bottomrule
  \end{tabular}
\end{table}

备用容量按“是否已并网”分为\textbf{热备用}（旋转备用，机组已并网、可立即增出力）与
\textbf{冷备用}（机组未启动或未并网，需启动时间）。热备用是调频与事故支援的主力。

\subsection{各类发电厂负荷的合理分配}

在满足安全约束的前提下，按\textbf{经济性}分配负荷，基本原则是：
\begin{enumerate}
  \item \textbf{核电厂}：一次投资大、效率高、不宜频繁调节，带\textbf{基荷}；
  \item \textbf{高效率大容量火电厂}：带基荷与腰荷；
  \item \textbf{中小容量火电厂}：带腰荷与峰荷，调节灵活；
  \item \textbf{水电厂}：启动快、调节灵活，宜带\textbf{峰荷}与调频、事故备用；
        但受水库调度（来水、灌溉、防洪）约束；
  \item \textbf{抽水蓄能电厂}：低谷抽水蓄能（消耗多余电能）、高峰发电，
        是调峰填谷的有效手段。
  \item 新能源（风电、光伏）：优先消纳，但其出力波动性要求更多备用与调节能力。
\end{enumerate}

\begin{table}[htbp]
  \centering
  \caption{各类电厂在负荷曲线上的位置}
  \small
  \begin{tabular}{@{}p{0.18\linewidth}p{0.20\linewidth}p{0.26\linewidth}p{0.26\linewidth}@{}}
    \toprule
    电厂类型 & 负荷位置 & 调节能力 & 分配要点 \\
    \midrule
    核电 & 基荷 & 差 & 满发、少调 \\
    大容量火电 & 基荷／腰荷 & 中等 & 按耗量特性排顺序 \\
    中小火电 & 腰荷／峰荷 & 较好 & 承担日负荷变化 \\
    水电 & 峰荷／备用 & 好（秒级） & 兼顾水库调度 \\
    抽水蓄能 & 填谷／峰荷 & 很好 & 低谷抽水、高峰发电 \\
    \bottomrule
  \end{tabular}
\end{table}

\corepoint{本章的主线：\textbf{频率偏差 }$\Delta f=-\Delta P/K_S$\textbf{ 是一切调频结论的母式}。
一次调整靠 $K_S$（有差），二次调整靠 $\Delta P_{G0}$（可无差），互联系统靠 ACE 分区平衡。
把“$K_L$ 不可调、$K_G$ 可调、$K_S$ 是二者之和”这三句记住，
再记住“备用容量四类、电厂分配按基荷—腰荷—峰荷排序”，本章的计算与简答就都有了骨架。}
  % 第 13 章
% ============================================================
% 第 14 章  电力系统的经济运行
% （教材：何仰赞、温增银《电力系统分析》下册 第 14 章）
% 结构对齐 14.1—14.4；例题、习题不录
% ============================================================

\chapter{电力系统的经济运行}

前两章解决“电压合格”“频率合格”，本章解决“花最少的燃料与最少的损耗，发出同样的电”。
经济运行有两条战线：\textbf{降低网损}（输配环节）与\textbf{负荷经济分配}（发输环节）。
两者的数学工具都是\textbf{等微增率准则}——一个由拉格朗日乘子法导出的、极为简洁的结论：
\textbf{各机组的耗量微增率相等时，总耗量最小}。

\section{电力网中的能量损耗}

\subsection{电力网的能量损耗和损耗率}

电力网的能量损耗由\textbf{线路}和\textbf{变压器}中的电阻损耗（可变损耗）与
\textbf{变压器励磁、线路电晕}等固定损耗构成。衡量指标是\textbf{网损率}（线损率）：
\keyformula{\text{网损率}=\frac{\Delta W}{W_{\text{供电}}}\times100\%
  =\frac{W_{\text{供电}}-W_{\text{售电}}}{W_{\text{供电}}}\times100\%}
我国电网的网损率一般在 $5\%\sim8\%$ 之间。降低网损是电力系统节能的主要着力点，
因为网损是“发了电却没用上”的部分。

网损可分解为两部分，二者的变化规律完全不同：

\begin{table}[htbp]
  \centering
  \caption{网损的两个组成部分}
  \small
  \begin{tabular}{@{}p{0.18\linewidth}p{0.32\linewidth}p{0.38\linewidth}@{}}
    \toprule
    类别 & 组成 & 随负荷的变化 \\
    \midrule
    可变损耗（铜耗） & 线路电阻、变压器绕组电阻中的损耗 & 与电流平方成正比，
      随负荷剧烈变化，$\Delta P\propto S^{2}$ \\
    固定损耗（铁耗） & 变压器励磁损耗、线路电晕损耗 & 与电压有关、与负荷几乎无关，
      只要设备带电就存在 \\
    \bottomrule
  \end{tabular}
\end{table}

这一分解解释了两个反直觉的现象：（1）轻载时网损率反而升高（固定损耗不变、售电量减少）；
（2）多台变压器并列轻载运行不如少投几台经济（见 14.1.3 的临界负载系数）。

\subsection{线路中能量损耗的计算方法}

线路的有功损耗随负荷变化，全年损耗须对时间积分：
\begin{equation}
  \Delta W=\int_{0}^{8760}\Delta P(t)\,\mathrm{d}t
  =\int_{0}^{8760}\frac{P^{2}(t)+Q^{2}(t)}{U^{2}(t)}R\,\mathrm{d}t .
\end{equation}
直接积分需全年负荷曲线，工程上用\textbf{最大负荷损耗时间}折算：
\keyformula{\Delta W=\Delta P_{\max}\,\tau}
其中 $\Delta P_{\max}$ 为全年最大负荷时的功率损耗，$\tau$ 称\textbf{最大负荷损耗时间}
（一年中若损耗始终等于 $\Delta P_{\max}$，经过 $\tau$ 小时所耗电能恰等于全年实际损耗）。
$\tau$ 与\textbf{最大负荷利用小时数} $T_{\max}$、功率因数 $\cos\varphi$ 有关，
可由曲线或表格查得，一般 $\tau<T_{\max}$（因损耗与电流平方成正比，轻载时损耗小）。

\begin{table}[htbp]
  \centering
  \caption{三个“时间”概念的区别}
  \small
  \begin{tabular}{@{}p{0.22\linewidth}p{0.30\linewidth}p{0.34\linewidth}@{}}
    \toprule
    量 & 定义 & 用途 \\
    \midrule
    $T_{\max}$（最大负荷利用小时数） & $W/P_{\max}$ & 衡量负荷平稳程度（第 9 章） \\
    $\tau$（最大负荷损耗时间） & $\Delta W/\Delta P_{\max}$ & 由负荷曲线折算全年电能损耗 \\
    $t$（运行小时数） & 实际运行时间 & 设备可用率统计 \\
    \bottomrule
  \end{tabular}
\end{table}

\subsection{降低网损的技术措施}

\begin{table}[htbp]
  \centering
  \caption{降低网损的措施及其机理}
  \small
  \begin{tabular}{@{}p{0.20\linewidth}p{0.34\linewidth}p{0.34\linewidth}@{}}
    \toprule
    措施 & 机理 & 备注 \\
    \midrule
    提高电压等级 & $\Delta P\propto1/U^{2}$ & 规划阶段措施，效果最显著 \\
    升高运行电压水平 & 同上 & 运行措施，受电压允许上限约束 \\
    无功就地补偿 & 减少无功传输，降低 $I$ 与 $\Delta P$ & 兼顾调压（第 12 章） \\
    合理确定运行方式 & 停运轻载变压器、解开环流、经济调度 & 运行措施，成本最低 \\
    增大导线截面 & 减小 $R$ & 需投资，按经济电流密度选择 \\
    减少变压器损耗 & 选用低损耗变压器、合理配置 & 固定损耗为主 \\
    降低环流（循环功率） & 消除变比失配、调整潮流 & 环流不传输功率却全部变成损耗 \\
    \bottomrule
  \end{tabular}
\end{table}

其中“\textbf{停运轻载变压器}”最容易被忽视：变压器空载损耗基本固定，
多一台变压器并网就多一份固定损耗，故轻载时应把负荷集中到少数变压器上。
反之，负载重时应多台并联以减小绕组损耗——这就是“变压器经济运行的临界负载系数”。

设两台同容量变压器，单台的额定铜耗为 $\Delta P_{cu}$、空载损耗为 $\Delta P_0$，
总负荷为 $S$。若只用一台，损耗为
\begin{equation}
  \Delta P_{(1)}=\Delta P_0+\Delta P_{cu}\left(\frac{S}{S_N}\right)^{2};
\end{equation}
若两台并联，损耗为
\begin{equation}
  \Delta P_{(2)}=2\Delta P_0+2\Delta P_{cu}\left(\frac{S}{2S_N}\right)^{2}.
\end{equation}
令两式相等，得\textbf{临界负载系数}
\keyformula{S_{cr}=S_N\sqrt{\frac{2\Delta P_0}{\Delta P_{cu}}}}
当 $S<S_{cr}$ 时投一台经济，$S>S_{cr}$ 时投两台经济。这一判据可直接推广到 $n$ 台并联的情形。

\section{火电厂间有功功率负荷的经济分配}

\subsection{耗量特性}

火电厂的燃料消耗 $F$（t/h 或元/h）与有功出力 $P$（MW）的关系称\textbf{耗量特性}，
通常呈下凸曲线。描述它的两个量：
\begin{itemize}
  \item \textbf{比耗量} $\mu=F/P$：单位出力的燃料消耗，$\mu$ 最小时对应最高效率点；
  \item \textbf{耗量微增率} $\lambda=\mathrm{d}F/\mathrm{d}P$：增加单位出力所增加的燃料消耗，
        它是经济分配的\textbf{判据量}。
\end{itemize}
耗量特性为下凸，意味着 $\lambda$ 随 $P$ 单调上升——这是等微增率准则有唯一解的前提。

\subsection{等微增率准则}

设系统有 $n$ 台机组（或电厂），总负荷 $P_L$ 给定，目标是最小化总耗量
\begin{equation}
  \min F_{\Sigma}=\sum_{i=1}^{n}F_i\left(P_i\right),
  \qquad \text{s.t.}\ \sum_{i=1}^{n}P_i=P_L .
\end{equation}
用拉格朗日乘子法，构造 $L=\sum F_i+\lambda\left(P_L-\sum P_i\right)$，令 $\partial L/\partial P_i=0$，得
\keyformula{\frac{\mathrm{d}F_1}{\mathrm{d}P_1}=\frac{\mathrm{d}F_2}{\mathrm{d}P_2}=\cdots
  =\frac{\mathrm{d}F_n}{\mathrm{d}P_n}=\lambda}
即\textbf{等微增率准则}：各机组的耗量微增率相等时，总耗量最小。
$\lambda$ 称\textbf{能量价格}（系统的边际成本），是负荷增加单位功率时总耗量的增量。

\textbf{适用条件}（必须写清）：
\begin{enumerate}
  \item 各机组的耗量特性为\textbf{下凸}（$\mathrm{d}^{2}F/\mathrm{d}P^{2}>0$），保证解唯一；
  \item 各机组出力满足上下限约束 $P_{i\min}\le P_i\le P_{i\max}$——
        越限的机组按极限运行，不参与等微增率分配；
  \item 不计网损（或已作网损修正）；
  \item 系统频率、电压维持额定，即只考虑有功的静态经济分配。
\end{enumerate}

\subsection{多个发电厂间的负荷经济分配}

对多个电厂，仍按等微增率分配，但须按\textbf{电厂总耗量特性}（由厂内机组等微增率先分配后
汇总得到）参与全网分配。做法分两层：
\begin{enumerate}
  \item 厂内：各机组按等微增率分配厂负荷；
  \item 全网：各厂按等微增率分配系统负荷。
\end{enumerate}
若某厂机组效率差异大，也可直接用“机组级”等微增率，效果更好。

\subsection{计及网损的有功负荷经济分配}

计及网损时，约束变为 $\sum P_i=P_L+\Delta P_{\Sigma}$，拉格朗日函数为
\begin{equation}
  L=\sum_{i}F_i+\lambda\left(P_L+\Delta P_{\Sigma}-\sum_{i}P_i\right),
\end{equation}
令 $\partial L/\partial P_i=0$ 得
\keyformula{\frac{\mathrm{d}F_i}{\mathrm{d}P_i}\cdot\frac{1}{1-\dfrac{\partial\Delta P_{\Sigma}}{\partial P_i}}
  =\lambda\qquad(i=1,2,\dots,n)}
$\partial\Delta P_{\Sigma}/\partial P_i$ 称\textbf{网损微增率}，
$1/(1-\partial\Delta P_{\Sigma}/\partial P_i)$ 称\textbf{网损修正系数}。
其物理含义：机组 $i$ 多发出 $1$~MW，扣除它引起的网损增量后，实际送到负荷的只有
$1-\partial\Delta P_{\Sigma}/\partial P_i$ MW，故其“有效”微增率被放大。
\textbf{离负荷中心近的机组网损微增率小，应多发}——这是“就近供电”的经济学解释。

\section{水、火电厂间有功功率负荷的经济分配}

\subsection{一个水电厂和一个火电厂间负荷的经济分配}

水电厂的特点是：\textbf{燃料不花钱，但水量有限}。故目标函数只含火电厂的燃料消耗，
约束除功率平衡外还有\textbf{水量约束}：
\begin{equation}
  \min F_T\left(P_T\right),\qquad
  \text{s.t.}\ P_T+P_H=P_L,\qquad \int_{0}^{T}W_H\left(P_H\right)\mathrm{d}t=W_{\Sigma},
\end{equation}
$W_H(P_H)$ 为水电厂的耗水量特性。构造拉格朗日函数并对时间积分，可得
\keyformula{\frac{\mathrm{d}F_T}{\mathrm{d}P_T}=\gamma\,\frac{\mathrm{d}W_H}{\mathrm{d}P_H}}
$\gamma$ 为\textbf{水的折算系数}（把水量折算为燃料量的乘子，单位 元/m$^3$ 或 t/m$^3$）。
其含义是：火电的燃料微增率应等于水电的“耗水微增率乘以水价”。
若 $\gamma$ 取得大，说明水宝贵，应少用水电、多带火电；反之则多用水电。

工程上常用\textbf{等微增率 $+$ 水火折算}的实用做法：先假定一个 $\gamma$，
按等微增率分配各时段出力，校验总用水量是否等于给定水量，不等则调整 $\gamma$ 重算
（即“$\gamma$ 迭代法”）。

\subsection{计及网损时若干个水、火电厂间负荷的经济分配}

含网损时，对每个电厂（火或水）都须乘网损修正系数，形式为
\begin{equation}
  \frac{\mathrm{d}F_i}{\mathrm{d}P_i}\cdot\frac{1}{1-\partial\Delta P_{\Sigma}/\partial P_i}
  =\lambda\quad(\text{火电厂}),\qquad
  \gamma\frac{\mathrm{d}W_j}{\mathrm{d}P_j}\cdot\frac{1}{1-\partial\Delta P_{\Sigma}/\partial P_j}
  =\lambda\quad(\text{水电厂}).
\end{equation}
这样，不同类型电厂的经济分配就统一到“\textbf{修正后的微增率相等}”这一条准则上。

\begin{table}[htbp]
  \centering
  \caption{三类经济分配问题的对照}
  \small
  \begin{tabular}{@{}p{0.18\linewidth}p{0.30\linewidth}p{0.34\linewidth}@{}}
    \toprule
    问题 & 目标函数 & 判据 \\
    \midrule
    火电之间（不计网损） & $\sum F_i$ & $\mathrm{d}F_i/\mathrm{d}P_i=\lambda$ \\
    火电之间（计网损） & $\sum F_i$ & 修正后微增率相等 \\
    水火之间 & $F_T$，水量约束 & $\mathrm{d}F_T/\mathrm{d}P_T=\gamma\,\mathrm{d}W_H/\mathrm{d}P_H$ \\
    \bottomrule
  \end{tabular}
\end{table}

\corepoint{等微增率准则的形式会变（加网损修正、加水火折算），但\textbf{结构不变}：
“各方的边际代价相等”。记住这一点，三类问题就只是一种问题的三种写法。
另外必须写清适用条件（耗量特性下凸、出力上下限、不计或已修正网损），
否则“相等”只是必要条件而非充分条件。}

\section{无功功率负荷的经济分配}

\subsection{等微增率准则的应用}

无功功率的经济分配与有功类似，但目标不同：无功不消耗燃料（除发电机组的附带损耗外），
分配无功的目的是\textbf{使网损最小}。设系统有若干无功电源（发电机、调相机、电容器），
其无功出力 $Q_i$，则网损 $\Delta P_{\Sigma}$ 是各 $Q_i$ 的函数，问题为
\begin{equation}
  \min\Delta P_{\Sigma}\left(Q_1,\dots,Q_n\right),
  \qquad \text{s.t.}\ \sum Q_i=Q_L+\Delta Q_{\Sigma}.
\end{equation}
用拉格朗日乘子法得
\keyformula{\frac{\partial\Delta P_{\Sigma}}{\partial Q_1}
  =\frac{\partial\Delta P_{\Sigma}}{\partial Q_2}=\cdots
  =\frac{\partial\Delta P_{\Sigma}}{\partial Q_n}}
即\textbf{各无功电源的网损微增率相等}时网损最小。这与有功的等微增率准则形式一致，
只是把“耗量微增率”换成了“网损微增率”。

\subsection{无功功率补偿的经济配置}

无功补偿装置（电容器）有投资成本，其经济配置须在“减少网损的收益”与“投资成本”之间权衡。
设第 $i$ 个节点装设补偿容量 $Q_{Ci}$，投资为 $K_i(Q_{Ci})$，年网损费用为 $C\Delta P_{\Sigma}$，
则目标为
\begin{equation}
  \min\left\{C\,\Delta P_{\Sigma}\left(Q_{C1},\dots,Q_{Cn}\right)+\sum_{i}K_i\left(Q_{Ci}\right)\right\}.
\end{equation}
对 $Q_{Ci}$ 求偏导并令其为零，得
\begin{equation}
  C\frac{\partial\Delta P_{\Sigma}}{\partial Q_{Ci}}+\frac{\mathrm{d}K_i}{\mathrm{d}Q_{Ci}}=0 ,
\end{equation}
即\textbf{各补偿点的“网损减少的边际收益”等于“补偿投资的边际成本”}。
工程近似结论：
\begin{itemize}
  \item 补偿应优先装设在\textbf{网损微增率大}的节点（一般是远离电源、无功缺额大的负荷点）；
  \item 补偿容量应按“\textbf{等网损微增率}”分配，而不是平均分配；
  \item 考虑投资时，单位容量投资低的地点应多装。
\end{itemize}

\begin{table}[htbp]
  \centering
  \caption{有功经济分配与无功经济分配的比较}
  \small
  \begin{tabular}{@{}p{0.16\linewidth}p{0.34\linewidth}p{0.34\linewidth}@{}}
    \toprule
    比较项 & 有功经济分配 & 无功经济分配 \\
    \midrule
    目标 & 燃料消耗最小 & 网损最小（或综合费用最小） \\
    判据量 & 耗量微增率 $\mathrm{d}F/\mathrm{d}P$ & 网损微增率 $\partial\Delta P/\partial Q$ \\
    约束 & 功率平衡、出力上下限 & 无功平衡、电压约束 \\
    价格乘子 & 能量价格 $\lambda$ & 网损电价 \\
    \bottomrule
  \end{tabular}
\end{table}

\corepoint{本章四个“微增率”不要混：\textbf{耗量微增率}（火电，$\mathrm{d}F/\mathrm{d}P$）、
\textbf{耗水微增率}（水电，$\mathrm{d}W/\mathrm{d}P$）、
\textbf{网损微增率}（$\partial\Delta P/\partial P$ 或 $\partial\Delta P/\partial Q$）、
\textbf{比耗量}（$F/P$，不是微增量）。等微增率准则比较的是“微增率”，不是“比耗量”——
这是最常见的概念错误。}
  % 第 14 章
% ============================================================
% 第 15 章  电力系统运行稳定性的基本概念
% （教材：何仰赞、温增银《电力系统分析》下册 第 15 章）
% 结构对齐 15.1—15.7；例题、习题不录
% ============================================================

\chapter{电力系统运行稳定性的基本概念}

前面十四章讲的是“稳态”——系统在一个平衡点附近运行。本章开始讲“稳定性”：
\textbf{系统受到扰动后能否回到（或进入）一个可接受的运行状态}。
稳定性问题的一切讨论都围绕一个变量展开——\textbf{功角 $\delta$}；
也围绕一个方程展开——\textbf{转子运动方程}。本章的任务是把这两个对象建立起来，
并区分清楚“静态稳定”“暂态稳定”“负荷稳定”“电压稳定”四个容易混淆的概念。

\section{概述}

电力系统稳定性的一般定义是：系统在某一运行状态下受到扰动后，
经过一段过渡过程，能够恢复到原来的运行状态（或进入一个新的可接受的运行状态），
则称系统在该运行状态下是\textbf{稳定}的；否则是\textbf{不稳定}的。

按扰动大小与物理机理，稳定性分为：

\begin{table}[htbp]
  \centering
  \caption{电力系统稳定性的分类}
  \small
  \begin{tabular}{@{}p{0.16\linewidth}p{0.20\linewidth}p{0.26\linewidth}p{0.26\linewidth}@{}}
    \toprule
    类型 & 扰动 & 判别对象 & 主要分析工具 \\
    \midrule
    静态稳定 & 小扰动 & 发电机间功角 & 小扰动法（线性化）、$K_p$ 储备系数 \\
    暂态稳定 & 大扰动（短路、切机、切负荷） & 发电机间功角 & 等面积定则、数值积分 \\
    负荷稳定 & 小扰动 & 异步电动机转差 & 电动机转矩—转差特性 \\
    电压稳定 & 小／大扰动 & 节点电压 & $P$-$U$ 曲线、$\mathrm{d}Q/\mathrm{d}U$ 判据 \\
    \bottomrule
  \end{tabular}
\end{table}

本章及第 16～19 章主要讨论\textbf{功角稳定}（发电机保持同步运行的能力），
这是电力系统稳定性中最经典、最主要的部分。电压稳定与负荷稳定在本章作概念性介绍。

\section{功角的概念}

\subsection{功角的定义}

同步发电机并网运行时，其转子转速与电网同步速一致（平均而言），
但转子磁极轴线与定子合成磁极轴线之间总有一个空间夹角，这个夹角就是\textbf{功角} $\delta$
（电角度）。它有两重含义：
\begin{itemize}
  \item \textbf{空间上}：转子磁极轴线领先定子（气隙）合成磁极轴线的电角度，
        反映发电机转子的实际位置；
  \item \textbf{时间上}：发电机空载电势 $\dot E_q$ 与机端电压 $\dot U$（或系统某参考点电压）
        之间的相位差，反映各发电机电动势相量的相对位置。
\end{itemize}
两者是同一件事的两种描述：转子位置决定电动势相位。

\subsection{功角与有功功率的关系}

对“发电机经电抗 $X$ 接无限大母线（电压 $U$ 恒定）”的简单系统，发电机输出的有功功率为
\keyformula{P_e=\frac{E_qU}{X}\sin\delta}
这就是\textbf{功角特性}。它说明：
\begin{itemize}
  \item 有功功率的传输必须靠功角差（第 10 章“有功靠相角差传输”的定量表述）；
  \item $\delta>0$ 表示发电机向系统送有功（发电机状态）；$\delta<0$ 表示吸收有功（电动机状态）；
  \item $P_e$ 是 $\delta$ 的正弦函数，$\delta=90^\circ$ 时达最大值 $P_{\max}=E_qU/X$。
\end{itemize}

\begin{table}[htbp]
  \centering
  \caption{功角的物理意义与运行状态}
  \small
  \begin{tabular}{@{}lll@{}}
    \toprule
    功角 $\delta$ & 状态 & 说明 \\
    \midrule
    $\delta>0$ & 发电机 & 转子磁极领先定子磁极，向系统送有功 \\
    $\delta=0$ & 空载 & 不送有功 \\
    $\delta<0$ & 电动机 & 吸收有功 \\
    $0<\delta<90^\circ$ & 静态稳定区 & $\mathrm{d}P/\mathrm{d}\delta>0$ \\
    $\delta=90^\circ$ & 静态稳定极限 & $P=P_{\max}$ \\
    $\delta>90^\circ$ & 不稳定区 & $\mathrm{d}P/\mathrm{d}\delta<0$ \\
    \bottomrule
  \end{tabular}
\end{table}

\corepoint{功角不是“几何角度”，而是\textbf{电磁功率与机械功率的平衡状态在角度上的表现}。
稳态时 $\delta$ 由 $P_e=P_T$ 唯一确定（$P_T$ 为原动机机械功率）；
暂态时 $\delta$ 由转子运动方程决定。因此讨论稳定就是讨论“$\delta$ 会不会越过临界点”。
这一句话把第 15～19 章全部串起来了。}

\section{静态稳定的初步概念}

\textbf{静态稳定}指系统在某一运行状态下受到\textbf{小扰动}后，能否恢复到原运行状态的能力。

由功角特性，设原动机机械功率 $P_T$ 不变，系统原运行于 $\delta_0$：
\begin{itemize}
  \item 若 $\delta_0<90^\circ$，则 $\mathrm{d}P_e/\mathrm{d}\delta>0$。当扰动使 $\delta$ 略增时，
        $P_e$ 增大， 电磁转矩大于机械转矩，使转子减速、$\delta$ 回落到 $\delta_0$——
        系统稳定；
  \item 若 $\delta_0>90^\circ$，则 $\mathrm{d}P_e/\mathrm{d}\delta<0$。$\delta$ 略增使 $P_e$ 减小，
        电磁转矩小于机械转矩，转子进一步加速、$\delta$ 继续增大——系统不稳定（失步）。
\end{itemize}
故\textbf{静态稳定的判据}为
\keyformula{\frac{\mathrm{d}P_e}{\mathrm{d}\delta}>0}
临界点为 $\delta=90^\circ$，对应 $P_{\max}=E_qU/X$。

工程上用\textbf{静态稳定储备系数}衡量裕度：
\keyformula{K_p=\frac{P_{\max}-P_0}{P_0}\times100\%}
正常运行时要求 $K_p\ge15\%\sim20\%$；事故后（或特殊运行方式）允许降低，但一般不低于 $10\%$。

\section{暂态稳定的初步概念}

\textbf{暂态稳定}指系统受到\textbf{大扰动}（短路、切除线路、切机、切负荷）后，
各发电机能否保持同步运行。分析方法与静态稳定不同：不能用线性化，
必须研究\textbf{大扰动后转子运动的整个过渡过程}。

以“单机—无限大系统”发生短路为例：
\begin{enumerate}
  \item \textbf{故障期间}：短路使发电机端电压降低，电磁功率 $P_e$ 大幅下降，
        而机械功率 $P_T$ 不变，故转子受到加速转矩，$\delta$ 不断增大，转子积累动能。
        这一段在 $P$-$\delta$ 平面上对应“加速面积”；
  \item \textbf{故障切除后}：电磁功率恢复到（或接近）正常值，若此时 $P_e>P_T$，
        转子减速，$\delta$ 继续增大但增速减慢，积累的动能被消耗。这一段对应“减速面积”。
\end{enumerate}
\textbf{暂态稳定的判据（等面积定则的雏形）}：若最大可能的减速面积大于加速面积，
则转子在 $\delta$ 达到临界角之前减速到零，系统保持稳定；否则失步。

\begin{table}[htbp]
  \centering
  \caption{静态稳定与暂态稳定的比较}
  \small
  \begin{tabular}{@{}p{0.12\linewidth}p{0.34\linewidth}p{0.40\linewidth}@{}}
    \toprule
    比较项 & 静态稳定 & 暂态稳定 \\
    \midrule
    扰动 & 小扰动（负荷波动） & 大扰动（短路、切线路） \\
    分析方法 & 线性化、小扰动法 & 非线性、数值积分、等面积定则 \\
    判据 & $\mathrm{d}P/\mathrm{d}\delta>0$ & 减速面积 $>$ 加速面积 \\
    指标 & $K_p$（储备系数） & 极限切除角 $\delta_{cm}$、极限切除时间 \\
    关注时段 & 扰动后无穷小邻域 & 扰动后 $1\sim5$ s \\
    \bottomrule
  \end{tabular}
\end{table}

\section{负荷稳定的概念}

\textbf{负荷稳定}主要指异步电动机运行的稳定性。异步电动机的电磁转矩 $M_e$ 与转差 $s$
的关系曲线中，$M_e$ 有一个最大值 $M_{\max}$（临界转差 $s_{cr}$ 处）。
\begin{itemize}
  \item 若负载转矩 $M_L<M_{\max}$，则存在两个平衡点：$s<s_{cr}$（稳定运行点）与
        $s>s_{cr}$（不稳定运行点）；
  \item 当电压下降使 $M_{\max}$ 减小到小于 $M_L$ 时，电动机\textbf{堵转}（失速），
        转差急剧增大，定子电流剧增，可能引起电压进一步下降——这是电压崩溃的重要诱因。
\end{itemize}
负荷稳定的判据可写成 $\mathrm{d}M_e/\mathrm{d}s<0$（在稳定运行点附近）。

\section{电压稳定性的概念}

\textbf{电压稳定}指系统在扰动后维持各节点电压在可接受范围的能力。
它与功角稳定不同：功角稳定关心“发电机能否保持同步”，
电压稳定关心“负荷节点电压能否维持在可接受水平”。

\paragraph{$P$-$U$ 曲线（鼻形曲线）}
对“电源经阻抗送负荷”的系统，负荷功率 $P$ 与负荷节点电压 $U$ 的关系是一条向右弯回的曲线
（形如鼻子）。曲线的顶点（鼻点）对应\textbf{电压稳定极限}：
\begin{itemize}
  \item 鼻点以上（$U$ 较高的一支）：$\mathrm{d}P/\mathrm{d}U<0$，是\textbf{可运行区}；
        负荷增加时电压下降，系统可自我调节；
  \item 鼻点以下：$\mathrm{d}P/\mathrm{d}U>0$，不可运行；
  \item 负荷超过极限功率后，不存在平衡点，电压\textbf{崩溃}（collapse）。
\end{itemize}

\paragraph{$\mathrm{d}Q/\mathrm{d}U<0$ 判据}
在电压稳定分析中，常把“负荷所需无功 $Q_L(U)$”与“系统可提供的无功 $Q_S(U)$”画在同一张图上。
稳定运行要求
\begin{equation}
  \frac{\mathrm{d}Q}{\mathrm{d}U}<0 ,
\end{equation}
其物理含义是：电压略降时，系统能提供的无功增量大于负荷需求的增量，故能把电压“托住”。
当该条件破坏（$\mathrm{d}Q/\mathrm{d}U>0$）时，任何微小扰动都会使电压持续下降，即电压崩溃。
第 18 章 18.4.3 节将给出这一判据的详细推导与工程用法。

\corepoint{电压稳定与功角稳定的区别要能说清：\textbf{功角稳定}是“有功—角度”问题，
失稳表现为发电机失步、功率振荡；\textbf{电压稳定}是“无功—电压”问题，
失稳表现为负荷节点电压持续下降、电动机堵转。
两者可相互诱发（电压下降 $\to$ 电动机堵转 $\to$ 无功需求剧增 $\to$ 电压更低），
故实际系统常表现为“电压崩溃引发功角失稳”的复合过程。}

\section{发电机转子运动方程}

\subsection{转子运动方程}

发电机转子是一个刚体，其转动服从牛顿第二定律。设转子转动惯量为 $J$，
机械角速度为 $\omega$，机械转矩为 $M_T$、电磁转矩为 $M_e$（均取加速方向为正），则
\keyformula{J\frac{\mathrm{d}\omega}{\mathrm{d}t}=M_T-M_e=M_a}
$M_a$ 称\textbf{加速转矩}。由于功角 $\delta$ 与电角速度 $\omega$ 的关系为
$\dfrac{\mathrm{d}\delta}{\mathrm{d}t}=\omega-\omega_N$（$\omega_N$ 为同步角速度），
故转子运动方程通常写成联立形式
\begin{equation}
  \begin{cases}
    J\dfrac{\mathrm{d}\omega}{\mathrm{d}t}=M_T-M_e,\\[4pt]
    \dfrac{\mathrm{d}\delta}{\mathrm{d}t}=\omega-\omega_N .
  \end{cases}
\end{equation}
用功率表示（$M=P/\omega$，在同步速附近 $\omega\approx\omega_N$）：
\begin{equation}
  J\omega_N\frac{\mathrm{d}\omega}{\mathrm{d}t}=P_T-P_e .
\end{equation}

\textbf{转矩与功率的对应}：取转矩基值 $M_B=S_B/\omega_N$（$\omega_N$ 为同步机械角速度），
则标幺转矩 $m=M/M_B$ 与标幺功率 $p=P/S_B$ 的关系为
\begin{equation}
  m=\frac{M}{M_B}=\frac{P/\omega}{S_B/\omega_N}=p\,\frac{\omega_N}{\omega}\approx p
\end{equation}
因暂态过程中 $\omega$ 偏离额定值很小（一般不超过 $\pm2\%$），故\textbf{标幺制下转矩与功率数值相等}，
这是转子运动方程可以直接用功率差表示加速度的依据。
若用有名值，则须除以 $\omega$，且须注意机械角速度与电角速度相差极对数 $p_n$ 倍。

\subsection{用标幺值表示的转子运动方程}

取转矩基值为 $M_B=S_B/\omega_N$、时间基值为 $1/\omega_N$（即用电角度为时间单位），
可得标幺形式
\keyformula{T_J\frac{\mathrm{d}^{2}\delta}{\mathrm{d}t^{2}}=P_T-P_e,\qquad
  \frac{\mathrm{d}\delta}{\mathrm{d}t}=\omega-1}
其中 $\delta$ 单位为弧度（电角度），$t$ 单位为秒，$\omega$ 为标幺角速度，
$T_J$ 为\textbf{惯性时间常数}。

\subsection{惯性时间常数的意义}

惯性时间常数定义为
\keyformula{T_J=\frac{J\omega_N^{2}}{S_N}\quad(\text{s})}
其物理意义是：\textbf{若在转子上施加额定转矩（即额定功率），转子从静止加速到额定转速
所需的时间}。因为角加速度为 $M_N/J$，从 $0$ 到 $\omega_N$ 需时
$t=J\omega_N/M_N=J\omega_N^{2}/S_N=T_J$。

\begin{table}[htbp]
  \centering
  \caption{惯性时间常数的典型值}
  \small
  \begin{tabular}{@{}lll@{}}
    \toprule
    机组类型 & $T_J$ 典型值 & 说明 \\
    \midrule
    汽轮发电机 & $6\sim12$ s & 转速高、转子细长，$T_J$ 较小 \\
    水轮发电机 & $4\sim10$ s & 转速低、转子粗大，$J$ 大但 $\omega_N$ 小 \\
    同步调相机 & $3\sim6$ s & 惯性小，对稳定不利 \\
    \bottomrule
  \end{tabular}
\end{table}

$T_J$ 越大，转子惯性越大，同样的功率不平衡引起的角加速度越小，$\delta$ 变化越慢，
\textbf{对暂态稳定越有利}（给保护切除争取时间）；但 $\omega$ 的变化也越慢，
对调频不利。这一矛盾在第 19 章“改善元件特性”中还要提到。

\corepoint{转子运动方程是本课程唯一的\textbf{微分方程}，也是全部稳定计算的出发点。
要记牢三点：（1） 右边是“机械功率减电磁功率”（加速转矩），符号不能反；
（2） 标幺形式下“功率差直接等于角加速度乘 $T_J$”，无需再乘 $\omega$；
（3） 方程是二阶的，须两个初始条件（$\delta_0$ 与 $\Delta\omega_0=0$），
稳态起始时 $\mathrm{d}\delta/\mathrm{d}t=0$。第 17 章的数值解法、第 18 章的小扰动法
都是解这一个方程，区别只在“扰动大小”与“是否线性化”。}
  % 第 15 章
% ============================================================
% 第 16 章  电力系统的电磁功率特性
% （教材：何仰赞、温增银《电力系统分析》下册 第 16 章）
% 结构对齐 16.1—16.4；例题、习题不录
% ============================================================

\chapter{电力系统的电磁功率特性}

功角稳定问题的核心是转子运动方程，而方程的右端是 $P_T-P_e$。
机械功率 $P_T$ 由原动机决定，相对简单；\textbf{电磁功率 $P_e$ 与功角 $\delta$ 的关系}
才是稳定分析的关键——它决定了“功率极限”和“稳定边界”。
本章系统地给出 $P_e=f(\delta)$：隐极与凸极、有无电阻、有无并联支路、有无励磁调节、
简单系统与复杂系统。

\section{简单电力系统的功率特性}

\subsection{隐极式发电机的功率特性}

设发电机经变压器、线路接无限大母线（电压 $U$ 恒定）。把所有串联电抗归算到同一侧，
记 $\boldsymbol{X_{d\Sigma}=X_d+X_{T1}+X_L/2+X_{T2}}$（含全部外接电抗），
略去电阻，则隐极机（$X_d=X_q$）的电磁功率为
\keyformula{P_e=\frac{E_qU}{X_{d\Sigma}}\sin\delta,\qquad P_{\max}=\frac{E_qU}{X_{d\Sigma}}}
$\delta$ 为 $\dot E_q$ 与 $\dot U$ 之间的相位差，即功角。
功率极限 $P_{\max}$ 与 $E_qU$ 成正比、与总电抗成反比——这给出提高输送能力的三条途径：
提高电压、减小电抗（串补、分裂导线、双回线）、提高励磁（增大 $E_q$）。

\subsection{凸极式发电机的功率特性}

凸极机 $X_d\ne X_q$，$d$、$q$ 轴磁阻不同，即使不计励磁（$E_q=0$）也能传送有功，
故功率特性多出一项\textbf{磁阻功率}（反应功率）：
\keyformula{P_e=\frac{E_qU}{X_{d\Sigma}}\sin\delta
  +\frac{U^{2}}{2}\left(\frac{1}{X_{q\Sigma}}-\frac{1}{X_{d\Sigma}}\right)\sin 2\delta}
两项的物理含义：
\begin{itemize}
  \item 第一项（\textbf{基本分量}，$\sin\delta$）：由励磁磁通与电枢反应共同作用，
        与 $E_q$ 成正比，是隐极机唯一的成分；
  \item 第二项（\textbf{磁阻分量}，$\sin2\delta$）：由 $X_d\ne X_q$ 引起的磁阻变化产生，
        与励磁无关，周期为 $\pi$（$\delta=45^\circ$ 附近最大）。
\end{itemize}
由于磁阻分量在 $\delta$ 较小时为正，凸极机的功率极限比同参数隐极机略大，
且最大功率出现的角度略小于 $90^\circ$。

\begin{table}[htbp]
  \centering
  \caption{两种机型的功率特性对比}
  \small
  \begin{tabular}{@{}p{0.12\linewidth}p{0.40\linewidth}p{0.38\linewidth}@{}}
    \toprule
    机型 & 功率特性 & 特点 \\
    \midrule
    隐极机 & $P_e=\dfrac{E_qU}{X_{d\Sigma}}\sin\delta$ & 单一正弦项；$X_d=X_q$ \\
    凸极机 & $P_e=\dfrac{E_qU}{X_{d\Sigma}}\sin\delta
      +\dfrac{U^{2}}{2}\left(\dfrac{1}{X_{q\Sigma}}-\dfrac{1}{X_{d\Sigma}}\right)\sin2\delta$
      & 有磁阻功率；$E_q=0$ 时仍能送有功 \\
    \bottomrule
  \end{tabular}
\end{table}

\corepoint{凸极机多出的 $\sin2\delta$ 项常被漏写。检验方法：令 $E_q=0$（失磁），
隐极机 $P_e=0$，而凸极机仍有 $P_e\propto\sin2\delta$——这就是“凸极机可以异步运行、
隐极机不能”的根据，也是第 17 章“异步运行”一节的基础。}

\section{网络接线及参数对功率特性的影响}

\subsection{串联电阻的影响}

设发电机经 $Z=R+jX$ 接无限大母线，不计凸极效应，则
\begin{equation}
  P_e=\frac{E_qU}{Z^{2}}\left(R\cos\delta+X\sin\delta\right)-\frac{U^{2}R}{Z^{2}} .
\end{equation}
由 $\mathrm{d}P_e/\mathrm{d}\delta=0$ 得最大功率对应的功角满足
\begin{equation}
  \tan\delta_m=\frac{X}{R},
\end{equation}
即 $\delta_m<90^\circ$（$R$ 越大，$\delta_m$ 越小，越偏离 $90^\circ$）。
电阻还使功率极限降低、并在 $\delta=0$ 时产生负的功率（$R$ 上消耗有功）。
结论：\textbf{串联电阻使功率极限下降、极限角减小，对稳定不利}。
高压电网 $R\ll X$，故常略去电阻，但做精确稳定计算时应计入。

\subsection{并联电阻的影响}

并联电阻（如用恒定阻抗表示的负荷、并联电抗器的有功损耗、线路电导）从网络中分走有功，
其影响有两个方面：
\begin{itemize}
  \item \textbf{直接消耗}：并联电阻上消耗的有功随其端电压平方变化，
        使“送到受端的功率”小于“发电机发出的功率”，功率特性整体下移；
  \item \textbf{改变网络等值参数}：并联支路使网络的等值电势与等值电抗都发生变化，
        相当于在发电机与系统之间多了一条有功通路，从而改变了 $P_e$ 对 $\delta$ 的依赖关系。
\end{itemize}
定量分析时，把并联电阻与负荷一并归入节点导纳矩阵（第 4 章），
再作网络化简，得到“等值电势 $E_{eq}$、等值电抗 $X_{eq}$”，功率特性仍写成
$P_e=E_{eq}U\sin\delta/X_{eq}$ 的形式，只是参数变了。
若并联电阻较大（负荷较轻），其影响可忽略；若并联电阻较小（重负荷、低电压），
则功率极限明显下降。

\subsection{并联电抗的影响}

中间变电所装设并联电抗器（吸收容性无功）会降低该点电压，从而降低输电线路的功率极限；
反之，并联电容器（或静止补偿器）提高该点电压，可提高功率极限。

设发电机经 $X_1$ 送到中间节点（电压 $U_m$），再经 $X_2$ 送到无限大母线（$U$），
中间节点接并联电抗 $X_p$。略去电阻时，中间节点电压满足
\begin{equation}
  U_m\approx\frac{U X_p}{X_2+X_p}+\text{（与传输功率有关的部分）},
\end{equation}
即 $X_p$ 越小（并联电抗越强），$U_m$ 越低。而功率极限
\begin{equation}
  P_{\max}\approx\frac{E_qU_m}{X_1}
\end{equation}
正比于 $U_m$，故并联电抗使 $P_{\max}$ 下降，并联电容使 $P_{\max}$ 上升。

这一结论有重要的工程含义：\textbf{无功补偿不仅调压，也直接改变稳定水平}——
并联电容器、静止补偿器在提高电压的同时提高了功率极限；
而长线空载或轻载时为限制工频过电压而投入的并联电抗器，则会降低输送能力，
故运行中须按输送功率协调投切。这是第 19 章“并联补偿提高稳定性”的理论根据。

\subsection{发电机与无限大系统复合联接时的功率特性}

实际系统中发电机经多级变压器、多回线路接系统，等值电抗为各元件串联之和：
\begin{equation}
  X_{\Sigma}=X_d+X_{T1}+\frac{X_{L1}}{2}+\cdots+X_{Tn},
\end{equation}
功率特性仍为 $P_e=E_qU\sin\delta/X_{\Sigma}$ 的形式，但有三点须注意：
\begin{enumerate}
  \item $\delta$ 是“等值发电机电动势”与“系统电压”之间的夹角，它\textbf{包含了变压器与
        线路的相角偏移}，不等于发电机的实际功角；
  \item 若中间有负荷或并联支路，须先把它们归入网络，再化简为“等值电势串等值电抗”；
  \item 若网络中有环网或多电源，须先作网络化简（第 4 章），把系统化为
        “单机—无穷大”或“少机”等值模型，才可用单机公式。
\end{enumerate}
复合联接的实质是：\textbf{把“多级电压、多回线路”折算成一条等值电抗}，
使稳定分析回到单机模型。这一步做对了，后面等面积定则、$\delta_{cm}$ 公式才能用。

\section{自动励磁调节器对功率特性的影响}

\subsection{无调节励磁时发电机端电压的变化}

若无自动励磁调节（$E_q$ 恒定），当发电机出力增大、$\delta$ 增大时，定子电流增大，
发电机电抗上的压降增大，故\textbf{端电压 $U_G$ 随出力增大而下降}。
端电压下降又使功率特性“变软”：$P_e$ 随 $\delta$ 的增长比 $E_q$ 恒定时更慢，
功率极限更低。

\subsection{自动励磁调节器对功率特性的影响}

自动励磁调节器（AVR）的作用是：端电压下降时迅速增大励磁电流，维持端电压。
其效果是把功率特性从“$E_q$ 恒定”变为“$U_G$ 恒定”：
\keyformula{P_e=\frac{U_GU}{X_{\text{外}}}\sin\delta,\qquad
  P_{\max}=\frac{U_GU}{X_{\text{外}}}}
由于 $X_{\text{外}}<X_{d\Sigma}$，且 $U_G$ 在运行中基本不变，故\textbf{功率极限显著提高}，
静态稳定极限角也可从 $90^\circ$ 提高到接近 $90^\circ$（理想励磁时理论上为 $90^\circ$，
实际因调节器有死区与滞后，略小于 $90^\circ$）。

\begin{table}[htbp]
  \centering
  \caption{励磁调节对功率极限的影响}
  \small
  \begin{tabular}{@{}p{0.20\linewidth}p{0.32\linewidth}p{0.36\linewidth}@{}}
    \toprule
    励磁方式 & 功率特性 & 功率极限 \\
    \midrule
    无调节（$E_q$ 恒定） & $P_e=E_qU\sin\delta/X_{d\Sigma}$ & 较低，$U_G$ 随出力下降 \\
    比例式调节（维持 $U_G$） & $P_e=U_GU\sin\delta/X_{\text{外}}$ & 明显提高 \\
    理想调节（$U_G$ 完全恒定） & 同上，且 $\delta$ 可达 $90^\circ$ & 最高 \\
    \bottomrule
  \end{tabular}
\end{table}

\subsection{用各种电势表示的功率特性}

同一台发电机可用不同电势与相应电抗表示，对应不同的时间尺度：

\begin{table}[htbp]
  \centering
  \caption{不同电势表示的功率特性及其适用场合}
  \small
  \begin{tabular}{@{}p{0.12\linewidth}p{0.34\linewidth}p{0.26\linewidth}p{0.22\linewidth}@{}}
    \toprule
    电势 & 功率特性 & 对应电抗 & 适用 \\
    \midrule
    $E_q$ & $P_e=E_qU\sin\delta/X_{d\Sigma}$ & $X_d$ & 稳态、静态稳定 \\
    $E_q'$ & $P_e=E_q'U\sin\delta'/X_{d\Sigma}'$ & $X_d'$ & 暂态过程（阻尼绕组已衰减） \\
    $E''$ & $P_e\approx E''U\sin\delta''/X_{d\Sigma}''$ & $X_d''$ & 次暂态、短路初瞬 \\
    $U_G$ & $P_e=U_GU\sin\delta/X_{\text{外}}$ & 外接电抗 & 有强励、端电压恒定时 \\
    \bottomrule
  \end{tabular}
\end{table}

由表可见一条规律：\textbf{电势取得越“靠外”（$E_q\to E_q'\to E''\to U_G$），
所用到的电抗越小，功率极限越高，但适用的时间尺度越短}。
这解释了为什么“短路初瞬的功率极限最高（$E''$、$X_d''$），稳态的最低（$E_q$、$X_d$）”。

\section{复杂电力系统的功率特性}

多机系统中，第 $i$ 台发电机的电磁功率不仅与自己的功角有关，还与其他所有发电机的功角有关。
设各机用 $E_i\angle\delta_i$ 表示、网络用节点导纳矩阵 $\boldsymbol Y=\boldsymbol G+j\boldsymbol B$ 表示，则
\keyformula{P_{ei}=E_i^{2}G_{ii}+\sum_{\substack{j=1\\ j\ne i}}^{n}E_iE_j
  \left(G_{ij}\cos\delta_{ij}+B_{ij}\sin\delta_{ij}\right)}
其中 $\delta_{ij}=\delta_i-\delta_j$。
\begin{itemize}
  \item $E_i^{2}G_{ii}$ 是“自功率”，与自身功角无关（相当于本机对网络电阻的损耗）；
  \item 求和项是“互功率”，由各机之间的功角差决定——这就是\textbf{多机系统功率特性的耦合性}，
        也是多机系统稳定分析远比单机复杂的根本原因。
\end{itemize}
工程上常作两种简化：
\begin{enumerate}
  \item \textbf{单机无穷大等值}：把被研究机组以外的系统等值为一台无穷大母线，
        适用于“该机远离其余机组”或“其余机组惯性极大”的场合；
  \item \textbf{两机等值}：把系统分为两个机群，各自等值为一台机，适用于区域间振荡分析。
\end{enumerate}

\subsection{两机系统的相对功角特性}

两机系统中，真正决定稳定的是\textbf{相对功角} $\delta_{12}=\delta_1-\delta_2$。
把两机各自向联络线送出的功率相减，可把两机系统化为“一台等值机对无穷大母线”的形式：
\begin{equation}
  P_{12}=\frac{E_1E_2}{X_{\Sigma}}\sin\delta_{12},
  \qquad X_{\Sigma}=X_1+X_{L}+X_2 ,
\end{equation}
其中 $X_1$、$X_2$ 为两机（含升压变）的电抗，$X_L$ 为联络线电抗。
惯性时间常数也须用两机的\textbf{等值惯性常数}
\begin{equation}
  T_{J12}=\frac{T_{J1}T_{J2}}{T_{J1}+T_{J2}}
\end{equation}
（相当于两机惯量的串联）。这样，两机系统的暂态稳定分析就完全回到单机模型，
等面积定则、极限切除角公式都可直接套用。这一“相对运动”的思想，
是多机系统稳定分析中最常用的降维手段。

\subsection{功率特性在工程中的三种用法}

\begin{table}[htbp]
  \centering
  \caption{功率特性的工程用法}
  \small
  \begin{tabular}{@{}p{0.16\linewidth}p{0.34\linewidth}p{0.40\linewidth}@{}}
    \toprule
    用法 & 问的问题 & 用到的量 \\
    \midrule
    定极限 & 这条线路最多能送多少功率 & $P_{\max}=EU/X_{\Sigma}$，配合 $K_p$ \\
    定运行点 & 当前出力对应的功角是多少 & $\delta_0=\arcsin\left(P_0X_{\Sigma}/(EU)\right)$ \\
    定措施 & 提高输送能力该改哪个参数 & $\partial P_{\max}/\partial U>0$、
      $\partial P_{\max}/\partial X<0$、$\partial P_{\max}/\partial E>0$ \\
    \bottomrule
  \end{tabular}
\end{table}

由 $P_{\max}=EU/X_{\Sigma}$ 可算出“每提高 $1\%$ 电压、每减小 $1\%$ 电抗”对输送能力的贡献：
$P_{\max}$ 与 $U$、$E$ 成正比，与 $X_{\Sigma}$ 成反比。
这解释了为什么\textbf{提高电压等级、采用串联电容补偿、增加并联回路}是提高输电能力的三大手段，
而“提高励磁（增大 $E$）”则是最经济的手段（不需新增一次设备）。

\corepoint{本章的核心是一张图（$P_e$-$\delta$ 曲线）和一串参数（用哪个电势配哪个电抗）。
答题时先明确“用哪个电势、对应哪个电抗、是隐极还是凸极、有没有电阻与并联支路”，
再写功率特性——这四问决定了公式的具体形式。复杂系统则要回到
$P_{ei}=E_i^2G_{ii}+\sum E_iE_j(G_{ij}\cos\delta_{ij}+B_{ij}\sin\delta_{ij})$ 这个通式，
它是一切多机稳定程序的功率计算内核。}
  % 第 16 章
% ============================================================
% 第 17 章  电力系统暂态稳定性
% （教材：何仰赞、温增银《电力系统分析》下册 第 17 章）
% 结构对齐 17.1—17.6；例题、习题不录
% ============================================================

\chapter{电力系统暂态稳定性}

暂态稳定研究\textbf{大扰动}后系统能否保持同步运行。大扰动（短路、切线路、切机）
使电磁功率突然大幅下降，转子在“机械功率 $-$ 电磁功率”的差值下加速，
功角迅速增大。若在功角越过临界值之前能把加速动能消耗掉，系统稳定；否则失步。
本章给出这套判断的定量工具：\textbf{等面积定则}、\textbf{极限切除角}、
以及求解转子运动方程的\textbf{数值方法}。

\section{暂态稳定分析计算的基本假设}

\subsection{电力系统机电暂态过程的特点}

短路等大扰动发生后，系统中同时存在两类暂态过程：
\begin{itemize}
  \item \textbf{电磁暂态}：以毫秒计，表现为定子电流中的直流分量、倍频分量，
        以及网络中的波过程；
  \item \textbf{机电暂态}：以秒计，表现为发电机转子转速与功角的变化。
\end{itemize}
暂态稳定关心的是后者（$\delta$ 是否单调增大到失步）。两者的时间尺度相差两个数量级以上，
故可以\textbf{在机电暂态分析中忽略电磁暂态的细节}，只保留基频分量。

\subsection{基本假设}

\begin{enumerate}
  \item \textbf{忽略定子电流的直流分量与倍频分量}，只计基频分量。
        这样发电机可用正序相量模型，网络可用稳态相量方程。
  \item \textbf{忽略阻尼绕组的暂态过程}，发电机用“$E'$ 恒定、$X_d'$”模型
        （称\textbf{经典模型}）。若要计入阻尼绕组，则用 $E''$、$X_d''$，但方程数增加。
  \item \textbf{不计频率变化对网络参数的影响}，网络元件仍用额定频率下的恒定阻抗，
        网络方程是线性的代数方程。
  \item \textbf{不计原动机调速器的作用}，机械功率 $P_T$ 在暂态过程中保持不变
        （调速器动作需数秒，比暂态过程慢）。
  \item \textbf{负荷用恒定阻抗表示}（或按恒定功率、恒定电流的静特性），
        以便把负荷并入节点导纳矩阵。
  \item \textbf{不计网络电磁暂态}，即认为网络中的功率传输“瞬时”跟随功角变化。
\end{enumerate}

\subsection{近似计算中的简化}

在上述假设下，发电机等值为“$\dot E'$ 恒定（暂态电抗后电势恒定）”的电压源串 $jX_d'$，
其电磁功率为
\keyformula{P_e=\frac{E'U}{X_{d\Sigma}'}\sin\delta',\qquad X_{d\Sigma}'=X_d'+X_{\text{外}}}
$\delta'$ 为 $\dot E'$ 与 $\dot U$ 的夹角。这一模型的优点是：整个暂态过程中只有一个
状态变量 $\delta$（每台机），计算量小，且能反映暂态稳定的主要特征。
需要更精确时，可改用 $E_q$ 变化模型（计及励磁绕组暂态）或 $E''$ 模型（计及阻尼绕组）。

\section{简单电力系统暂态稳定的分析计算}

\subsection{各种运行情况下的功率特性}

以“单机—无限大系统”双回线中一回线首端短路为例，暂态过程分三个阶段，
每个阶段对应一条功率特性曲线：

\begin{table}[htbp]
  \centering
  \caption{三个阶段的功率特性}
  \small
  \begin{tabular}{@{}p{0.18\linewidth}p{0.26\linewidth}p{0.46\linewidth}@{}}
    \toprule
    阶段 & 功率特性 & 说明 \\
    \midrule
    正常运行 & $P_{\text{I}}=P_{\text{I}m}\sin\delta$ & 双回线并联，电抗最小，$P_{\text{I}m}$ 最大 \\
    故障期间 & $P_{\text{II}}=P_{\text{II}m}\sin\delta$ & 短路使等值电抗增大（或电势降低），
      $P_{\text{II}m}<P_{\text{I}m}$ \\
    故障切除后 & $P_{\text{III}}=P_{\text{III}m}\sin\delta$ & 单回线运行，$P_{\text{III}m}$ 介于两者之间 \\
    \bottomrule
  \end{tabular}
\end{table}

故障期间 $P_{\text{II}m}$ 最小，转子加速最剧烈；切除后 $P_{\text{III}m}$ 较大，
转子开始减速。三条曲线的相对高低决定了稳定与否。

\textbf{三个功率极限的求法}（这一步是暂态稳定计算中最容易出错的地方）：
\begin{enumerate}
  \item 在正序等值网络中确定短路点，按故障类型在故障点接入附加电抗（第 8 章正序等效定则），
        然后求“发电机电动势到无限大母线”的总电抗 $X_{\text{I}}$、$X_{\text{II}}$、$X_{\text{III}}$；
  \item 由 $P_{m}=E'U/X_{\Sigma}'$ 得三个阶段各自的功率极限；
  \item 注意故障期间若短路点靠近机端，$X_{\text{II}}$ 可能非常大（$P_{\text{II}m}$ 很小），
        加速面积随之剧增——这正是“近端短路最危险”的原因。
\end{enumerate}
例如单回线首端发生三相短路：故障期间发电机端电压接近于零，$P_{\text{II}}\approx0$，
加速面积达到最大；故障切除后单回线运行，$P_{\text{III}m}$ 只有双回线时的一半。
若为单相接地短路，则须在正序网中串入 $Z_{\Delta}=Z_{2\Sigma}+Z_{0\Sigma}$，
$X_{\text{II}}$ 不会太大，故\textbf{单相接地对暂态稳定的威胁远小于三相短路}。

\subsection{大扰动后发电机转子的相对运动}

短路瞬间，$P_T$ 不变而 $P_e$ 骤降至 $P_{\text{II}}$，转子受加速转矩：
\begin{equation}
  T_J\frac{\mathrm{d}^{2}\delta}{\mathrm{d}t^{2}}=P_T-P_{\text{II}m}\sin\delta>0 .
\end{equation}
$\delta$ 增大、转速上升，转子积累动能。故障切除后，若 $P_{\text{III}m}\sin\delta>P_T$，
则加速转矩变负，转子减速，但 $\delta$ 仍继续增大（惯性），直到转速恢复到同步速、
$\mathrm{d}\delta/\mathrm{d}t=0$，此时 $\delta$ 达最大值 $\delta_{\max}$。
若此后 $\delta$ 能回落，系统稳定；若 $\delta$ 越过 $180^\circ-\delta_0$（功率特性与 $P_T$ 的
另一个交点）仍未减速到零，则失步。

\subsection{等面积定则}

把转子运动方程两边对 $\delta$ 积分：
\begin{equation}
  T_J\frac{\mathrm{d}^{2}\delta}{\mathrm{d}t^{2}}=P_T-P_e
  \;\Longrightarrow\;
  \frac{1}{2}T_J\left(\frac{\mathrm{d}\delta}{\mathrm{d}t}\right)^{2}
  =\int_{\delta_0}^{\delta}\left(P_T-P_e\right)\mathrm{d}\delta .
\end{equation}
转速回到同步速时 $\mathrm{d}\delta/\mathrm{d}t=0$，故积分为零。这就是\textbf{等面积定则}：
\keyformula{\int_{\delta_0}^{\delta_c}\left(P_T-P_{\text{II}}\right)\mathrm{d}\delta
  =\int_{\delta_c}^{\delta_{\max}}\left(P_{\text{III}}-P_T\right)\mathrm{d}\delta}
左端为故障期间的\textbf{加速面积} $A_{acc}$，右端为切除后的\textbf{减速面积} $A_{dec}$。
\textbf{稳定判据}：
\begin{itemize}
  \item $A_{dec}>A_{acc}$：$\delta$ 在到达 $\delta_{\max}$ 之前减速到零，系统\textbf{稳定}；
  \item $A_{dec}=A_{acc}$：\textbf{临界状态}，对应极限切除角；
  \item $A_{dec}<A_{acc}$：转子无法消耗全部动能，$\delta$ 持续增大，\textbf{失步}。
\end{itemize}

\subsection{极限切除角}

令 $A_{acc}=A_{dec}$，即临界状态，解出\textbf{极限切除角} $\delta_{cm}$：
\keyformula{\cos\delta_{cm}
  =\frac{P_{\text{III}m}\cos\delta_{\max}-P_{\text{II}m}\cos\delta_0
  +P_T\left(\delta_0-\delta_{\max}\right)}{P_{\text{III}m}-P_{\text{II}m}}}
其中
\begin{equation}
  \delta_{\max}=\pi-\arcsin\frac{P_T}{P_{\text{III}m}}
\end{equation}
为故障切除后功率特性与 $P_T$ 的第二个交点（超过它必失步）。
由 $\delta_{cm}$ 可进一步求\textbf{极限切除时间} $t_{cm}$：把 $\delta=\delta_{cm}$ 代入转子运动方程
的数值解，反求对应的时间。工程上要求实际切除时间（保护动作时间 $+$ 断路器开断时间）
小于 $t_{cm}$，并留裕度。

\begin{table}[htbp]
  \centering
  \caption{暂态稳定的三个临界量}
  \small
  \begin{tabular}{@{}p{0.20\linewidth}p{0.30\linewidth}p{0.40\linewidth}@{}}
    \toprule
    量 & 含义 & 工程用途 \\
    \midrule
    $A_{acc}$、$A_{dec}$ & 加速、减速面积 & 等面积定则判稳（手算） \\
    $\delta_{cm}$ & 极限切除角 & 由角度判稳 \\
    $t_{cm}$ & 极限切除时间 & 整定保护动作时间、校验断路器 \\
    \bottomrule
  \end{tabular}
\end{table}

\subsection{简单电力系统暂态稳定判断的极值比较法}

除等面积定则外，另一种判断方法是\textbf{比较极值}：求出 $\delta$-$t$ 曲线的最大值 $\delta_{\max}$，
与临界角 $\delta_{cr}$（功率特性与 $P_T$ 的另一交点）比较：
\begin{itemize}
  \item $\delta_{\max}<\delta_{cr}$：稳定；
  \item $\delta_{\max}\ge\delta_{cr}$：失步。
\end{itemize}
极值比较法的优点是物理意义直观（“摆动到多大角度”），缺点是须先数值求解 $\delta(t)$，
不能像等面积定则那样直接由面积判断。

\corepoint{等面积定则只适用于“单机—无限大系统”或“两机系统”这类单自由度问题。
多机系统没有统一的“加速面积”，必须用数值积分判断各机相对角是否发散。
另外等面积定则的前提是“$P_T$ 不变、网络方程代数化”，若计入调速器与励磁暂态，
只能用数值方法。}

\section{发电机转子运动方程的数值解法}

\subsection{分段计算法}

把时间分成等步长 $\Delta t$，用差分代替微分。由转子运动方程的标幺形式，
在功率不变的假定下，第 $k$ 段的功角增量为
\begin{equation}
  \Delta\delta_k=\Delta\delta_{k-1}+K\,\Delta P_{k-1},
\end{equation}
\begin{equation}
  \delta_k=\delta_{k-1}+\Delta\delta_k,
\end{equation}
其中
\keyformula{K=\frac{360f_N\Delta t^{2}}{T_J}\quad\left(\text{度}\right)}
$\Delta P_{k-1}=P_T-P_e(\delta_{k-1})$ 为第 $k-1$ 段的加速功率。
分段计算法的精度取决于 $\Delta t$：$\Delta t$ 越小越准，但计算量越大；
工程上常取 $\Delta t=0.05$~s。其误差主要来自“假定每段内加速功率不变”。

\subsection{改进欧拉法}

把二阶方程化为一阶方程组：
\begin{equation}
  \frac{\mathrm{d}\delta}{\mathrm{d}t}=\Delta\omega,\qquad
  \frac{\mathrm{d}\Delta\omega}{\mathrm{d}t}=\frac{1}{T_J}\left(P_T-P_e(\delta)\right).
\end{equation}
\textbf{欧拉法（预测）}：
\begin{equation}
  \delta_{k+1}^{(0)}=\delta_k+\Delta\omega_k\Delta t,\qquad
  \Delta\omega_{k+1}^{(0)}=\Delta\omega_k+\frac{1}{T_J}\left(P_T-P_e(\delta_k)\right)\Delta t .
\end{equation}
\textbf{改进欧拉法（校正）}：用预测点的功率重新计算斜率，取平均
\begin{align}
  \delta_{k+1}&=\delta_k+\frac{1}{2}\left(\Delta\omega_k+\Delta\omega_{k+1}^{(0)}\right)\Delta t,\\
  \Delta\omega_{k+1}&=\Delta\omega_k+\frac{1}{2T_J}
  \left[P_T-P_e\left(\delta_k\right)+P_T-P_e\left(\delta_{k+1}^{(0)}\right)\right]\Delta t .
\end{align}
改进欧拉法是二阶精度的显式方法，比分段计算法准确，且不必迭代，
是暂态稳定计算中最常用的基本方法。更精确时可用龙格—库塔法。

\begin{table}[htbp]
  \centering
  \caption{三种数值解法的比较}
  \small
  \begin{tabular}{@{}p{0.16\linewidth}p{0.24\linewidth}p{0.24\linewidth}p{0.26\linewidth}@{}}
    \toprule
    方法 & 精度 & 计算量 & 特点 \\
    \midrule
    分段计算法 & 一阶（误差较大） & 最小 & 手算方便，须取小步长 \\
    改进欧拉法 & 二阶 & 中等 & 显式，不需迭代，工程常用 \\
    龙格—库塔法 & 四阶 & 较大 & 精度高，程序计算常用 \\
    \bottomrule
  \end{tabular}
\end{table}

\section{复杂电力系统暂态稳定的分析计算}

\subsection{大扰动后各发电机转子运动的特点}

多机系统中，大扰动后各机转子运动不再是“单摆”，而是：
\begin{itemize}
  \item 各机加速程度不同（电气距离近、惯量小的机组加速快），
        功角差 $\delta_{ij}$ 迅速拉开；
  \item 系统可能分化为\textbf{两个机群}（一群加速、一群减速），表现为\textbf{区域间振荡}；
  \item 也可能出现\textbf{多群振荡}或“一台机对系统”的模式（该机远离其余机组）。
\end{itemize}
稳定判断的实质是：\textbf{是否存在一个机群相对其余机群持续发散}。

\subsection{复杂电力系统暂态稳定的近似计算}

实用近似方法：
\begin{enumerate}
  \item \textbf{单机等值}：把被研究机组（或机群）等值为一台机，其余系统等值为无穷大母线，
        用等面积定则估算；
  \item \textbf{两机等值}：把系统分为两个机群，各自聚合为等值机（惯量、电抗按容量加权），
        化为两机系统后用相对角方程求解；
  \item \textbf{数值积分}：对全部机组联立求解转子运动方程与网络方程（交替求解），
        逐点判断各机相对角是否发散——这是工程标准做法。
\end{enumerate}

\subsection{复杂电力系统暂态稳定的判断}

工程判据通常取：
\begin{itemize}
  \item 任两台机的相对角 $\delta_{ij}$ 超过某阈值（如 $180^\circ$ 或 $360^\circ$）且持续增大，
        判为失步；
  \item 或各机相对角在仿真时段内（如 $5$~s）保持有界且趋于新的平衡，判为稳定；
  \item 也可用\textbf{能量函数法}（暂态能量裕度）作快速判断，但需构造李雅普诺夫函数。
\end{itemize}

\section{暂态稳定实际计算中系统各元件的数学模型}

（选学内容）精确的暂态稳定计算须为各元件建立动态模型，并与网络方程联立：

\begin{table}[htbp]
  \centering
  \caption{暂态稳定计算中各元件的模型}
  \small
  \begin{tabular}{@{}p{0.18\linewidth}p{0.38\linewidth}p{0.34\linewidth}@{}}
    \toprule
    元件 & 模型 & 说明 \\
    \midrule
    发电机 & 二阶（$E'$ 恒定）到五阶（计 $E_q'$、$E_d'$、$E_q''$、$E_d''$）
      & 阶数越高越精确，计算量越大 \\
    励磁系统 & 一阶或二阶传递函数（含强励顶值与时间常数） & 影响暂态过程中的 $E_q$ \\
    原动机及调速器 & 一阶惯性环节 $+$ 调差特性 & 影响 $P_T$ 的变化（慢过程） \\
    负荷 & 恒定阻抗、恒定功率、或“阻抗 $+$ 电动机”组合 & 影响扰动后的功率平衡 \\
    网络 & 代数方程（节点导纳矩阵） & 与发电机方程交替求解 \\
    \bottomrule
  \end{tabular}
\end{table}

求解思路是“\textbf{交替求解}”：每一时步先用当前 $\delta$、$E$ 解网络方程得各机电磁功率，
再积分转子运动方程得新的 $\delta$，逐步推进。

\section{电力系统异步运行的概念}

（选学内容）异步运行指发电机\textbf{失去同步但仍与系统并联}运行的状态（转差不为零）。

\subsection{发电机异步运行时的功率特性}

失步后转子转速偏离同步速，转差 $s\ne0$，定子电流产生相对转子滑动的旋转磁场，
在转子绕组中感应出转差频率电流，产生\textbf{异步转矩}。异步转矩随转差增大而增大，
故异步运行时的平均功率由“异步转矩—转差”特性决定，而不是由功角特性决定。
凸极机因有磁阻功率，异步运行能力较强；隐极机较弱。

\subsection{发电机由失步过渡到稳态异步运行的过程}

失步瞬间，转子相对角持续增大（或减小），电磁功率与机械功率在“转差”意义上重新平衡，
经过一段振荡后进入稳态异步运行：转速稳定在某个偏离同步速的值，功率以异步转矩传递。
这一过程伴随大电流、转子发热与机械振动，对机组危害很大。

\subsection{实现再同步的必要条件和促使再同步的措施}

再同步指失步机组重新被拉入同步。必要条件是：
\begin{enumerate}
  \item 系统中要有足够的\textbf{同步功率}（即功率特性足够高）把转子拉回；
  \item 转差不能太大，否则异步转矩过大、难以牵入；
  \item 机组须能承受异步运行期间的电流与发热。
\end{enumerate}
促使再同步的措施：增大励磁（提高同步功率）、减小原动机出力（降低转速）、
切除部分负荷、投入制动电阻等。

\corepoint{本章的计算骨架只有三步：（1） 写出三个阶段的功率特性 $P_{\text{I}}$、$P_{\text{II}}$、
$P_{\text{III}}$（关键是把故障期间的等值电抗画对）；（2） 用等面积定则或
$\delta_{cm}$ 公式判断稳定；（3） 若要求时间，用分段计算法或改进欧拉法积分。
判断“稳定与否”用面积，判断“切除时间够不够”用极限切除角 $+$ 数值积分，
这两件事不要混。}
  % 第 17 章
% ============================================================
% 第 18 章  电力系统静态稳定性
% （教材：何仰赞、温增银《电力系统分析》下册 第 18 章）
% 结构对齐 18.1—18.5；例题、习题不录
% ============================================================

\chapter{电力系统静态稳定性}

暂态稳定研究“大扰动后能否撑住”，静态稳定研究“小扰动后能否站住”。
后者是系统在运行点附近的\textbf{局部性质}，可以用线性化方法严格分析：
把非线性方程在工作点线性化，得到状态方程 $\Delta\dot{\boldsymbol x}=\boldsymbol A\Delta\boldsymbol x$，
再由 $\boldsymbol A$ 的特征值判断稳定性。本章的主线是：\textbf{李雅普诺夫意义下的稳定性
$\to$ 线性近似判据 $\to$ 简单系统的 $S=\mathrm{d}P_e/\mathrm{d}\delta$ 判据
$\to$ 励磁调节的影响 $\to$ 复杂系统与电压稳定判据}。

\section{运动稳定性的基本概念和小扰动法原理}

\subsection{未受扰运动与受扰运动}

设系统的运动方程为 $\dot{\boldsymbol x}=\boldsymbol f(\boldsymbol x)$，
$\boldsymbol x_0$ 为某一平衡点（$\boldsymbol f(\boldsymbol x_0)=\boldsymbol 0$），
对应的运动称\textbf{未受扰运动}。若有扰动使状态偏离到 $\boldsymbol x_0+\Delta\boldsymbol x_0$，
此后的运动称\textbf{受扰运动}。
\begin{itemize}
  \item 若受扰运动能回到 $\boldsymbol x_0$，称未受扰运动\textbf{稳定}；
  \item 若受扰运动趋于 $\boldsymbol x_0$，称\textbf{渐近稳定}；
  \item 若偏离越来越大，称\textbf{不稳定}。
\end{itemize}

\subsection{李雅普诺夫运动稳定性定义}

对任意给定的 $\varepsilon>0$，若存在 $\eta(\varepsilon)>0$，使得当
$\left\|\Delta\boldsymbol x_0\right\|<\eta$ 时，对一切 $t\ge0$ 有
$\left\|\Delta\boldsymbol x(t)\right\|<\varepsilon$，则称未受扰运动\textbf{稳定}（李雅普诺夫意义）；
若进一步有 $\lim_{t\to\infty}\Delta\boldsymbol x(t)=\boldsymbol 0$，则称\textbf{渐近稳定}。

这一定义的关键是：\textbf{稳定性是“初值扰动的影响是否有界”的性质}，
它不要求系统回到原状态，也不涉及扰动的大小（那是暂态稳定的问题）。

\subsection{非线性系统的线性近似稳定性判断法}

在平衡点附近作泰勒展开并略去高阶项，得线性化方程
\begin{equation}
  \Delta\dot{\boldsymbol x}=\boldsymbol A\Delta\boldsymbol x,
  \qquad \boldsymbol A=\left.\frac{\partial\boldsymbol f}{\partial\boldsymbol x}\right|_{\boldsymbol x_0}.
\end{equation}
设 $\boldsymbol A$ 的特征值为 $\lambda_i=\alpha_i+j\beta_i$，则
\begin{table}[htbp]
  \centering
  \caption{特征值与稳定性的关系（线性近似判据）}
  \small
  \begin{tabular}{@{}lll@{}}
    \toprule
    特征值情况 & 判定 & 说明 \\
    \midrule
    全部 $\alpha_i<0$ & \textbf{渐近稳定} & 所有模态衰减 \\
    存在 $\alpha_i>0$ & \textbf{不稳定} & 至少一个模态发散 \\
    存在 $\alpha_i=0$（其余 $<0$） & \textbf{临界}，须看高阶项 & 线性近似失效 \\
    有 $\alpha_i=0$ 且 $\beta_i\ne0$ & 等幅振荡 & 实际由阻尼决定 \\
    \bottomrule
  \end{tabular}
\end{table}

\textbf{注意}：线性近似只能判断“渐近稳定”与“不稳定”两种情形；
当出现零实部特征值时，线性化结论不可靠，必须回到非线性方程分析（李雅普诺夫第二方法
或中心流形理论）。这一条在电力系统静态稳定中正对应“$\delta_0=90^\circ$”的临界情形。

\subsection{用小扰动法分析电力系统静态稳定的步骤}

\begin{enumerate}
  \item 建立系统的完整数学模型（发电机方程 $+$ 励磁方程 $+$ 网络方程 $+$ 负荷方程）；
  \item 求出待研究的稳态运行点（潮流解）；
  \item 在运行点处线性化，消去代数变量，得到状态方程
        $\Delta\dot{\boldsymbol x}=\boldsymbol A\Delta\boldsymbol x$；
  \item 求 $\boldsymbol A$ 的特征值，按上表判断稳定性；
  \item 若不稳定，找出主导特征值（实部最大的），分析其物理来源（通常是阻尼不足或
        励磁调节引起的负阻尼）。
\end{enumerate}

\section{简单电力系统的静态稳定}

\subsection{不计发电机组的阻尼作用}

单机—无限大系统，发电机用“$E_q$ 恒定、$X_d$”模型，不计阻尼绕组与励磁调节。
由转子运动方程
\begin{equation}
  T_J\frac{\mathrm{d}^{2}\Delta\delta}{\mathrm{d}t^{2}}=-\Delta P_e
  =-\left.\frac{\mathrm{d}P_e}{\mathrm{d}\delta}\right|_{\delta_0}\Delta\delta ,
\end{equation}
记 \textbf{同步功率系数}
\keyformula{S=\left.\frac{\mathrm{d}P_e}{\mathrm{d}\delta}\right|_{\delta_0}
  =P_{m}\cos\delta_0}
则特征方程为
\begin{equation}
  T_Jp^{2}+S=0,\qquad p=\pm\sqrt{-S/T_J}.
\end{equation}
\begin{itemize}
  \item $S>0$（即 $\delta_0<90^\circ$）：$p=\pm j\sqrt{S/T_J}$，特征值为纯虚数，
        系统作\textbf{等幅振荡}——线性意义下是“临界稳定”，工程上因实际存在阻尼，
        认为是稳定的；
  \item $S<0$（即 $\delta_0>90^\circ$）：$p$ 有一正一负实根，$\Delta\delta$ 按指数增长，
        系统\textbf{单调失稳}（非周期性失步）。
\end{itemize}
故静态稳定判据为
\keyformula{S=\frac{\mathrm{d}P_e}{\mathrm{d}\delta}>0}
与第 15 章的物理判断完全一致。极限情况 $\delta_0=90^\circ$ 对应 $P_{\max}=E_qU/X_{d\Sigma}$。

\subsection{计及发电机组的阻尼作用}

实际系统中存在阻尼（阻尼绕组、机械摩擦、风阻、负荷的频率效应等），
用阻尼系数 $D$ 表示，则
\begin{equation}
  T_J\frac{\mathrm{d}^{2}\Delta\delta}{\mathrm{d}t^{2}}
  +D\frac{\mathrm{d}\Delta\delta}{\mathrm{d}t}+S\Delta\delta=0 ,
\end{equation}
特征方程
\keyformula{T_Jp^{2}+Dp+S=0,\qquad
  p=\frac{-D\pm\sqrt{D^{2}-4T_JS}}{2T_J}}
稳定条件为
\begin{equation}
  S>0\quad\text{且}\quad D>0 .
\end{equation}
\begin{itemize}
  \item $S>0$、$D>0$：特征根实部为负，$\Delta\delta$ 衰减振荡，\textbf{稳定}；
  \item $S>0$、$D<0$：实部为正，\textbf{振荡失稳}（负阻尼振荡）——这是低频振荡的机理；
  \item $S<0$：单调失稳。
\end{itemize}

\begin{table}[htbp]
  \centering
  \caption{阻尼对静态稳定的影响}
  \small
  \begin{tabular}{@{}p{0.10\linewidth}p{0.10\linewidth}p{0.34\linewidth}p{0.30\linewidth}@{}}
    \toprule
    $S$ & $D$ & 特征根 & 结论 \\
    \midrule
    $>0$ & $>0$ & 实部为负的共轭复根或两个负实根 & 稳定（衰减振荡／单调衰减） \\
    $>0$ & $<0$ & 实部为正的共轭复根 & 振荡失稳（负阻尼） \\
    $<0$ & 任意 & 一正一负实根 & 单调失稳 \\
    \bottomrule
  \end{tabular}
\end{table}

\corepoint{$S=\mathrm{d}P_e/\mathrm{d}\delta$ 是“同步力矩”系数，$D$ 是“阻尼力矩”系数。
前者决定“有没有恢复力”，后者决定“振荡能不能衰减”。
实际系统中 $D$ 可能因励磁调节、原动机调节而变为负值——这就是“有恢复力却仍然失稳”的原因，
也是第 18.3 节要讨论的核心。}

\section{自动励磁调节器对静态稳定的影响}

\subsection{按电压偏差调节的比例式调节器}

比例式调节器（AVR）按机端电压偏差 $\Delta U_G$ 调节励磁：
\begin{equation}
  \Delta E_{fd}=-K_U\Delta U_G ,
\end{equation}
其动态用一阶惯性环节表示：
\begin{equation}
  T_e\frac{\mathrm{d}\Delta E_{fd}}{\mathrm{d}t}+\Delta E_{fd}=-K_U\Delta U_G ,
\end{equation}
$T_e$ 为励磁系统等值时间常数（$0.02\sim0.5$~s），$K_U$ 为放大倍数。

把发电机方程（含 $\Delta E_q'$）、励磁方程与网络方程联立并线性化，
消去代数变量后得到三阶（或更高）特征方程：
\begin{equation}
  a_3p^{3}+a_2p^{2}+a_1p+a_0=0 .
\end{equation}
稳定性由劳斯—胡尔维茨判据判断（各系数为正且 $a_1a_2>a_0a_3$）。

\subsection{比例式调节器对静态稳定的影响}

\begin{itemize}
  \item \textbf{正面作用}：提高励磁、维持电压，使功率极限从 $E_qU/X_{d\Sigma}$ 提高到
        $U_GU/X_{\text{外}}$，\textbf{大幅提高静态稳定极限}；
  \item \textbf{负面作用}：励磁调节有滞后（$T_e$），当 $K_U$ 过大时，
        调节器的响应会滞后于功角摆动，产生\textbf{负阻尼}，
        引起\underline{低频振荡}（频率 $0.2\sim2.5$~Hz）。
\end{itemize}
因此比例式调节器存在一对矛盾：\textbf{放大倍数越大，稳态精度越高、功率极限越高，
但动态阻尼越差}。

\subsection{改进励磁调节器的几种途径}

\begin{enumerate}
  \item \textbf{引入微分环节}（PD 或 PID）：用 $\mathrm{d}U_G/\mathrm{d}t$ 的超前作用补偿滞后，
        增加正阻尼；
  \item \textbf{电力系统稳定器（PSS）}：以转速偏差 $\Delta\omega$、功率偏差 $\Delta P$
        或频率偏差为输入，产生与转速同相的附加励磁，专门提供正阻尼；
  \item \textbf{采用快速励磁系统}：减小 $T_e$，使调节更及时；
  \item \textbf{优化参数}：合理选择 $K_U$ 与各时间常数，在稳态精度与动态阻尼之间折中。
\end{enumerate}

\subsection{电力系统静态稳定的简要述评}

\begin{itemize}
  \item 无励磁调节时，静态稳定极限由 $E_q$、$X_d$ 决定，$\delta<90^\circ$；
  \item 有比例式调节时，极限提高到由 $U_G$、$X_{\text{外}}$ 决定，$\delta$ 可接近 $90^\circ$，
        但可能出现负阻尼振荡；
  \item 加装 PSS 后，既保留高功率极限，又保证正阻尼，是现代电网的标准配置；
  \item 静态稳定的\textbf{本质}是“小扰动下的恢复能力”，其判据最终归结为
        “特征值实部为负”——所有具体判据（$S>0$、劳斯判据、$\mathrm{d}Q/\mathrm{d}U<0$）
        都是这一条在不同模型下的表现。
\end{itemize}

\section{电力系统静态稳定实际分析计算的概念}

\subsection{小扰动法在复杂电力系统中的应用}

多机系统的状态方程阶数很高（每台机 $2\sim5$ 阶，加上励磁、调速、负荷），
矩阵 $\boldsymbol A$ 的维数可达数百甚至数千。实际计算：
\begin{itemize}
  \item 用\textbf{稀疏特征值算法}（如 QR 法、Arnoldi 法）只求主导特征值（实部最大的若干个）；
  \item 对每个特征值分析其\textbf{模态}（参与因子），找出“哪台机、哪个控制器”应对
        该模态负责——这是振荡源定位的依据；
  \item 特征值的虚部给出振荡频率，实部给出阻尼比 $\zeta=-\alpha/\left|\lambda\right|$，
        工程上要求 $\zeta>3\%\sim5\%$。
\end{itemize}

\subsection{静态稳定储备系数 $K_p$ 的计算问题}

工程上用\textbf{静态稳定储备系数}衡量裕度：
\keyformula{K_p=\frac{P_{\max}-P_0}{P_0}\times100\%}
计算中的几个问题：
\begin{itemize}
  \item $P_{\max}$ 应按“最不利的运行方式”（如线路检修、机组停运）计算，
        而不是按正常运行方式；
  \item 对多机系统，“$P_{\max}$”通常指“保持给定运行方式下可增加的最大输送功率”，
        须逐台机、逐条断面校核；
  \item 正常方式要求 $K_p\ge15\%\sim20\%$，事故后允许降至 $10\%$；
  \item $K_p$ 只反映“有功功率裕度”，不反映阻尼好坏，故须与阻尼比一起考核。
\end{itemize}

\subsection{$\mathrm{d}Q/\mathrm{d}U<0$ 判据的意义及电压崩溃}

对负荷节点，把“负荷所需无功”$Q_L(U)$ 与“系统可提供的无功”$Q_S(U)$ 画在同一 $Q$-$U$ 图上，
运行点为两曲线交点。稳定运行要求
\begin{equation}
  \frac{\mathrm{d}Q}{\mathrm{d}U}<0 ,
\end{equation}
其含义是：\textbf{电压略降时，系统能提供的无功增量大于负荷需求的增量}，
故电压能被“托住”。反之，若 $\mathrm{d}Q/\mathrm{d}U>0$，电压稍降就使无功缺口更大，
电压持续下降，直至\textbf{电压崩溃}。

这一判据与 $P$-$U$ 曲线的鼻点对应：鼻点处 $\mathrm{d}P/\mathrm{d}U=0$，
$\mathrm{d}Q/\mathrm{d}U$ 也变号。故\textbf{电压稳定的临界点就是 $P$-$U$ 曲线的鼻点}。

\subsection{应用 $\mathrm{d}Q/\mathrm{d}U$ 判据计算静态稳定}

实用做法：
\begin{enumerate}
  \item 用潮流计算逐点增加某区域负荷，记录负荷节点电压 $U$ 与所需无功 $Q$；
  \item 作 $Q$-$U$ 曲线（或 $P$-$U$ 曲线），找到 $\mathrm{d}Q/\mathrm{d}U=0$ 的临界点；
  \item 临界点的负荷功率即\textbf{电压稳定极限}；
  \item 定义电压稳定储备系数
        $K_U=\dfrac{P_{cr}-P_0}{P_0}\times100\%$，工程上要求 $K_U\ge10\%\sim15\%$。
\end{enumerate}
与功角稳定储备系数 $K_p$ 相比，$K_U$ 反映的是“无功与电压”的裕度，
在受端系统、大负荷中心、新能源高渗透区域尤为重要。

\section{有关电力系统运行稳定性问题的几个基本概念}

\subsection{电力系统次同步谐振和扭转振荡}

当输电线路装有\textbf{串联电容补偿}时，线路的 $LC$ 回路可能在某一次同步频率下与
发电机轴系的某一扭振模态发生谐振，称\textbf{次同步谐振}（SSR）。
其后果是轴系承受很大的交变扭矩，可能导致大轴疲劳损伤甚至断裂。
防范措施：避开谐振点（调整补偿度）、加装阻塞滤波器、采用可控串补（TCSC）。

\subsection{电力系统的低频振荡}

低频振荡指并列运行的发电机之间以 $0.2\sim2.5$~Hz 的频率发生的功率振荡，
分为：
\begin{itemize}
  \item \textbf{区域间振荡}（$0.1\sim0.7$~Hz）：机群之间通过长距离弱联系线路振荡，
        幅值大、影响范围广；
  \item \textbf{局部振荡}（$0.7\sim2.5$~Hz）：一台机（或一个电厂）对系统振荡。
\end{itemize}
主要诱因是\textbf{阻尼不足}（快速励磁引起的负阻尼、重负荷长距离输电）。
抑制手段：加装 PSS、优化励磁参数、增设附加阻尼控制（如储能、直流调制）。

\subsection{电力系统弱联的功率振荡}

“弱联”指两个系统之间仅通过一条（或少数）长距离、大电抗的联络线相连。
弱联系统的特点是：
\begin{itemize}
  \item 同步功率系数 $S$ 小（因电抗大），静态稳定储备低；
  \item 阻尼弱，容易发生区域间低频振荡；
  \item 联络线功率波动大，任一侧故障都可能引起联络线过载或失步。
\end{itemize}
因此弱联系统须严格控制联络线潮流，并配置足够的阻尼控制与解列措施。

\begin{table}[htbp]
  \centering
  \caption{三类振荡问题的比较}
  \small
  \begin{tabular}{@{}p{0.16\linewidth}p{0.16\linewidth}p{0.28\linewidth}p{0.30\linewidth}@{}}
    \toprule
    问题 & 频率范围 & 机理 & 主要对策 \\
    \midrule
    次同步谐振 & $10\sim50$ Hz（轴系扭振） & 串补电容与轴系扭振模态谐振 & 避开谐振、阻塞滤波、TCSC \\
    区域间低频振荡 & $0.1\sim0.7$ Hz & 弱联系、阻尼不足 & PSS、附加阻尼控制 \\
    局部低频振荡 & $0.7\sim2.5$ Hz & 单机对系统、励磁负阻尼 & 优化励磁、PSS \\
    \bottomrule
  \end{tabular}
\end{table}

\corepoint{本章要能把四个判据串成一条线：\textbf{$S=\mathrm{d}P_e/\mathrm{d}\delta>0$}（物理判据）、
\textbf{特征值实部为负}（数学判据）、\textbf{$K_p$ 储备系数}（工程指标）、
\textbf{$\mathrm{d}Q/\mathrm{d}U<0$}（电压稳定判据）。
前两个说的是“功角稳定”，第四个说的是“电压稳定”，不要混为一谈；
$K_p$ 与阻尼比是同一系统的两个侧面，须同时满足。
另记住“励磁调节提高极限但可能引入负阻尼，PSS 用来补阻尼”这一条因果链。}
  % 第 18 章
% ============================================================
% 第 19 章  提高电力系统稳定性的措施
% （教材：何仰赞、温增银《电力系统分析》下册 第 19 章）
% 结构对齐 19.1—19.4；例题、习题不录
% ============================================================

\chapter{提高电力系统稳定性的措施}

前四章讲“怎样判断稳定”，本章讲“怎样提高稳定”。所有措施都可以从两个式子找到落点：
\[
  P_e=P_{\max}\sin\delta=\frac{EU}{X}\sin\delta
  \qquad\text{与}\qquad
  T_J\frac{\mathrm{d}^{2}\delta}{\mathrm{d}t^{2}}=P_T-P_e .
\]
前者说明“提高功率极限”的途径（提高 $E$、$U$，减小 $X$）；
后者说明“减小加速面积、增大减速面积、缩短故障持续时间”的途径。
本章按“元件—附加装置—运行方式”三个层次把这些途径展开。

\section{提高稳定性的一般原则}

由等面积定则（第 17 章），暂态稳定的条件是减速面积大于加速面积：
\begin{equation}
  A_{dec}=\int_{\delta_c}^{\delta_{\max}}\left(P_{\text{III}}-P_T\right)\mathrm{d}\delta
  >A_{acc}=\int_{\delta_0}^{\delta_c}\left(P_T-P_{\text{II}}\right)\mathrm{d}\delta .
\end{equation}
由此得到三条一般原则：
\begin{enumerate}
  \item \textbf{提高故障后的功率特性} $P_{\text{III}}$（增大减速面积）：
        提高电压、减小电抗、增加并联回路、投入强励；
  \item \textbf{减小故障期间的功率降落}（减小加速面积）：
        采用单相重合闸、快速切除、减小故障期间的等值电抗；
  \item \textbf{缩短故障持续时间}（减小加速面积）：
        加快继电保护动作、加快断路器开断、采用快速开关。
\end{enumerate}
此外还有两条辅助原则：
\begin{itemize}
  \item \textbf{增大惯性时间常数} $T_J$（减缓 $\delta$ 变化）；
  \item \textbf{增大阻尼}（抑制振荡，提高静态稳定与动态稳定）。
\end{itemize}

\begin{table}[htbp]
  \centering
  \caption{提高稳定性的三条基本途径}
  \small
  \begin{tabular}{@{}p{0.22\linewidth}p{0.30\linewidth}p{0.38\linewidth}@{}}
    \toprule
    途径 & 作用对象 & 典型措施 \\
    \midrule
    提高功率极限 $P_{\max}$ & $E$、$U$、$X$ & 强励、串补、分裂导线、双回线、中间开关站 \\
    减小加速面积 & 故障期间的 $P_{\text{II}}$、$\delta_c$ & 快速保护、快速开关、单相重合闸、电气制动 \\
    增大减速面积 & 故障后的 $P_{\text{III}}$ & 强励、快速励磁、并联补偿、切机切负荷 \\
    \bottomrule
  \end{tabular}
\end{table}

\subsection{三条途径的定量比较}

把等面积定则写成“临界切除角”的形式（第 17 章），可以定量看出各措施的相对效果：
\begin{equation}
  \cos\delta_{cm}
  =\frac{P_{\text{III}m}\cos\delta_{\max}-P_{\text{II}m}\cos\delta_0
  +P_T\left(\delta_0-\delta_{\max}\right)}{P_{\text{III}m}-P_{\text{II}m}} .
\end{equation}
\begin{itemize}
  \item \textbf{提高 $P_{\text{III}m}$}（强励、并联补偿、增加回路）：分母增大、分子中的
        $P_{\text{III}m}\cos\delta_{\max}$ 项也增大，$\cos\delta_{cm}$ 减小，
        $\delta_{cm}$ 增大——允许切除角变大，效果显著且“线性”；
  \item \textbf{提高 $P_{\text{II}m}$}（减小故障期间的电抗、单相重合闸）：分子中
        $-P_{\text{II}m}\cos\delta_0$ 项增大（$\cos\delta_0>0$ 时使分子减小）、
        分母增大，$\delta_{cm}$ 增大——效果也很显著，且不增加稳态设备投资；
  \item \textbf{缩短切除时间}：不改变 $\delta_{cm}$，而是使实际切除角 $\delta_c<\delta_{cm}$，
        相当于“把工作点移到安全区内”——这是最直接的措施，
        且保护与开关的投资远低于一次设备。
\end{itemize}
工程上的优先级通常是：\textbf{先加快切除（最便宜），再加强励与单相重合闸，
最后才考虑串补、并补等一次设备}。

\section{改善电力系统基本元件的特性和参数}

\subsection{改善发电机及其励磁调节系统的特性}

\begin{itemize}
  \item \textbf{快速励磁与强行励磁}：故障时迅速把励磁电压升到顶值（$1.8\sim2$ 倍额定），
        维持 $E_q$、提高 $P_{\text{III}}$，是最有效、最经济的措施之一；
  \item \textbf{装设电力系统稳定器（PSS）}：提供正阻尼，抑制低频振荡，
        使高放大倍数的快速励磁系统既能提高极限又不致振荡；
  \item \textbf{减小 $X_d$}：增大发电机短路比（减小同步电抗），可提高功率极限，
        但会增加电机尺寸与造价；
  \item \textbf{增大惯性时间常数 $T_J$}：转子加重，$\delta$ 变化变慢，有利于暂态稳定，
        但对调频不利。
\end{itemize}

\subsection{改善原动机的调节特性}

\begin{itemize}
  \item \textbf{快速关闭汽门（快关）}：故障时迅速减小汽轮机进汽量，降低 $P_T$，
        从而减小加速面积；故障切除后再快速恢复。这是火电机组提高暂态稳定的重要手段；
  \item \textbf{改善调速器特性}：使原动机出力能快速跟随，但调速器过快的响应可能
        引起机械功率振荡，须与系统阻尼协调；
  \item \textbf{水轮机的特点}：水锤效应使水轮机出力在导叶快速关闭时先增后减，
        调节不当会恶化稳定，故水轮机快关须谨慎。
\end{itemize}

\subsection{减小变压器的电抗}

增大变压器容量、降低短路电压百分数 $U_s\%$，可减小 $X_T$，
从而提高功率极限、减小电压损耗。但 $U_s\%$ 过小会使短路电流增大，
须与开关设备的开断能力协调。

\subsection{改善继电保护和开关设备的特性}

\begin{itemize}
  \item \textbf{快速切除故障}：缩短 $\delta_c$（极限切除角不变，切除越早越安全），
        是最直接有效的措施；
  \item \textbf{单相自动重合闸}：单相接地故障时只切除故障相，其余两相继续运行，
        保持了线路的功率传输能力，显著减小加速面积；
  \item \textbf{综合重合闸}：根据故障类型自动选择单相或三相重合，兼顾可靠性与稳定性；
  \item \textbf{快速开关设备}：缩短开断时间（从 $3\sim5$ 周波降到 $2$ 周波）。
\end{itemize}

\subsection{改善输电线路的特性}

\begin{itemize}
  \item \textbf{采用分裂导线}：减小线路电抗（$x_1$ 从 $0.4$ 降到 $0.28\sim0.33\ \Omega$/km），
        提高功率极限；
  \item \textbf{串联电容补偿}：直接抵消部分线路电抗，是提高远距离输电能力最有效的手段；
  \item \textbf{增大导线截面}：减小电阻与电抗（电抗变化不大），主要降低网损；
  \item \textbf{增加并联回路数}：双回线使等值电抗减半，可靠性也提高。
\end{itemize}

\subsection{采用直流输电}

直流输电在提高稳定性方面有独特优势：
\begin{itemize}
  \item 直流联络线\textbf{不传输同步功率}，两端交流系统不必保持同步，
        从根本上消除了“联络线失步”问题；
  \item 直流输电可以\textbf{快速调节功率}（毫秒级），故障时迅速降低输送功率，
        减轻送端加速、受端功率缺额；
  \item 直流输电可以\textbf{隔断}两侧系统的故障传递，避免连锁反应；
  \item 直流线路本身无电抗压降、无充电功率，适合远距离大容量输电。
\end{itemize}

\section{采用附加装置提高电力系统的稳定性}

\subsection{输电线路采用串联电容补偿}

在线路中串入电容器，等值电抗变为 $X-X_C$，功率极限提高为
\begin{equation}
  P_{\max}=\frac{EU}{X-X_C}=\frac{EU}{X\left(1-K_C\right)},
  \qquad K_C=\frac{X_C}{X}.
\end{equation}
补偿度 $K_C$ 一般取 $0.3\sim0.5$（过高会引起次同步谐振与过电压）。
串补的附加好处是减小电压损耗、提高输电能力；须配置的保护有：
电容器过电压保护、间隙保护、旁路开关。

\subsection{输电线路的并联电抗补偿}

并联电抗器主要用于\textbf{限制空载或轻载时的工频过电压}（长线电容效应），
它吸收容性无功、降低电压。虽然并联电抗本身降低电压、不利于功率传输，
但在长距离超高压线路中是必需的；运行时可按需投切。
与并联电抗相反，\textbf{并联电容器与静止补偿器}可提高电压、增大功率极限，
是提高稳定性的有效手段（尤其在受端系统）。

\subsection{输电线路设立开关站}

长距离线路中途设立开关站，把线路分成若干段：
\begin{itemize}
  \item 故障时只切除一段，其余段仍可运行，减小了故障影响范围与加速面积；
  \item 开关站可装设并联补偿装置，改善电压分布；
  \item 相当于缩短了电气距离，提高了功率极限。
\end{itemize}

\subsection{中继同步调相机}

在线路中途装设同步调相机，可维持中间点电压、提供无功支撑，
从而提高功率极限（相当于把长线分成两段、中间电压恒定）。
现代系统中多由静止补偿器（SVC／STATCOM）取代，因后者响应更快、维护更简单。

\subsection{变压器中性点经小阻抗接地}

单相接地短路时，故障电流经中性点接地阻抗流通。在变压器中性点串入小阻抗 $Z_n$，
相当于在零序网中串入 $3Z_n$，增大了零序阻抗、减小了单相接地短路电流，
从而减小故障期间的功率降落（减小加速面积）。这一措施对以单相接地为主的
超高压系统效果明显，但须注意它对零序保护整定的影响。

\subsection{发电机采用电气制动}

故障发生时，迅速投入接在发电机端的\textbf{制动电阻}，把转子的加速动能转化为热能消耗掉，
相当于人为增大减速面积。制动电阻在故障切除后退出。
电气制动的优点是响应快（毫秒级）、可控；缺点是需附加设备与散热措施，
且制动过强可能引起过制动（反向失步）。

\begin{table}[htbp]
  \centering
  \caption{附加装置类措施的比较}
  \small
  \begin{tabular}{@{}p{0.20\linewidth}p{0.24\linewidth}p{0.22\linewidth}p{0.26\linewidth}@{}}
    \toprule
    措施 & 主要作用 & 响应速度 & 代价／限制 \\
    \midrule
    串联电容补偿 & 减小电抗、提高 $P_{\max}$ & 固定（不可调） & 次同步谐振、过电压保护 \\
    并联电抗器 & 限制工频过电压 & 可投切 & 降低电压，不利于输送 \\
    并联电容器／SVC & 提高电压与 $P_{\max}$ & 快（SVC 毫秒级） & 投资、谐波 \\
    中间开关站 & 缩小故障范围 & — & 占地、投资 \\
    中性点小阻抗 & 减小单相短路电流 & — & 影响零序保护 \\
    电气制动 & 消耗加速动能 & 毫秒级 & 设备与散热，可能过制动 \\
    \bottomrule
  \end{tabular}
\end{table}

\section{改善运行条件及其他措施}

\subsection{正确规定电力系统运行参数的数值}

\begin{itemize}
  \item \textbf{运行角}：正常运行时功角应留有足够裕度（$K_p\ge15\%\sim20\%$），
        避免长期在接近 $90^\circ$ 处运行；
  \item \textbf{电压水平}：维持较高的运行电压（但不超过设备允许值），
        可提高功率极限；
  \item \textbf{备用容量}：留足旋转备用，事故时可快速支援；
  \item \textbf{线路潮流}：控制重载线路的输送功率，避免任一元件的故障引起连锁过载。
\end{itemize}

\subsection{合理选择电力系统的运行接线}

\begin{itemize}
  \item 避免\textbf{弱联}（单回长线联络），必要时增设第二回线或改用直流联络；
  \item 合理配置\textbf{环网}，使任一元件退出后潮流仍能合理分布；
  \item 重要负荷采用双电源供电，重要断面避免“卡脖子”；
  \item 检修方式下须重新校核稳定，必要时限制输送功率。
\end{itemize}

\subsection{切除部分发电机及部分负荷}

\begin{itemize}
  \item \textbf{送端切机}：故障后送端功率过剩、机组加速，快速切除部分机组可减小加速面积，
        是最有效的紧急控制措施之一；
  \item \textbf{受端切负荷}：受端功率缺额、频率下降时，切除部分负荷可维持系统平衡，
        防止频率崩溃；
  \item 切机切负荷须与保护、安控装置配合，动作要快（$100\sim200$ ms 级）。
\end{itemize}

\subsection{高压直流输电功率的快速调节}

直流输电的功率可以快速、独立地调节（毫秒级）。故障时：
\begin{itemize}
  \item 送端直流功率快速降低，减小送端加速；
  \item 受端直流功率快速提升（若有备用），支援受端；
  \item 直流功率调制还可为交流系统提供附加阻尼，抑制低频振荡。
\end{itemize}

\subsection{减少系统稳定破坏所带来的损失和影响}

稳定破坏（失步、振荡、电压崩溃）一旦发生，须尽量缩小影响范围：
\begin{enumerate}
  \item \textbf{解列}：在预定地点（解列点）把系统分成若干独立部分，
        防止事故扩大；解列点应使各部分功率尽量平衡；
  \item \textbf{低频减载}：频率下降到整定值时按轮次切除负荷，防止频率崩溃；
  \item \textbf{振荡解列}：检测到失步振荡时在振荡中心解列；
  \item \textbf{黑启动}：事故后依靠具备自启动能力的机组恢复供电。
\end{enumerate}

\begin{table}[htbp]
  \centering
  \caption{提高稳定性措施的分类总结}
  \small
  \begin{tabular}{@{}p{0.14\linewidth}p{0.36\linewidth}p{0.40\linewidth}@{}}
    \toprule
    层次 & 措施 & 特点 \\
    \midrule
    元件参数 & 减小 $X$（分裂导线、双回线、低 $U_s\%$ 变压器）、增大 $T_J$、
      快速励磁与强励、PSS & 一次投资，长期有效 \\
    附加装置 & 串补、并补、开关站、调相机、中性点小阻抗、电气制动 & 针对性强，需附加设备 \\
    运行方式 & 控制运行角与潮流、合理接线、留备用 & 不花钱，效果有限 \\
    紧急控制 & 快速保护、单相重合闸、切机切负荷、直流调制、解列 & 事故时动作，是“最后一道防线” \\
    \bottomrule
  \end{tabular}
\end{table}

\corepoint{本章是全书最“工程”的一章，答题时不必罗列全部措施，而要\textbf{按机理归类}：
凡提高 $P_{\max}$ 的（减小 $X$、提高 $E$ 与 $U$）归一类；
凡减小加速面积的（快速切除、单相重合闸、快关汽门、电气制动）归一类；
凡增大减速面积的（强励、并补）归一类；
凡防止事故扩大的（解列、低频减载）归一类。
再补一句“各措施的经济性与适用场合”，答案就完整了。}
  % 第 19 章
\end{document}